\documentclass{uai2026} % for initial submission
%\documentclass[accepted]{uai2026} % after acceptance

%% Choose your variant of English; be consistent
\usepackage[american]{babel}

%% Bibliography / citations
\usepackage{natbib}
\bibliographystyle{plainnat}
\renewcommand{\bibsection}{
\subsubsection*{References}}

%% Math / tables / figures
\usepackage{mathtools}
\usepackage{amssymb}
\usepackage{amsthm}

\usepackage{microtype}
\usepackage{booktabs}
\usepackage{graphicx}
\usepackage{float}
\usepackage{subcaption}
\usepackage{xcolor}
\graphicspath{{figures/}{figures_new/}}
\usepackage{bm}
\emergencystretch=1em
\hfuzz=1pt

%% Todo notes
% \usepackage{fix_todonotes_marginparwidth}
\usepackage[colorinlistoftodos,prependcaption,textsize=tiny]{todonotes}
\presetkeys{todonotes}{inline}{}
%% Theorem environments
\newtheorem{theorem}{Theorem}[section]
\newtheorem{proposition}[theorem]{Proposition}
\newtheorem{corollary}[theorem]{Corollary}
\newtheorem{definition}[theorem]{Definition}
\newtheorem{assumption}[theorem]{Assumption}
\newtheorem*{remark}{Remark}
\newtheorem{lemma}{Lemma}
\newtheorem{example}[theorem]{Example}

%% Self-defined macros
\newcommand{\N}{\mathbb{N}}
\newcommand{\R}{\mathbb{R}}

\newcommand{\nset}{\N}
\newcommand{\rset}{\R}
\newcommand{\1}{\mathbf{1}}
\newcommand{\ind}[1]{\mathbf{1}\{#1\}}
\newcommand{\prob}{\mathbb{P}}
\newcommand{\expec}{\mathbb{E}}
\newcommand{\cprob}[2]{\prob\left(#1 \mid #2\right)}
\newcommand{\cexpec}[2]{\expec\left[#1 \mid #2\right]}
\newcommand{\test}{\mathrm{test}}
% Instance, label, and score spaces
\newcommand{\xcal}{\mathcal{X}}
\newcommand{\ycal}{\mathcal{Y}}
\newcommand{\scal}{\mathcal{S}}
\newcommand{\scalt}{\mathcal{S}_{t}}
% Definitions and bracketing
\renewcommand{\epsilon}{\varepsilon}
\newcommand{\eps}{\varepsilon}
\newcommand{\eqdef}{\coloneqq}
\newcommand{\ens}[1]{\left\{#1\right\}}
\newcommand{\open}[1]{\left(#1\right)}
\newcommand{\closed}[1]{\left[#1\right]}
\newcommand{\modu}[1]{\left|#1\right|}
\newcommand{\ecal}{\mathcal{E}}
\newcommand{\err}{\mathrm{err}}
\newcommand{\indep}{\perp \!\!\! \perp}
% Filtrations and masks
\newcommand{\F}{\mathcal{F}}
\newcommand{\G}{\mathcal{G}}
\newcommand{\fcal}{\F}
\newcommand{\gcal}{\G}
\newcommand{\mcal}{\mathcal{M}}
\newcommand{\cali}{\mathrm{cal}}
% Datasets
\newcommand{\scalcal}{\scal_{\cali}}
\newcommand{\scaltest}{\scal_{\test}}
% Conditional on $\fcal_{t}$
\newcommand{\probt}[1]{\prob\left(#1 \mid \fcal_{t}\right)}
\newcommand{\expect}[1]{\expec\left[#1 \mid \fcal_{t}\right]}

% Draft comments toggle
\newcommand{\mhegz}[1]{\todo[linecolor=blue,backgroundcolor=blue!10,bordercolor=blue,inline]{MH: #1}}
\newcommand{\mh}[1]{{\color{blue}\textbf{MH: #1}}}

\usepackage[dvipsnames]{xcolor}
\hypersetup{
  colorlinks=true,
  linkcolor=blue,
  citecolor=ForestGreen,
  urlcolor=Maroon
}

% Citation placeholder
\newcommand{\citeneeded}[1][]{{\color{red}\textbf{[#1]}}}

% Wealth/capital (Section 5 baseline)
\newcommand{\K}{K}
\newcommand{\Bnd}{B}

% Dual / subdifferential notation
\newcommand{\subdiff}{\partial}
\newcommand{\bcal}{\mathcal{B}}
\newcommand{\acal}{\mathcal{A}}
\newcommand{\argmin}{\operatorname*{argmin}}
\newcommand{\argmax}{\operatorname*{argmax}}

% Numbered assumptions environment
\newenvironment{assumptions}{%
  \begingroup
  \renewcommand{\labelenumi}{\textbf{A\arabic{enumi}}.}%
  \renewcommand{\theenumi}{A\arabic{enumi}}%
  \begin{enumerate}
    }{%
  \end{enumerate}
  \endgroup
}

\newcommand{\aymeric}[1]{\textcolor{red}{ad: #1}}

\title{Beyond Anytime: Mask-Valid Conformal Prediction}

% Filled for convenience; removed automatically for anonymous submission.
\author[1]{Anonymous Authors}
\affil[1]{Anonymous Affiliation}

\begin{document}
\maketitle
%\listoftodos % Uncomment to see all todos at once

\begin{abstract}
  When conformal prediction is deployed over sequential batches, practitioners need coverage guarantees that localize to flexible evaluation windows, renew after failures, and do not require past test labels. We introduce mask-valid conformal prediction, a framework that interpolates between marginal and anytime guarantees through binary monitoring masks. We formalize two desiderata, restart validity and calibration-data measurability, identify a class of admissible threshold policies satisfying both, and show that the design of cost-optimal schedules reduces to an optimization problem with linear constraints. When costs are forecasted rather than known in advance, we propose a replanning procedure and bound its hindsight suboptimality in terms of forecast errors. In the anytime setting, under a parametric statistical model, we show that admissible threshold schedules remain informative for polynomially many steps in the calibration size, whereas e-processes can collapse to uninformative prediction sets after only logarithmically many. Experiments on classification and regression tasks validate the framework.
\end{abstract}

%

% \section{Introduction}
% Conformal prediction provides distribution-free prediction sets with finite-sample coverage guarantees~\citep{vovk2005algorithmic,angelopoulos2024theoretical,shafer2008tutorial}.
% Given a predictor $f:\xcal\to\ycal$, a calibration dataset $\scalcal=\ens{(X_1,Y_1),\ldots,(X_n,Y_n)}\subseteq\xcal\times\ycal$ and a test point $X$, split conformal prediction constructs a set-valued map $C_\alpha:\xcal \to 2^{\ycal}$ satisfying $  \prob\ens{Y\not\in C_\alpha(X)}\le \alpha$ for the unobserved test label $Y$, whenever $(X,Y)$ is exchangeable with $\scalcal$ and $\alpha\in(0,1)$.

% In practice, data may arrive in a batch-sequential manner. For example, consider a factory where parts are continuously produced and inspected in batches at every timestep $t\in\{1,\dots,T\}$. Each new batch provides a small labeled calibration set $\scalcal^t$, alongside one or more unlabeled test features $X_t$ whose unobserved labels $Y_t$ we wish to cover with confidence sets $\acal_t(X_t)$. Similarly, in a clinical setting, a diagnostic AI may be deployed and expert radiologists definitively label a subset of medical scans at the end of each weekly shift.

% In these scenarios, information obtained from earlier batches may be used to adapt parameters for the current batch, potentially leading to arbitrary dependencies across time. Nevertheless, the underlying environment typically remains stable during the processing of an individual batch. We thus assume that, conditional on the history of past batches, the test points are exchangeable with $\scalcal^t$. Under such conditions, for each time $t$, a decision maker may directly employ split conformal prediction using $\scalcal^t$ to guarantee that
% \begin{equation}
%   \label{eq:marginal_horizon}
%   \prob\ens{Y_t\not\in \acal_t(X_t)}\le \alpha,\qquad \forall t\in\{1,\dots,T\}.
% \end{equation}
% for an unlabeled test point $X_t$ with true label $Y_t$.

% However, marginal coverage at each timestep does not address two requirements that arise naturally in continual monitoring. First, practitioners often need guarantees localized to specific evaluation windows rather than a single timestep. For instance, a quality regulator might audit a rolling window of recent predictions, or an analyst might evaluate coverage during a specific production shift. Second, coverage failures rarely terminate the system. A self-driving car that makes one incorrect prediction continues driving. Guarantees should therefore remain valid after each failure, provide gaurantees on the probability of subsequent errors.

% % if the monitoring process restarted. We call this property \emph{restart validity} and formalize it in Section~\ref{sec:admissible}.

% To address the first of these two issues,~\citet{gauthier2025values} propose anytime-valid conformal prediction, which generates the conformal predictors $\acal_1, \ldots, \acal_T$ such that
% \begin{equation}
%   \label{eq:anytime_conformal_guarantee}
%   \prob\ens{\exists t: Y_t\not\in \acal_t(X_t)} \le \alpha.
% \end{equation}
% In words, the guarantee of Eq.~\ref{eq:anytime_conformal_guarantee}, which we thereafter call the \emph{anytime guarantee}, states that the probability of a single miscoverage error across the entire sequence of predictions is controlled at level $\alpha$. To provide such guarantees, even under unknown horizons $T$,~\citet{gauthier2025values} propose leveraging conformal e-values to construct an e-process. Informally, in our setting, for a sequence of test points $(X_1, Y_1), \ldots$, an e-process is a nonnegative supermartingale $(W_t)_{t\ge 0}$ constructed from initial capital $W_0=1$. Then for each $t$, a conformal e-value $E_t(X_t, Y_t)$ is computed and used to transition the e-process to $W_t$ from $W_{t-1}$. Crucially, $(W_t)_{t\ge 0}$ is constructed such that it is a supermartingale under the assumption that, conditional on past batches, $(X_t,Y_t)$ is exchangeable with $\scalcal^t$. For the confidence sets
% \begin{equation*}
%   \acal_t(X_t) = \ens{y \in \ycal : W_{t} < 1/\alpha},
% \end{equation*}
% an application of Ville's inequality~\citep{ville1939etude} guarantees that the e-process safely remains below $1/\alpha$ with high probability, securing the anytime-valid coverage of Eq.~\eqref{eq:anytime_conformal_guarantee}.

% E-processes provide an essential and general framework for sequential testing~\citep{shafer2011test,ramdas2023game, ramdas2025hypothesis}. In the batch-sequential conformal setting, however, we argue that there exists a lot of additional structure that e-processes do not benefit from capturing and present several limitations. For example, the e-process update requires the realized test label at each step, whereas labels are often delayed or unavailable in practice. In addition, anytime validity controls a single boundary-crossing event and does not entail restart validity (see Proposition~\ref{prop:e-process-not-restart}), i.e., after a coverage failure, it is uninformative about the probability of subsequent errors. Moreover, the anytime guarantee requires that a miscoverage never happens, which can be stronger than what most settings require, and as such leads to larger prediction sets.

% \paragraph{Contributions and Outline.} In this paper, we show that in the batch-sequential conformal setting, the independence of batch-conformal $p$-values~\citep{papadopoulos2002inductive} provides considerable structure that addresses the limitations above. Unlike online adaptive methods~\citep{gibbs2021adaptive,gibbs2024conformal} that provide asymptotic frequency guarantees under distribution shift, and unlike e-process-based anytime-valid inference~\citep{grunwald2020safe,ramdas2023game,gauthier2025values}, our approach yields exact, finite-sample coverage control over flexible monitoring windows without requiring test labels or a known horizon.

% In particular, Section~\ref{sec:setup} formalizes batch-sequential conformal prediction and proves that batch-conformal $p$-values are mutually independent under conditional within-batch exchangeability.
% Section~\ref{sec:admissible} defines two desiderata---restart validity and calibration-data measurability---defines \emph{admissible} algorithms as those satisfying both, and shows that e-process-based procedures fail each.
% Exploiting $p$-value independence, Section~\ref{sec:finite} casts algorithmic design as a convex program over a log-error budget---in the spirit of multivalid conformal prediction~\citep{jung2023batch}---analyzes a sequential replanning policy, and proves pathwise suboptimality bounds controlled by forecast errors.
% Section~\ref{sec:ville} demonstrates that even without prior knowledge of $T$, a simple $\alpha$-spending schedule~\citep{lan1983discrete} provides anytime-valid coverage, and that e-process-based methods can collapse exponentially faster under a tractable parametric model encompassing Gaussian errors.
% Related work is discussed in context throughout. Experiments are presented in Section~\ref{sec:experiments}.

\section{Introduction}
Conformal prediction provides distribution-free prediction sets with finite-sample coverage guarantees~\citep{vovk2005algorithmic,angelopoulos2024theoretical,shafer2008tutorial}.
Given a predictor $f:\xcal\to\ycal$, a calibration dataset $\scalcal=\ens{(X_1,Y_1),\ldots,(X_n,Y_n)}\subseteq\xcal\times\ycal$ and a test point $X$, split conformal prediction constructs a set-valued map $C_\alpha:\xcal \to 2^{\ycal}$ satisfying $  \prob\ens{Y\not\in C_\alpha(X)}\le \alpha$ for the unobserved test label $Y$, whenever $(X,Y)$ is exchangeable with $\scalcal$ and $\alpha\in(0,1)$.

In practice, data may arrive in a batch-sequential manner. For example, consider a factory where parts are continuously produced and inspected in batches at every timestep $t\in\{1,\dots,T\}$. Each new batch provides a small labeled calibration set $\scalcal^t$, alongside one or more unlabeled test features $X_t$ whose unobserved labels $Y_t$ we wish to cover with confidence sets $\acal_t(X_t)$. Similarly, in a clinical setting, a diagnostic AI may be deployed, and expert radiologists definitively label a subset of medical scans at the end of each weekly shift.

In these scenarios, information obtained from earlier batches may be used to adapt parameters for the current batch, potentially leading to arbitrary dependencies across time. Nevertheless, the underlying environment typically remains stable during the processing of an individual batch. We thus assume that, conditional on the history of past batches, the test points are exchangeable with $\scalcal^t$. Under such conditions, for each time $t$, a decision maker may directly employ split conformal prediction using $\scalcal^t$ to guarantee that
\begin{equation}
  \label{eq:marginal_horizon}
  \prob\ens{Y_t\not\in \acal_t(X_t)}\le \alpha,\qquad \forall t\in\{1,\dots,T\}.
\end{equation}
for an unlabeled test point $X_t$ with true label $Y_t$.

% Aymeric: would push this paragraph
However, marginal coverage at each timestep does not address two requirements that arise naturally in continual monitoring.
First, practitioners often need guarantees localized to specific evaluation windows rather than a single timestep.
For instance, a quality regulator might audit a rolling window of recent predictions, or an analyst might evaluate coverage during a specific production shift.
Second, coverage failures rarely terminate the system: a self-driving car that makes one incorrect prediction continues driving.
Guarantees should therefore renew after each failure, bounding the probability of subsequent errors.

To address the first of these two issues,~\citet{gauthier2025values} propose anytime-valid conformal prediction, which generates the conformal predictors $\acal_1, \ldots, \acal_T$ such that
\begin{equation}
  \label{eq:anytime_conformal_guarantee}
  \prob\ens{\exists t: Y_t\not\in \acal_t(X_t)} \le \alpha.
\end{equation}
In words, the guarantee of Eq.~\ref{eq:anytime_conformal_guarantee}, which we thereafter call the \emph{anytime guarantee}, states that the probability of a single miscoverage error across the entire sequence of predictions is controlled at level $\alpha$. To provide such guarantees, even under unknown horizons $T$,~\citet{gauthier2025values} propose leveraging conformal e-values to construct an e-process. Informally, in our setting, for a sequence of test points $(X_1, Y_1), \ldots$, an e-process is a nonnegative supermartingale $(W_t)_{t\ge 0}$ constructed from initial capital $W_0=1$. Then for each $t$, a conformal e-value $E_t(X_t, Y_t)$ is computed and used to transition the e-process from $W_{t-1}$ to $W_t$. Crucially, $(W_t)_{t\ge 0}$ is constructed such that it is a supermartingale under the assumption that, conditional on past batches, $(X_t,Y_t)$ is exchangeable with $\scalcal^t$. For the confidence sets
\begin{equation*}
  \acal_t(X_t) := \ens{y \in \ycal : W_{t} < 1/\alpha},
\end{equation*}
an application of Ville's inequality~\citep{ville1939etude} guarantees that the e-process safely remains below $1/\alpha$ with  probability $1-\alpha$, securing the anytime-valid coverage of Eq.~\eqref{eq:anytime_conformal_guarantee}.

% Aymeric:
Both the marginal guarantee~\eqref{eq:marginal_horizon} and the anytime guarantee~\eqref{eq:anytime_conformal_guarantee} can be viewed as attempts to strengthen per-step coverage to coverage across time: the former controls each step separately, the latter controls the entire trajectory simultaneously.
In practice, many monitoring settings call for intermediate notions, e.g., coverage over a rolling window of the last $K$ batches, or over a specific audit period, that neither extreme captures.

E-values and e-processes are powerful tools for sequential inference~\citep{shafer2011test,ramdas2023game,ramdas2025hypothesis} that have opened fruitful directions in conformal prediction, including anytime-valid and adaptive conformal sets~\citep{gauthier2025values,gauthier2026adaptive} and conformal e-prediction for cross-conformal settings~\citep{vovk2025conformal}.
The batch-sequential setting, however, possesses additional structure that e-processes, designed for worst-case dependence, do not exploit.  Namely, conformal $p$-values across different timesteps can be shown to be mutually exchangeable.  Moreover, the e-process update violates multiple desiderata, which are detailed in Section~\ref{sec:admissible}. This motivates a complementary approach tailored to the statistical structure of the batch-sequential setting.

\paragraph{Contributions and Outline.} In this paper, we show that in the batch-sequential conformal setting, the independence of batch-conformal $p$-values~\citep{papadopoulos2002inductive} provides considerable structure that addresses the limitations above.
Unlike online adaptive methods~\citep{gibbs2021adaptive,gibbs2024conformal} that provide asymptotic frequency guarantees under distribution shift, and unlike e-process-based anytime-valid inference~\citep{grunwald2020safe,ramdas2023game,gauthier2025values}, our approach yields exact, finite-sample coverage control over flexible monitoring windows without requiring test labels or a known horizon.

In particular, Section~\ref{sec:setup} formalizes batch-sequential conformal prediction. Section~\ref{sec:admissible} defines two desiderata: (1) restart validity and (2) calibration-data measurability, and shows that e-process-based procedures satisfy neither.
Exploiting $p$-value independence, Section~\ref{sec:finite} casts algorithmic design as a convex program over a log-error budget. Section~\ref{sec:ville} demonstrates that even without prior knowledge of $T$, a simple $\alpha$-spending schedule~\citep{lan1983discrete} provides anytime-valid coverage, and compares the longevity of informative predictions against e-process-based methods under a parametric model encompassing Gaussian errors. Related work is discussed in context throughout. Experiments are presented in Section~\ref{sec:experiments}.

\section{Setup and Background}
\label{sec:setup}
\label{sec:batch}

This section formalizes the batch-sequential conformal prediction setting, introduces masked monitoring, and reviews anytime monitoring via e-values and e-processes.

Fix a horizon $T\in\N$ and consider batches $t=1,\dots,T$. At time $t$, we receive a calibration set $\scalcal^t \coloneqq \{(X_{t,i},Y_{t,i})\}_{i=1}^n$ of size $n>1$, which we assume to be constant across $t$ purely to ease notation. Alongside, we consider a test set~$\scaltest^t$ whose labels are unobservable at time $t$. Let $\F_0$ be a $\sigma$-algebra representing any prior data available before monitoring begins. For $t\ge 1$, let $\F_t$ be the $\sigma$-algebra generated by $\F_{t-1}$ and all data from batch $t$, including unobserved test labels. We define the sub-filtration $(\G_t)_{t\ge 0}$ where $\G_0 \subseteq \F_0$ captures observable prior data, and for $t \ge 1$, $\G_t \subset \F_t$ is generated by $\G_{t-1}$ along with the calibration set and test features of batch $t$, but not the test labels. Conditional on $\F_{t-1}$, we assume every test point $(X, Y) \in \scaltest^t$ is exchangeable with $\scalcal^t$. This assumption permits cross-batch dependence between $\scal^t \eqdef \scaltest^t \cup \scalcal^t$ and $\F_{t-1}$.

At time $t$, the conformal algorithm uses a $\fcal_{t-1}$-measurable score function $s_t:\xcal\times\ycal\to\R$. Define calibration scores $S_{t,i}:=s_t(X_{t,i},Y_{t,i})$ for $(X_{t,i},Y_{t,i})\in \scalcal^t$. Without loss of generality and to lighten notation, we assume that all scores are almost surely unique. In practice, if ties occur, they can be broken uniformly at random. For a test feature $X$, the split conformal $p$-value~\citep{papadopoulos2002inductive,papadopoulos2008inductive}, which underpins split conformal prediction~\citep{lei2018distribution,romano2019conformalized} and conformal false discovery rate control~\citep{bates2023testing,jin2023selection}, for a candidate label $y$ is
\begin{equation}
  \label{eq:pt-def}
  p_t(X,y):=\frac{1+\sum_{i=1}^{n}\ind{S_{t,i}>s_t(X,y)}}{n+1}.
\end{equation}
By construction, $p_t(X_t, Y_t)$ is sub-uniform, satisfying $\prob(p_t(X_t, Y_t) \le \alpha \mid \F_{t-1}) \le \alpha$ for any significance level $\alpha \in [0,1]$. This naturally defines the $\alpha$-level conformal prediction set $\{y \in \ycal : p_t(x,y) > \alpha\}$. We now proceed to recall the definition of conformal e-values.

\begin{definition}[Conformal e-function]
  At time $t$, a \emph{conformal e-function} is $E_t:\xcal\times\ycal\to\R_{\ge 0}$ is a $\sigma(\F_{t-1} \cup \G_t)$-measurable function such that $E_t(X,Y)$ is a \emph{conformal e-value}, i.e.,  $\expec[E_t(X,Y) \mid \F_{t-1}] \le 1$ whenever $(X,Y)$ is conditionally exchangeable with $\scalcal^t$ given $\F_{t-1}$.
\end{definition}

Conformal e-functions can be constructed in two ways, either using conformal p-values or directly from the scores. First, one can apply an $p$-to-$e$ calibrator~\citep{vovk2021values} to the conformal $p$-value $p_t(X,Y)$. A calibrator is a decreasing function $f:[0,1]\to\R_{\ge 0}$ such that $\int_0^1 f(u) du \le 1$. As such $(X,Y) \to f(p_t(X,Y))$ can be shown to be an e-function. Second, as developed recently~\citep{balinsky2024enhancing,gauthier2025values}, one can construct an e-function directly from the non-conformity scores without passing through $p$-values. Specifically, if $s_t(\cdot, \cdot) \ge 0$, a valid conformal e-function is given by
\begin{equation*}
  E_t(x, y) = \frac{(n+1)s_t(x,y)}{\left( s_t(x,y) + \sum_{i=1}^n S_{t,i} \right)}.
\end{equation*}
With conformal p-values and e-values defined, we now define sequential monitoring algorithms.

\begin{definition}[Sequential Monitoring Algorithms]
  Let $\ecal = (\epsilon_t)_{t=1,\ldots,T}$ be a sequence of algorithmic random seeds, independent of all observed data, i.e.,  $\ecal \indep \F_T$. A sequential monitoring algorithm $\acal$ outputs a set-valued mapping $X \to \acal_t(X, \epsilon_t) \subseteq \ycal$ at each step $t$. For $(X, Y) \in \scaltest^t$, define the incurred error
  \begin{equation*}
    \err_t(\acal, X, Y) \eqdef \ind{Y \not\in \acal_t(X,\epsilon_t)}.
  \end{equation*}
\end{definition}

In the remainder of the paper, we omit the explicit dependence of $\acal_t$ on $\epsilon_t$ to lighten notation; $\epsilon_t$ simply represents the internal randomness of the algorithm, which we formalize through the filtration $\mathcal{H}_t \eqdef \sigma(\epsilon_1, \dots, \epsilon_t)$. For example, $\epsilon_t$ can represent randomness used to break ties in the case of equal scores in calibration data. In addition, for $\err_t$, when clear, we omit the notational dependence on $(X,Y)$.

A direct example of such sequential algorithms is the anytime-valid conformal prediction framework introduced by~\citet{gauthier2025values}, which builds on the broader anytime-valid inference and e-value based testing literature~\citep{grunwald2020safe,vovk2021values,ramdas2023game,howard2021time} and its recent conformal extensions~\citep{vovk2025conformal,balinsky2024enhancing}.
Evaluating sequential conformal e-functions on the observed test stream yields the e-values $E_t(X_t, Y_t)$ which are nonnegative and satisfy $\expec[E_t(X_t, Y_t) \mid \F_{t-1}] \le 1$.

To construct an e-process, one introduces a predictable learning rate $\lambda_t \in [0,1]$, which is $\F_{t-1}$-measurable\citep{waudby2025universal}. The e-process $(W_t)_{t\ge 0}$ is defined recursively, starting with $W_0=1$, via the update $W_t = W_{t-1} \left(1 - \lambda_t + \lambda_t E_t(X_t, Y_t)\right)$. Then, $(W_t)_{t \ge 0}$ is a nonnegative supermartingale, and Ville's inequality~\citep{ville1939etude} implies $\prob\left(\sup_{t\ge 1} W_t\ge \alpha^{-1}\right)\le \alpha.$

At each step $t$, the conformal prediction set is constructed by inverting this e-process over all candidate $y \in \ycal$:
\begin{equation*}
  \acal_t(X) = \{y \in \ycal \colon W_{t-1} (1 - \lambda_t + \lambda_t E_t(X, y)) < \alpha^{-1}\}.
\end{equation*}
Because both $\lambda_t$ and $W_{t-1}$ depend on past features and true test labels via $\F_{t-1}$, updating $W_t$ requires the realized test labels from all previous steps.

\subsection{Masked Monitoring}

While the e-process approach of~\citet{gauthier2025values} provides anytime validity and effectively bounds the probability of an error ever occurs across the entire horizon, such guarantee can be overly stringent for many applications. In practice, a decision maker may only need guarantees over specific evaluation windows. For example, one may wish to ensure that over any contiguous window of $K$ timesteps, the probability of any miscoverage is bounded above by $\alpha$. More generally, different stakeholders may monitor different subsets of the timeline: a regulator auditing a rolling window, an analyst evaluating a particular production shift, or a system that restricts monitoring to flagged timesteps.

To capture this flexibility, we introduce the notion of \emph{mask validity}, where boolean mask $M \in \{0,1\}^T$ specifies which timesteps are actively monitored.  Mask validity requires coverage control simultaneously over an entire family of such masks and enables smooth interpolation between the marginal guarantee of Eq.~\eqref{eq:marginal_horizon} and the anytime guarantee of Eq.~\eqref{eq:anytime_conformal_guarantee}, recovering each as a special case.

Before proceeding to the definition of mask validity, we note that multiple notion of conformal prediction under sequential data arrival have been studied. For example, under online data arrival without exchangeability, Adaptive Conformal Inference~\citep{gibbs2021adaptive} and subsequent works~\citep{bhatnagar2023improved,gibbs2024conformal, angelopoulos2023conformal} provide empirical frequency guarantees that bound the time-averaged miscoverage rate. While  some works~\citep{bhatnagar2023improved} share our goal of localized monitoring and operate under weaker assumptions, their frequency control is strictly weaker than ours. In this work, we aim to bound the probability of at least one error within a given window, whereas they bound the average miscoverage rate within such window. Generalizations beyond exchangeability~\citep{barber2023conformal} address the complementary direction of non-exchangeable data within batches, whereas we assume within-batch exchangeability but allow arbitrary cross-batch distribution shifts.

\begin{definition}[Mask validity]
  \label{def:mask-validity}
  Let $M\in \ens{0,1}^T$ be a mask and $S(M) := \{t \in \{1,\dots,T\} \colon M_t=1\}$ denote the set of monitored timesteps. For mask $M$ and an algorithm $\acal$, we define the masked failure event as
  \begin{equation}
    \label{eq:mask_valid_norestart}
    F_T(M, \acal):=\left\{\exists t\in S(M)\colon \err_t(\acal) = 1\right\}.
  \end{equation}
  For a family of masks $\mcal \subseteq\{0,1\}^T$, an algorithm $\acal$ is \emph{mask-valid at level $\alpha$ over $\mcal$} if
  \begin{equation*}
    \sup_{M\in\mcal}\prob(F_T(M, \acal))\le\alpha.
  \end{equation*}
\end{definition}

We now illustrate how specific mask families recover and interpolate across the marginal and anytime guarantees.

\begin{example}[All-ones mask]
  Let $\mcal = \{(1, \ldots, 1)\}$. Then mask validity recovers standard anytime monitoring (Eq.~\ref{eq:anytime_conformal_guarantee}), controlling the error rate across the entire horizon.
\end{example}

\begin{example}[Singleton masks]
  The family of singleton masks $\mcal = \{ M^{(k)} \}_{k=1}^T$, where each $M^{(k)}$ equals $1$ at time $k$ and $0$ otherwise, recovers marginal coverage guarantees (Eq.~\ref{eq:marginal_horizon}), as it requires validity w.r.t. individual timesteps.
\end{example}

The above two examples represent extreme notions and suggest a natural interpolation. Instead of monitoring every timestep or only isolated ones, one may require that coverage holds over every contiguous window of $K$ consecutive timesteps. Such windowed guarantees arise naturally when a regulator audits rolling time intervals, or when a deployed system must be reliable over contiguous stretches.

\begin{example}[Windowed masks]
  For a window length $K \le T$, the family $\mcal = \{ M^{(j)} \}_{j=1}^{T-K+1}$, where $M^{(j)}_t = \ind{j \le t \le j+K-1}$, requires that within every contiguous window of $K$ consecutive timesteps, the probability of at least one monitored error is at most $\alpha$.
\end{example}

\subsection{Distribution of Conformal \texorpdfstring{$p$}{p}-values}
\label{sec:distribution-pvalues}

With mask validity as the target guarantee, we now identify the key structural property making it tractable. A fundamental property of the batch-sequential setup is that conditional exchangeability is sufficient to establish that the conformal $p$-values computed via \eqref{eq:pt-def} are mutually independent.

\begin{proposition}[Independenceconformal $p$-values]
  \label{prop:iid-pvalues} Consider a sequence of test points $$(X_1,Y_1), \ldots, (X_T,Y_T) \in  \scaltest^1\times \ldots\times \scaltest^T.$$

  Then, under conditional within-batch exchangeability, the corresponding conformal $p$-values $(p_1(X_1,Y_1), \ldots, p_T(X_T, Y_T))$ are mutually independent.
  Moreover, for each $t\ge 1$, the random variable $p_t(X_t, Y_t)$ is uniform on $\{1/(n+1),2/(n+1),\ldots,1\}$.
\end{proposition}

The independence established in Proposition~\ref{prop:iid-pvalues} is the key structural fact our framework exploits; it holds even when the score function $s_t$ is trained adaptively on past batches, since the dependence on the past is absorbed into $\F_{t-1}$.
%In a concurrent work,~\citet{hultberg2026anytime} study the related problem of anytime-valid conformal risk control, where the calibration dataset is accumulated over time.

%%%%%%%%%%%%%%%%%%%%%%%%%%%%%%%%%%%%%%%%%%%%%%%%%%%%%%%%%%%%%%%%%%%%%%%%%%%%%%%

% %%%%%%%%%%%%%%%%%%%%%%%%%%%%%%%%%%%%%%%%%%%%%%%%%%%%%%%%%%%%%%%%%%%%%%%%%%%%%%%

\section{Admissible Algorithms}
\label{sec:admissible}
\label{sec:beyond}

While mask validity constrains the probability of a miscoverage for any monitoring mask $M\in\mcal$, two additional requirements often arise in practice.
First, an algorithm should not depend on information that is unavailable at deployment.
Second, its guarantees should survive past failures by not turning vacuous once the first miscoverage occurs. We motivate and formalize each desideratum and identify algorithmic class satisfying both. We first define different information models for the algoirthms.

\begin{definition}[Sequential Information Models]
  \label{def:info-models}
  At each step $t$, a sequential monitoring algorithm is classified by the measurability of~$\acal_t$:
  \textbf{(1) Open-loop}: measurable with respect to $\sigma(\scalcal^t, \mathcal{H}_t)$;
  \textbf{(2) Closed-loop}: measurable with respect to $\sigma(\F_{t-1}, \scalcal^t, \mathcal{H}_t)$;
  \textbf{(3) Restricted closed-loop}: measurable with respect to $\sigma(\G_t, \mathcal{H}_t)$.
\end{definition}

In words, at time $t$, an open-loop algorithm ignores all past information $\F_{t-1}$ and relies solely on the current calibration set $\scalcal^t$ and its internal randomness.
A closed-loop algorithm may depend on the full history, including the labels of past test points in $\scaltest^1, \ldots, \scaltest^{t-1}$.
In many deployments, however, test labels are delayed, expensive, or entirely unavailable when the next prediction must be issued. Meanwhile, past calibration data and unlabeled test features are generally accessible and can be exploited to build sequentially better estimators. Our first desideratum is therefore that the algorithmic class we consider should be usable in at least a restricted closed-loop fashion, i.e., it can be used without access to previous test labels.

Beyond the information model, solely requiring mask validity can give rise to pathological behavior.
Consider a single-mask family $\mcal = \{(M)\}$ and take any mask-valid algorithm~$\acal$ with singleton test sets $(\scaltest^t \coloneqq \{(X_t, Y_t)\})_{t=1,\ldots,T}$.
Define a modified algorithm~$\acal'$ that outputs $\acal_t(X_t)$ if no miscoverage has already occurred on timestep monitored by $M$, and $\emptyset$ otherwise.
Because $\acal$ and $\acal'$ agree on all sample paths up to the first monitored error, $\acal'$ is mask-valid at the same level as~$\acal$.
Yet no practitioner would deploy~$\acal'$ as after a single miscoverage it deterministically predicts the empty set, offering no subsequent guarantees.

To rule out such behavior, we introduce our second desideratum, \emph{restart validity}. Informally, if a miscoverage occurs at a monitored time~$s$, restart validity requires that the algorithm remain mask-valid as if monitoring were restarted from time~$s+1$.

\begin{definition}[Restart validity]
  \label{def:restart-valid}
  For a mask $M\in\{0,1\}^{T}$ and algorithm~$\acal$, define the successive monitored error times by $\sigma_0^{M}\coloneqq 0$ and
  \begin{equation*}
    \sigma_{j+1}^{M}
    \coloneqq
    \inf\left\{t > \sigma_{j}^{M} \colon M_t=1,\ \err_t(\acal)=1 \right\},
  \end{equation*}
  with $\inf\emptyset=\infty$.
  Let $\mathcal{I}_t\coloneqq \sigma(\F_t,\mathcal{H}_t)$.
  An algorithm is \emph{restart-valid at level $\alpha$} over $\mcal$ if, for every $M, M'\in\mcal$, every $j\ge 0$, and every $s\in\{0,\dots,T-1\}$,
  \begin{equation*}
    \prob(\sigma_{j+1}^{M}\le T \mid \mathcal{I}_s,\; \sigma_j^{M'}=s)
    \le \alpha.
  \end{equation*}
\end{definition}

Restart validity implies geometric decay of repeated errors.

\begin{proposition}[Geometric decay of repeated errors]
  \label{prop:geom-decay}
  If $\acal$ is restart-valid at level $\alpha$ over $\mcal$, then for every $M\in\mcal$ and every integer $k\ge 1$, $\prob(\sigma_k^{M}\le T)\le \alpha^k$.
\end{proposition}

This is considerably stronger than first-error control and captures the practical reality that monitoring continues after failures.
We call an algorithm class \emph{admissible} if every member is mask-valid and restart-valid, and the class contains members that are at most restricted closed-loop, allowing for usage even when test labels are unavailable.

We now exhibit a natural class of admissible algorithms built from conformal $p$-values.
Recall that the conformal $p$-value $p_t(x,y)$ defined in~\eqref{eq:pt-def} satisfies $\prob(p_t(X_t, Y_t) \le \tau \mid \F_{t-1}) \le \tau$ for any $\tau \in [0,1]$.
A threshold policy is naturally constructed through a sequence of thresholds $(\tau_t)_{t=1,\ldots, T}$ used to convert the conformal $p$-values to confidence sets.

\begin{definition}[$p$-value thresholding policy]
  \label{def:pval-threshold-policy}
  A \emph{$p$-value thresholding policy} is parameterized by a \emph{threshold schedule} $(\tau_t)_{t=1}^{T}$ with each $\tau_t \in [0,1]$. At time $t$, the prediction set for a test feature $x$ is $\acal_t(x) = \{y \in \ycal : p_t(x, y) > \tau_t\}.$
\end{definition}

First, we note that thresholding policy class encompasses multiple information models as $\acal$ can be open-loop if each $\tau_t$ and $s_t$ are $\G_0$-measurable, restricted closed-loop if $\tau_t$ and $s_t$ are $\G_{t-1}$-measurable, or closed-loop if they are $\fcal_{t-1}$ measurable For any threshold schedule, the probability of no monitored error under mask~$M$ factorizes by independence (Proposition~\ref{prop:iid-pvalues}) and $\prob\bigl(\forall t \in S(M): p_t(X_t, Y_t) > \tau_t\bigr)$ is equal to
\begin{equation}
  \label{eq:product-form}
  \prod_{t \in S(M)} \prob(p_t(X_t, Y_t) > \tau_t)
  \ge \prod_{t \in S(M)} (1 - \tau_t).
\end{equation}
The policy is therefore mask-valid at level~$\alpha$ over~$\mcal$ whenever $\prod_{t \in S(M)} (1 - \tau_t) \ge 1-\alpha$ for every $M \in \mcal$.
The following proposition shows that mask validity automatically implies restart validity.

\begin{proposition}[Threshold policies are restart-valid]
  \label{prop:open-loop-restart}
  Under Proposition~\ref{prop:iid-pvalues}, let $(\tau_t)_{t=1}^T$ be any threshold schedule where each $\tau_t$ is $\F_{t-1}$-measurable.
  If the induced thresholding policy is mask-valid at level $\alpha$ over~$\mcal$, then it is automatically restart-valid at the same level.
\end{proposition}

The key observation is that mask validity requires $\prod_{t \in S(M)} (1-\tau_t) \ge 1-\alpha$ for each $M \in \mcal$.
Since $1-\tau_t \le 1$, any suffix of the product satisfies $\prod_{t \in S(M), t > s} (1-\tau_t) \ge \prod_{t \in S(M)} (1-\tau_t) \ge 1-\alpha$.
Because $p_t$ is independent of~$\F_{t-1}$ (Proposition~\ref{prop:iid-pvalues}), this suffix bound yields a conditional coverage guarantee from any time~$s$ onward, which is exactly restart validity. Thus, set of thresholding policies induced by such constraints form an admissible class as every element is mask-valid and restart-valid. In addition, they can be used without knowldge of past test labels.  r

% estricted closed-loop subclass can be deployed without test labels.
% Moreover, threshold policies handle multiple test points per batch naturally, since all points in~$\scaltest^t$ share the same calibration set and threshold, and their $p$-values remain independent across both points and batches.

\paragraph{E-process-based procedures are not admissible.}
By contrast, the e-process-based procedure of~\citet{gauthier2025values} fails both desiderata.
The e-process update $W_t = W_{t-1}(1-\lambda_t + \lambda_t E_t(X_t, Y_t))$ is undefined without the realized test label~$Y_t$: one cannot even compute $W_t$, let alone select thresholds for subsequent steps.
Furthermore, as formalized by the following proposition, Ville's inequality only controls a the first miscoverage and can be silent after.

\begin{proposition}
  \label{prop:e-process-not-restart}
  There exist e-process-based procedures that are anytime-valid at level $\alpha$ but not restart-valid for the all-ones mask.
\end{proposition}

\section{Optimal Scheduling}
\label{sec:finite}
\label{sec:optimal_schedule}

Having identified threshold policies as an admissible class, we turn to the problem of how to choose a threshold schedule optimally. By Eq.~\eqref{eq:product-form}, a threshold schedule $(\tau_t)_{t=1}^T$ is mask-valid at level~$\alpha$ whenever $1 - \prod_{t \in S(M)}(1-\tau_t) \le \alpha$ for every $M \in \mcal$. A smaller threshold at time~$t$ yields a larger, safer prediction set but less informative output.
We encode this tradeoff through per-step cost functions $c_t(\tau)$ and seek the schedule minimizing total cost.

Without loss of generality, we assume that $\tau_t \in [0, \tau_{\max}]$ for a fixed upper bound $\tau_{\max} \in (0,1)$. Then, if the cost functions $c_1,\ldots,c_T$ are available apriori, the optimal cost allocation problem reduces to.

\begin{equation}
  \label{eq:design-problem}
  \begin{aligned}
    \min_{0 \le \tau_{1},\dots,\tau_{T}\le \tau_{\max}}\quad & \sum_{t=1}^{T} c_t(\tau_t)\\
    \text{s.t.}\quad  1-\prod_{t\in S(M)}&(1-\tau_t)\le \alpha,\quad \forall M\in\mcal.
  \end{aligned}
\end{equation}
Substituting $s_t := -\log(1-\tau_t)$ and $L := -\log(1-\alpha)$ linearizes the product constraints and mask validity becomes $\sum_{t=1}^{T} M_t s_t \le L$ for every $M \in \mcal$.
Writing $s_{\max} := -\log(1-\tau_{\max})$ and $\tilde{c}_t(s) := c_t(1 - e^{-s})$, the design problem reduces to
\begin{equation}
  \label{eq:design-problem-log}
  \begin{aligned}
    \min_{0 \le s_{1},\dots,s_{T}\le s_{\max}}\quad & \sum_{t=1}^{T} \tilde{c}_t(s_t)\\
    \text{s.t.}\quad & \sum_{t=1}^{T} M_t s_t \le L,\qquad \forall M\in\mcal
  \end{aligned}
\end{equation}
Whenever $\tilde{c}_1,\ldots\tilde{c}_T$ are convex, the program in \ref{eq:design-problem-log} is convex with separable objective and one linear constraint per mask.
Unlike a simple Bonferroni splitting, which assigns $\alpha/T$ to each step, the optimization distributes the budget according to the cost geometry and mask structure, concentrating it where the marginal reduction in set size is largest.

\subsection{Replanning under Forecasted Costs}
\label{sec:replanning}

The main limitation of Program~\eqref{eq:design-problem-log} is the assumption that the cost functions are known at the start of deployment. In practice, the mapping from thresholds to expected set sizes may need to be learned online and may shift as the data distribution evolves. Committing to a fixed schedule based on initial, potentially inaccurate estimates can lead to highly suboptimal behavior.

Instead, an alternative approach is to dynamically adapt the allocation as new data arrives and estimates of future costs change. To this end, at time $t$, let $\hat{c}_{t,u}$ denote the forecasted cost function for a future step $u \ge t$.
At each timestep, the \emph{replanning policy} tracks the remaining log-budget $b_t(M) := L - \sum_{u < t} M_u s_u$ for each mask, and solves the forecasted tail problem
\begin{equation}
  \label{eq:replan-action}
  \hat{\bm{s}}_t \in \argmin_{\substack{0 \le s_{t:T} \le s_{\max}\\
  \sum_{u=t}^{T} M_u s_u \le b_t(M),\; \forall M \in \mcal}}
  \sum_{u=t}^{T} \hat{c}_{t,u}(s_u),
\end{equation}
commits to the immediate allocation $s_t = \hat{s}_{t|t}$, and discards the rest.
Because the policy strictly respects the budget constraints at every step, validity is immediate and restart validity follows from Proposition~\ref{prop:open-loop-restart}. In particular, the accuracy of the forecast costs affect efficiency, not safety.

\begin{corollary}[Validity of replanning]
  \label{cor:replan-validity}
  The replanning policy is mask-valid and restart-valid at level~$\alpha$ over~$\mcal$.
\end{corollary}
% \begin{proof}
%   The replanning policy produces a $\G_{t-1}$-measurable threshold schedule satisfying $\sum_{t=1}^T M_t s_t \le L$ for every $M \in \mcal$ by construction.
%   The result follows from Proposition~\ref{prop:open-loop-restart}.
% \end{proof}
%

\subsection{Suboptimality Bounds}
\label{sec:suboptimality}

To evaluate the efficiency of the replanning policy, we bound its total cost relative to the offline optimal schedule. This ground-truth baseline represents the sequence of allocations one would compute with perfect hindsight of all true realized cost functions $(\tilde{c}_t)_{t=1}^T$ prior to deployment.
The primary factor driving the suboptimality gap is the accuracy of the cost forecasts, which we measure via the uniform derivative forecast error $e_{t,u} := \sup_{s \in [0,s_{\max}]} |\tilde{c}_u'(s) - \hat{c}_{t,u}'(s)|$ with $e_{t}\eqdef e_{t,t}$.

Our formal guarantees rely on three regularity conditions (detailed in Appendix~\ref{app:replanning}): (1) strong convexity and smoothness of the cost functions, ensuring well-posed optimization admitting unique solution; (2) strictly decreasing costs, reflecting that larger log-budgets result in smaller prediction sets; and (3) trajectory interiority, which avoids degenerate replaning trajectories by assuming that the both chosen and optimal actions never deplete the budget prematurely.

To informally state the decomposition, at time $t$, let $\epsilon_t$ be the maximum distance between the dual solutions of Program \ref{eq:design-problem-log} when solved using the forecasted costs and the true costs. We provide a precise definition in Appendix~\ref{app:replanning} Eq.~\ref{eq:eps-def-app}. Then, the suboptimality of using the forecasted costs can be controlled as follows.

\begin{theorem}[Suboptimality decomposition, informal]
  \label{thm:regret-decomp-informal}
  Under the regularity conditions, the suboptimality gap is bounded by $O\!\left( \sum_{t=1}^T (e_t + \epsilon_t)^2 \right)$.
\end{theorem}

Using the above theorem as a building block, we provide further detailed analysis in the case of the all-ones masks. In particular, Theorem~\ref{thm:epsilon-allones-app} bounds $\epsilon_t$ in terms of the raw derivative errors $e_{t,u}$. Let $$\bar{e}_{t,>t} := \frac{1}{T-t} \sum_{u=t+1}^{T} e_{t,u}$$ denote the average future derivative forecast error. Then, we show that
\begin{corollary}[Suboptimality for the all-ones mask, informal]
  \label{cor:regret-allones-informal}
  For the all-ones mask $\mcal = \{\1\}$, the total suboptimality gap is bounded by
  $O\!\left( \sum_{t=1}^{T-1} \bar{e}_{t,>t}^2 + \sum_{t=1}^{T} e_{t,t}^2 \right)$.
\end{corollary}

The replanning policy connects to receding-horizon control~\citep{lin2021perturbation,lin2022bounded,li2020leveraging,vera2021online}, where per-step suboptimality is controlled via mismatch between realized and forecasted quantities.
Our setting is simpler---no dynamics or switching costs---and the suboptimality is expressed through forecast-induced discrepancies in the dual multipliers rather than primal trajectories.
Complete formal statements and proofs are in Appendix~\ref{app:replanning}.

\section{Anytime Setting Analysis}
\label{sec:ville}

In the preceding sections, we have assumed that the horizon $T$ and the mask family $\mcal$ are specified in advance. However, in practice, the decision maker may not know these aspects prior to deployment, necessitating a cautious approach that ensures validity over an indefinite horizon and for any possible monitoring mask. This consideration motivates the examination of anytime-valid e-process procedures, such as those proposed by~\citet{gauthier2025values}. To facilitate a comparative analysis of these approaches under an indefinite horizon, we adopt a tractable parametric model and compare the behvaior of the e-process against a simple anytime-valid $\alpha$-spending schedule.

Consider a regression setting with $Y = f(X) + \xi$ and squared-loss score $S(x,y) = \|y-f(x)\|^2$.
Under Gaussian errors, the squared residuals follow a chi-squared distribution, which is a special case of the Gamma family.
The Gamma model is therefore a canonical parametric setting for squared-loss conformal prediction. This makes it appealing to study the e-process already under this well-behaved, commonly used,  stationary model.

% Consider a simpler regression setting with $\ycal=\rset^d$ and $Y = f(X) + \xi$ with the squared-loss score $S(x,y) = \|y-f(x)\|^2$. If the errors $\xi$ are Gaussian with zero mean, the squared errors follow a chi-squared distribution, which is a special case of the Gamma family. More generally, for any error distribution whose squared residuals are well-modeled by a Gamma, the following assumption provides a clean analytical framework.

\begin{assumption}[I.i.d.\ Gamma scores within and across batches]
  \label{ass:gamma-scores}
  For each $t\ge 1$, the calibration scores $S_{t,1}, \dots, S_{t,n}$ and the test score $S_{t,0} = s_t(X_t, Y_t)$ are independent across batches and  satisfy
  \begin{equation*}
    S_{t,0}, S_{t,1}, \ldots, S_{t,n} \stackrel{\text{i.i.d.}}{\sim} \mathrm{Gamma}(a, \theta).
  \end{equation*}
\end{assumption}

Under this model, we compare how quickly each approach produces \emph{uninformative} prediction sets, which include every candidate label and therefore are uninformative about $Y_t$.

For a $p$-value thresholding policy with schedule $(\tau_t)_{t \ge 1}$, the conformal $p$-value $p_t(X_t, y)$ lies on the grid $\{1/(n{+}1), \ldots, 1\}$.
If $\tau_t < 1/(n{+}1)$, then $p_t(X_t, y) > \tau_t$ for every candidate label~$y$, and the prediction set is the full label space.
A similar condition can occur for the e-process. In particular, using the construction in \citet{gauthier2025values}, let $W_t \eqdef
\prod_{u=1}^t E_u$ where $E_u(X_u, y) = (n+1) S_{u}(X_u, y) / \sum_{i=0}^n S_{u,i}$. Since $S_{u}(X_u, y) \le \sum_{i=0}^n S_{u,i}$ because the test score is included in the sum, the e-values naturally satisfy $E_u(X_u,y) \le n+1$ for all candidate labels. Therefore, the prediction set $\{y : W_{t-1} E_t(X_t, y) < \alpha^{-1}\}$ becomes the full label space once $W_{t-1} < 1/(\alpha(n{+}1))$. We formalize this notation below.

\begin{definition}[Time to uninformative sets]
  \label{def:time-uninformative}
  For a $p$-value thresholding policy with schedule $(\tau_t)_{t \ge 1}$, the \emph{$p$-value time to uninformative sets} is $T_p := \inf\{t \ge 1 : \tau_t < 1/(n{+}1)\}.$
  For the e-process $(W_t)_{t \ge 0}$, the \emph{e-value time to uninformative sets} is $T_e \eqdef
  \inf\{t \ge 1 : W_{t-1} < 1/(\alpha(n{+}1))\}.$
\end{definition}

For an open-loop thresholding policy, the quantity $T_p$ is deterministic, while $T_e$ is generally random. Comparing these two quantities reveals how efficiently each approach allocates its error budget over an indefinite horizon.

We first exhibit an anytime-valid $\alpha$-spending schedule and characterize its $T_p$.

\begin{definition}[Log-squared $\alpha$-spending schedule]
  \label{def:log2-spending}
  Fix $\alpha\in(0,1)$ and let $L = -\log(1-\alpha)$.
  Define $Z := \sum_{t=1}^{\infty} (t{+}1)^{-1}(\log(t{+}1))^{-2}$.
  For $t \ge 1$, set
  \begin{equation*}
    s_t := \frac{L}{Z} \cdot \frac{1}{(t{+}1)(\log(t{+}1))^2},
    \qquad
    \tau_t := 1 - \exp(-s_t).
  \end{equation*}
\end{definition}

The sum $\sum_{t \ge 1} s_t = L$ converges by construction, and since $1 - \prod_{t=1}^\infty (1-\tau_t) = 1 - \exp(-L) = \alpha$, the sequence of thresholds $(\tau_t)_{t \ge 1}$ is a valid anytime $\alpha$-spending schedule.
The thresholds $\tau_t$ decrease monotonically to zero, so at some point $\tau_t$ drops below the conformal grid spacing~$1/(n{+}1)$ and the sets become uninformative.
The following proposition shows this schedule is anytime-valid and quantifies how long it remains informative.

\begin{proposition}[Anytime validity and informativeness of the $\alpha$-spending schedule]
  \label{prop:log2-spending}
  Under Proposition~\ref{prop:iid-pvalues}, the log-squared $\alpha$-spending schedule (Definition~\ref{def:log2-spending}) is anytime-valid at level~$\alpha$: $\prob(\exists t \ge 1 : p_t(X_t, Y_t) \le \tau_t) \le \alpha$.
  Moreover, the $p$-value time to uninformative sets satisfies $T_p = \Omega(n / (\log n)^2)$.
\end{proposition}

Under Assumption~\ref{ass:gamma-scores}, the score-based e-value is produced through scaling a Gamma random variable and normalizing by itself and the sum of other independent Gamma random variables. As such, using the folklore relation between Gamma and Dirichlet random variables, the conformal e-value is equal in distribution to scaled single coordinate of a Dirichlet random variable with uniform scale parameter and as such it is equal in distribution to a Beta random variable.

\begin{lemma}[Law of e-values under Gamma scores]
  \label{lem:beta-law}
  Under Assumption~\ref{ass:gamma-scores},
  $E_t / (n{+}1) \sim \mathrm{Beta}(a, na)$, i.i.d.\ across~$t$.
  For any $q > -a$,
  \begin{equation}
    \label{eq:Eq-moment}
    \expec[E_t^q]
    = (n{+}1)^q \frac{\Gamma(a{+}q)\,\Gamma((n{+}1)a)}{\Gamma(a)\,\Gamma((n{+}1)a{+}q)}.
  \end{equation}
\end{lemma}

Using this law, we show that the e-process collapses to the uninformative regime exponentially faster.

\begin{theorem}[Collapse bound for score-e-value e-process]
  \label{thm:Tfull-bound}
  Under Assumption~\ref{ass:gamma-scores}, fix $q \in (0,1)$ and define $c_{a,q,n} := \expec[E_1^q]$.
  Then for every $t \ge 1$,
  \begin{equation}
    \label{eq:Tfull-tail}
    \prob(T_e > t) \le (\alpha(n{+}1))^q\, c_{a,q,n}^{t-1}.
  \end{equation}
  Moreover, $c_{a,q,n} \to c_{a,q} := \Gamma(a{+}q)/(\Gamma(a)\,a^q) \in (0,1)$ as $n \to \infty$.
  Thus, $T_e = O(\log n)$ with high probability.
\end{theorem}

The contrast is stark.
Even in this favorable stationary setting with i.i.d.\ scores, the multiplicative e-process shrinks rapidly and forces uninformative prediction sets with high probability after $O(\log n)$ steps.
The $\alpha$-spending schedule, by contrast, remains informative for $\Omega(n / (\log n)^2)$ steps under the same anytime risk level.
This gap in $T_p$ versus $T_e$ reflects a potential inefficiency of e-process-based methods in the discrete setting. As e-values are bounded by $n{+}1$, the e-process can grow by at most a factor of $n{+}1$ per step, but it can shrink much faster leading to rapid depletion.

Our log-budget formulation draws on classical $\alpha$-spending~\citep{lan1983discrete} and $\alpha$-wealth allocation~\citep{foster2008alpha} from sequential analysis.
The key difference is that independence of batch-conformal $p$-values, paralleling structures long exploited in batched online testing e.g.~\citep{zrnic2020power}, allows tighter risk evaluation rather than conservative bounds. This also relates to interactive testing frameworks such as STAR~\citep{lei2017star}, which operate under general structural constraints on the rejection process.

\section{Experiments}
\label{sec:experiments}

\begin{table*}[htbp]
  \centering
  \caption{%
    Results on California Housing (Reg.) and CIFAR-10 (Clf.) under no-drift (ND) and drift (D) conditions, $T = 20$ batches, 100 runs, $\alpha = 0.15$.
    \textbf{W}: window failure rate (fraction of windows with $\ge 1$ miscoverage; target $\le 0.15$).
    \textbf{S}: mean prediction set size (interval width for regression; set cardinality for classification); $\pm$ denotes standard deviation across runs.
    Closed-loop (CL) replanning uses an exponentially weighted cost forecaster (Section~\ref{sec:experiments}).
  }
  \label{tab:results}
  \small
  \begin{tabular}{@{}l|cc|cc|cc|cc@{}}
    \toprule
    & \multicolumn{2}{c|}{Reg.\ ND}
    & \multicolumn{2}{c|}{Reg.\ D}
    & \multicolumn{2}{c|}{Clf.\ ND}
    & \multicolumn{2}{c}{Clf.\ D}                                                                                                                       \\
    \cmidrule(lr){2-3}\cmidrule(lr){4-5}\cmidrule(lr){6-7}\cmidrule(lr){8-9}
    Method                        & W                             & S            & W               & S            & W               & S             & W               & S             \\
    \midrule
    Anytime                       & --                            & $273{\pm}68$ & --              & $279{\pm}57$ & --              & $4.0{\pm}3.4$ & --              & $7.8{\pm}3.5$ \\
    Mask Valid CL ($K{=}1$) & $0.14{\pm}0.08$               & $51{\pm}28$  & $0.13{\pm}0.08$ & $104{\pm}27$ & $0.15{\pm}0.08$ & $0.9{\pm}0.3$ & $0.15{\pm}0.08$ & $3.9{\pm}2.9$ \\
    Mask Valid CL ($K{=}3$) & $0.14{\pm}0.14$               & $83{\pm}24$  & $0.13{\pm}0.13$ & $166{\pm}31$ & $0.18{\pm}0.14$ & $1.0{\pm}0.2$ & $0.14{\pm}0.13$ & $6.8{\pm}3.3$ \\
    Mask Valid CL ($K{=}5$) & $0.15{\pm}0.18$               & $101{\pm}26$ & $0.13{\pm}0.17$ & $191{\pm}35$ & $0.16{\pm}0.18$ & $1.1{\pm}0.5$ & $0.17{\pm}0.21$ & $7.7{\pm}3.1$ \\
    Mask Valid CL ($K{=}7$) & $0.12{\pm}0.20$               & $118{\pm}29$ & $0.14{\pm}0.21$ & $206{\pm}36$ & $0.15{\pm}0.22$ & $1.3{\pm}0.7$ & $0.17{\pm}0.25$ & $8.1{\pm}2.8$ \\
    Mask Valid CL ($K{=}9$) & $0.10{\pm}0.21$               & $131{\pm}32$ & $0.14{\pm}0.26$ & $216{\pm}37$ & $0.15{\pm}0.26$ & $1.4{\pm}0.9$ & $0.20{\pm}0.30$ & $8.3{\pm}2.7$ \\
    \bottomrule
  \end{tabular}
\end{table*}

We evaluate our framework against the anytime-valid baseline of~\citet{gauthier2025values} on two benchmark datasets, California Housing~\citep{california_housing} for regression and CIFAR-10~\citep{cifar10} for classification, under both stationary (ND) and drifting  (D) conditions, using a miscoverage level $\alpha = 0.15$ across $T = 20$ sequential batches with $100$ independent runs.

Full details of the experimental setup, including models, nonconformity scores, drift protocols, and results for the open-loop (OL) variant, are provided in Appendix~\ref{app:additional_material}.
\label{sec:exp-setup}

\emph{Mask Valid closed-loop (CL)} solves the log-budget program~\eqref{eq:design-problem-log} under sliding-window mask families with $K \in \{1, 3, 5, 7, 9\}$.
At each step $t$, the current cost $\hat{c}_{t,t}$ is built from the observed calibration scores, while future costs $\hat{c}_{t,u}$ for $u > t$ are forecast via an exponentially weighted average of all past unconditioned cost functions, $\hat{c}_{t,u} = \sum_{v=1}^{t} w_{t,v}\, \tilde{c}_v$, where $w_{t,v} \propto \gamma^{t-v}$ with decay $\gamma = 0.7$.
Thresholds depend only on past calibration data (restricted closed-loop), not on test labels.
\emph{Mask Valid open-loop (OL)} pre-computes a fixed threshold schedule via the static log-budget program~\eqref{eq:design-problem-log} before any data are observed, using only the prior cost model; it requires no replanning and no access to past data at deployment time.

% \textbf{Metrics.}
We report the \emph{window failure rate} (W): the fraction of audit windows, across all runs and all starting positions, that contain at least one miscoverage.
For windowed masks with parameter $K$, each window is a contiguous block of $K$ timesteps; W directly estimates $\frac{1}{|\mcal|}\sum_{M \in \mcal}
\prob(F_T(M, \acal))$, the mean window failure probability, which is bounded
above by $\alpha = 0.15$ since each term is controlled at that level.
We also report the mean prediction \emph{set size} (S) as a measure of informativeness.

% \paragraph{Results} \label{sec:results} Table~\ref{tab:results} summarises all results, and Appendix~\ref{app:figures} provides per-batch visualisations.

\textbf{Anytime baseline: structural limitations.}
Anytime CP~\citep{gauthier2025values} builds a product-martingale from conformal e-values and provides an anytime guarantee (Eq.~\ref{eq:anytime_conformal_guarantee}).
While it achieves $W{=}0$ (no monitored errors), its prediction sets are large: $\mathrm{S}{=}273{\pm}68$ for regression ND and $4.0{\pm}3.4$ for classification ND.
The per-batch traces in Appendix~\ref{app:figures} reveal the mechanism: the martingale $W_t$ collapses to near zero within the first few batches, forcing prediction sets to span the full label space for the remaining horizon.
This empirically confirms Theorem~\ref{thm:Tfull-bound}, which predicts an $O(\log n)$ collapse time, versus the $\Omega(n/\log^2 n)$ informative horizon of the $\alpha$-spending schedule (Proposition~\ref{prop:log2-spending}).
Furthermore, updating the e-process requires realized test labels at each step, while Mask Valid CL operates in restricted closed-loop mode (no test labels needed).

\textbf{Window masks: smooth interpolation between marginal and anytime validity.}
The sliding-window families with $K \in \{1, 3, 5, 7, 9\}$ smoothly interpolate between marginal coverage and anytime validity.
As $K$ decreases, the per-window guarantee weakens, requiring only that some $K$-step stretch is error-free, and the algorithm concentrates its budget, yielding smaller, more informative sets.
This is reflected directly in S: for regression ND, mean set size drops from $131$ ($K{=}9$) to $118$ ($K{=}7$), $101$ ($K{=}5$), $83$ ($K{=}3$), and $51$ ($K{=}1$), a reduction of up to $61\%$.
The same monotone pattern holds under drift and on classification, where $K{=}1$ produces near point-prediction sets (mean size ${\approx}1$ class).
Per-batch traces in Appendix~\ref{app:figures} show how the threshold $\tau_t$ rises with $K$, visually capturing the budget-allocation tradeoff that drives set-size reduction.

\textbf{Window failure rates validate the guarantee.}
The W column estimates the mean window failure probability, which our theory bounds at $\alpha{=}0.15$.
Across all CL configurations, W lies in $[0.10, 0.17]$, confirming empirical validity.
The constraint is neither vacuous (W is bounded away from zero, indicating the budget is actively consumed) nor violated (W $\le \alpha$ within Monte Carlo noise for 100 runs).
Under drift, W remains within the same range, consistent with the finite-sample guarantee that requires only within-batch exchangeability and not stationarity across batches.

% \paragraph{Anytime validity: comparable efficiency, structural advantages.}
% Both Anytime CP and Anyhow (all-ones) target the same anytime guarantee (Eq.~\ref{eq:anytime_conformal_guarantee}).
% Under stationary conditions their mean set sizes are nearly indistinguishable
% ($287{\pm}45$ vs.\ $287{\pm}37$ for regression), confirming that the $\alpha$-spending schedule
% does not sacrifice efficiency relative to wealth-based methods at this calibration-set size.
% Under drift the gap is equally small ($289{\pm}42$ vs.\ $296{\pm}20$ for regression), and the
% $W{=}0.0$ failure rates for Anyhow (all-ones) confirm validity while also indicating the
% uniform budget schedule is slightly conservative, the budget is not fully consumed, consistent
% with the $\Omega(n/\log^2 n)$ informative-horizon guarantee of Proposition~\ref{prop:log2-spending}.

% These structural advantages are achieved without accessing test labels and with the
% restart-validity and calibration-data-measurability guarantees of Section~\ref{sec:admissible},
% properties the wealth-based baseline cannot provide.
% The per-batch traces in Appendix~\ref{app:figures} reveal a mechanistic distinction:
% the Anytime martingale $W_t$ collapses to near zero within the first five batches under drift,
% forcing prediction sets to span the full label space for essentially the remaining horizon.
% Under no-drift the collapse is slower, Anytime attains $\mathrm{S}{=}7.3$ on CIFAR-10 ND versus
% $9.6$ under drift, because the wealth is sustained longer when genuine errors are rare.
% Both observations empirically confirm Theorem~\ref{thm:Tfull-bound}: collapse time scales as
% $O(\log n)$ and is sensitive to the error frequency induced by distribution shift.

% \paragraph{Window masks: smooth interpolation between marginal and anytime validity.}
% The windowed families with $K \in \{1, 3, 5\}$ smoothly bridge marginal coverage and anytime validity.
% As $K$ increases, the per-window guarantee strengthens, any $K$-step stretch must be error-free
% with probability at least $1-\alpha$, but each batch now participates in more overlapping
% windows, tightening the per-step budget and necessarily increasing prediction set sizes.
% This is reflected directly in S: for regression ND, mean set size rises from $60$
% ($K{=}1$) to $99$ ($K{=}3$), $144$ ($K{=}5$), and $287$ (all-ones), a span of $4.8\times$.
% The same monotone pattern holds under drift and on classification,
% where $K{=}1$ produces near point-prediction sets (mean size ${\approx}1$ class).

% \paragraph{Window failure rates validate the guarantee.}
% The W column estimates $\sup_{M\in\mcal}\prob(F_T(M,\acal))$, which our theory bounds at
% $\alpha{=}0.1$.
% Across all Anyhow configurations W lies in $[0.04, 0.10]$, consistent with validity.
% Two subtleties deserve mention.
% First, W is bounded away from zero for windowed masks, confirming the budget is actively
% consumed and the guarantee is non-vacuous.
% Second, W is near zero for the all-ones mask, reflecting the conservative uniform schedule;
% this is not a sign of an inactive constraint but of a schedule that spends its budget slowly
% by design.
% Under drift, W remains within the same range as under no-drift, consistent with the
% finite-sample guarantee that requires only within-batch exchangeability and
% not stationarity across batches.
% We note that with 20 independent runs the estimation uncertainty on W is substantial
% (the reported standard deviations range from $0.03$ to $0.30$); the failure-rate validation
% is therefore qualitative rather than a precise empirical certificate.
\section{Discussion}
\label{sec:discussion}

Mask-valid conformal prediction replaces anytime monitoring by a family of structured monitoring rules specified ex ante, while treating the realized mask as unobserved online.
In the batched sequential model under conditional exchangeability, randomized conformal $p$-values are grid-uniform conditional on the past, which implies independence from the past and i.i.d.\ $p$-values across time even without batch independence.
This i.i.d.\ structure yields an exact finite-horizon evaluation formula and a log-budget reduction of mask-validity to linear constraints.
Computational aspects of large mask families are deferred.
Appendix~\ref{app:replanning} provides the complete proofs underlying the replanning bounds, and the section-specific appendices provide other proofs deferred from the main text.

\bibliography{references}
\appendix
\onecolumn

\section*{Organization of the Appendix}

Appendix~\ref{app:proofs_2_3} provides proofs deferred from Sections~2 and~3. Appendix~\ref{app:proofs_4} provides proofs deferred from Section~4. Appendix~\ref{app:proofs_5} provides proofs for Section~5. Appendix~\ref{app:datasets}--\ref{app:figures} contain dataset details and extended experimental results.

\section{Proofs of Sections 2 and 3}
\label{app:proofs_2_3}

\begin{proof}[Proof of Proposition~\ref{prop:iid-pvalues}]
  Fix $t\ge 1$ and a test point $(X_t,Y_t)\in\scaltest^t$.
  Conditional on $\F_{t-1}$, $(X_t,Y_t)$ is exchangeable with the $n$ calibration points in $\scalcal^t$, and $s_t$ is $\F_{t-1}$-measurable.

  Write $S_{t,0}\coloneqq s_t(X_t,Y_t)$ and $S_{t,i}\coloneqq s_t(X_{t,i},Y_{t,i})$ for $i\in\{1,\dots,n\}$.
  Assuming almost-sure uniqueness of the $n+1$ scores, exchangeability implies that the rank
  \begin{equation*}
    R_t
    \coloneqq
    1+\sum_{i=1}^{n}\ind{S_{t,i}>S_{t,0}}
    \in\{1,\dots,n+1\}
  \end{equation*}
  is uniform on $\{1,\dots,n+1\}$ conditional on $\F_{t-1}$.
  Therefore $p_t(X_t,Y_t)=R_t/(n+1)$ is uniform on $\{1/(n+1),\dots,1\}$ conditional on $\F_{t-1}$.

  Note that this conditional distribution does not depend on $\F_{t-1}$. Thus, for any $T\ge 2$, any measurable sets $A_1,\dots,A_T\subseteq\{1/(n+1),\dots,1\}$, and any $t\ge 2$, applying the tower property yields
  \begin{equation*}
    \begin{aligned}
      \prob\left(\bigcap_{u=1}^{t}\{p_u\in A_u\}\right)
      &= \expec\left[ \expec\left[ \ind{\bigcap_{u=1}^{t}\{p_u\in A_u\}} \middle| \F_{t-1} \right] \right] \\
      &= \expec\left[\ind{\bigcap_{u=1}^{t-1}\{p_u\in A_u\}}\prob(p_t\in A_t\mid \F_{t-1})\right] \\
      &= \prob(p_t\in A_t)\prob\left(\bigcap_{u=1}^{t-1}\{p_u\in A_u\}\right),
    \end{aligned}
  \end{equation*}
  because $\bigcap_{u=1}^{t-1}\{p_u\in A_u\}\in\F_{t-1}$ and $\prob(p_t\in A_t\mid \F_{t-1})=\prob(p_t\in A_t)$ almost surely.
  Iterating this recursion yields $\prob(\bigcap_{u=1}^{T}\{p_u\in A_u\})=\prod_{u=1}^{T}\prob(p_u\in A_u)$, which is mutual independence.
\end{proof}

\begin{proof}[Proof of Proposition~\ref{prop:e-process-not-restart}]
  Let $T=2$ and consider the all-ones mask.
  Fix $\alpha \in (0, 1)$ and choose the calibration size $n$ large enough so that $1/(n+1) \le \alpha^2 / 2$.
  Define the thresholds
  \begin{equation*}
    \beta = \max\left\{p \in \left\{\frac{1}{n+1}, \dots, 1\right\} : p \le \alpha^2\right\},
    \qquad
    \gamma = \max\left\{p \in \left\{\frac{1}{n+1}, \dots, 1\right\} : p \le \sqrt{\alpha}\right\}.
  \end{equation*}
  Since $1/(n+1) \le \alpha^2/2$, both $\beta$ and $\gamma$ are strictly positive, well-defined grid values, satisfying $\beta \le \alpha^2$ and $\gamma \le \sqrt{\alpha}$.
  Moreover, since $\alpha \in (0,1)$, for large enough $n$ we strictly have $\gamma > \alpha$.

  Define the valid all-or-nothing $p$-to-$e$ calibrators $f_1(p) = \beta^{-1}\ind{p \le \beta}$ and $f_2(p) = \gamma^{-1}\ind{p \le \gamma}$, and let $E_t(x, y) = f_t(p_t(x, y))$.
  By Proposition~\ref{prop:iid-pvalues}, $p_t(X_t, Y_t)$ is uniform on $\{1/(n+1), \dots, 1\}$ and independent of the past, so $\expec[f_t(p_t(X_t, Y_t)) \mid \mathcal{F}_{t-1}] \leq 1$.

  An error occurs at $t=1$ if $W_1 = f_1(p_1(X_1, Y_1)) \ge \alpha^{-1}$, which happens if and only if $p_1(X_1, Y_1) \le \beta$ (because $\beta^{-1} \ge \alpha^{-2} > \alpha^{-1}$).
  Since $\beta \ge 1/(n+1)$, this event $\{\sigma_1 = 1\}$ occurs with strictly positive probability $\beta > 0$.
  On this event, the e-process is exactly $W_1 = \beta^{-1}$.

  Conditional on $\mathcal{I}_1$, an error occurs at $t=2$ if $W_1 f_2(p_2(X_2, Y_2)) \ge \alpha^{-1}$.
  Substituting $W_1 = \beta^{-1}$, this condition is $\beta^{-1} \gamma^{-1} \ind{p_2(X_2, Y_2) \le \gamma} \ge \alpha^{-1}$.
  Because $\beta \le \alpha^2$ and $\gamma \le \sqrt{\alpha}$, we have $\beta^{-1} \gamma^{-1} \ge \alpha^{-2} \alpha^{-1/2} = \alpha^{-5/2} > \alpha^{-1}$, meaning the condition holds if and only if $p_2(X_2, Y_2) \le \gamma$.
  Since $p_2(X_2, Y_2)$ is independent of $\mathcal{I}_1$ and uniform on the grid, the conditional probability of a second error is
  \begin{equation*}
    \prob(\sigma_2 \le 2 \mid \mathcal{I}_1) = \prob(p_2(X_2, Y_2) \le \gamma \mid \mathcal{I}_1) = \gamma
    \qquad \text{almost surely on } \{\sigma_1 = 1\}.
  \end{equation*}
  Because $\gamma > \alpha$ for large enough $n$, and the conditioning event $\{\sigma_1 = 1\}$ has strictly positive probability, this strictly violates the restart-validity requirement that $\prob(\sigma_2 \le 2 \mid \mathcal{I}_1) \le \alpha$ almost surely on $\{\sigma_1 = 1\}$.
\end{proof}\par\medskip

\begin{proof}[Proof of Proposition~\ref{prop:open-loop-restart}]
  Fix $M\in\mcal$ and $s\in\{0,\dots,T-1\}$.
  Write $\tau_t$ for the predictable thresholds and $s_t := -\log(1-\tau_t)$ for the log-budget coordinate.
  Let $L := -\log(1-\alpha)$.
  Since each $s_t \ge 0$, the assumed log-budget bound $\sum_{t=1}^{T} M_t s_t \le L$ implies the tail bound $\sum_{t=s+1}^{T} M_t s_t \le L$ almost surely.

  For $t\in\{s,\dots,T\}$, let
  \begin{equation*}
    \mathcal{E}_{s:t}^{M}
    \coloneqq
    \left\{\forall u\in\{s+1,\dots,t\}\colon M_u=0\ \text{or}\ \err_u(\acal)=0\right\}
  \end{equation*}
  be the event that no monitored errors occur in the interval $(s,t]$ under mask~$M$.
  Define
  \begin{equation*}
    W_t \coloneqq \ind{\mathcal{E}_{s:t}^{M}} \exp\Big(\sum_{u=s+1}^{t} M_u s_u\Big),
    \qquad t\in\{s,\dots,T\},
  \end{equation*}
  so $W_s=1$. For every $t>s$,
  \begin{equation*}
    W_t = W_{t-1}\ind{M_t=0\ \text{or}\ p_t(X_t,Y_t)>\tau_t}\exp(M_t s_t).
  \end{equation*}
  By Proposition~\ref{prop:iid-pvalues}, conditional on $\mathcal{I}_{t-1}$ the $p$-value $p_t(X_t,Y_t)$ is super-uniform.
  Since $W_{t-1}$, $M_t$, and $s_t$ are $\mathcal{I}_{t-1}$-measurable, taking conditional expectations gives
  \begin{equation*}
    \expec[W_t \mid \mathcal{I}_{t-1}]
    \ge W_{t-1}\exp(M_t s_t)(1-\tau_t)^{M_t}
    = W_{t-1},
  \end{equation*}
  using $1-\tau_t=\exp(-s_t)$.
  Iterating the tower property yields $\expec[W_T \mid \mathcal{I}_s]\ge W_s=1$.
  By the tail log-budget bound, $W_T \le \exp(L)\ind{\mathcal{E}_{s:T}^{M}}$, hence
  \begin{equation*}
    1 \le \expec[W_T \mid \mathcal{I}_s] \le \exp(L)\prob(\mathcal{E}_{s:T}^{M}\mid \mathcal{I}_s),
  \end{equation*}
  which implies
  \begin{equation*}
    \prob\Big(\exists t>s\colon M_t=1,\ \err_t(\acal)=1 \Big| \mathcal{I}_s\Big)\le \alpha
    \qquad \text{almost surely.}
  \end{equation*}
  Taking $s=0$ proves mask validity.
  On the event $\{\sigma_j^{M}=s\}$ we have $\{\exists t>s\colon M_t=1,\ \err_t(\acal)=1\}=\{\sigma_{j+1}^{M}\le T\}$, so the conditional bound is exactly Definition~\ref{def:restart-valid}.
\end{proof}\par\medskip

\begin{proof}[Proof of Proposition~\ref{prop:geom-decay}]
  Fix a mask $M\in\mcal$.
  We prove by induction on $k\ge 1$ that $\prob(\sigma_k^{M}\le T)\le \alpha^k$.

  For $k=1$, Definition~\ref{def:restart-valid} with $j=0$ and $s=0$ gives $\prob(\sigma_1^{M}\le T)\le \alpha$.

  Now fix $k\ge 2$ and assume $\prob(\sigma_{k-1}^{M}\le T)\le \alpha^{k-1}$.
  Since $\sigma_{k-1}^{M}\in\{0,\dots,T\}\cup\{\infty\}$,
  \begin{equation*}
    \prob(\sigma_k^{M}\le T)
    =\sum_{s=0}^{T-1}\expec\Big[\prob(\sigma_k^{M}\le T \mid \mathcal{I}_s)\ind{\sigma_{k-1}^{M}=s}\Big].
  \end{equation*}
  By Definition~\ref{def:restart-valid},
  $\prob(\sigma_k^{M}\le T \mid \mathcal{I}_s)\le \alpha$
  almost surely on $\{\sigma_{k-1}^{M}=s\}$ for every $s\in\{0,\dots,T-1\}$.
  Therefore,
  \begin{equation*}
    \prob(\sigma_k^{M}\le T)
    \le \alpha\sum_{s=0}^{T-1}\prob(\sigma_{k-1}^{M}=s)
    = \alpha\prob(\sigma_{k-1}^{M}\le T)
    \le \alpha^k,
  \end{equation*}
  which closes the induction.
\end{proof}\par\medskip

\section{Deferred Results and Proofs of Section 4}
\label{app:proofs_4}

% \begin{proof}[Proof of Corollary~\ref{cor:replan-validity}]
%   Fix $M\in\mcal$ and let $\mathcal{E}_t^{M} \coloneqq \{\forall u\in\{1,\dots,t\}\colon M_u=0 \text{ or } \err_u(\acal)=0\}$ be the event that no errors occur up to time $t$.
%   Consider the random variable
%   \begin{equation*}
%     W_t \coloneqq \ind{\mathcal{E}_t^{M}} \exp\Big(\sum_{u=1}^{t} M_u s_u\Big),
%     \qquad t\in\{0,\dots,T\},
%   \end{equation*}
%   with $W_0=1$.
%   We have $W_t = W_{t-1} \ind{M_t=0 \text{ or } p_t > \tau_t} \exp(M_t s_t)$.
%   By Proposition~\ref{prop:iid-pvalues}, conditional on $\mathcal{I}_{t-1}$, the $p$-value $p_t$ is super-uniform.
%   Since $W_{t-1}$, $M_t$, and $s_t$ are $\mathcal{I}_{t-1}$-measurable, taking the conditional expectation yields
%   \begin{equation*}
%     \expec[W_t \mid \mathcal{I}_{t-1}]
%     \ge W_{t-1} \exp(M_t s_t) (1-\tau_t)^{M_t}
%     = W_{t-1},
%   \end{equation*}
%   because $1-\tau_t = \exp(-s_t)$.
%   Applying the tower property iteratively gives $\expec[W_T] \ge \expec[W_0] = 1$.

%   By the log-budget constraint, $\sum_{t=1}^{T} M_t s_t \le L$ almost surely.
%   Therefore,
%   \begin{equation*}
%     1 \le \expec[W_T] \le \expec\Big[\ind{\mathcal{E}_T^{M}} \exp(L)\Big] = \exp(L) \prob(\mathcal{E}_T^{M}),
%   \end{equation*}
%   which implies $\prob(\mathcal{E}_T^{M}) \ge \exp(-L) = 1-\alpha$, and thus $\prob(F_T(M,\acal)) \le \alpha$, proving mask validity.

%   For suffix-validity, fix any $s\in\{0,\dots,T-1\}$ and apply the same argument to the tail event
%   \begin{equation*}
%     \mathcal{E}_{>s}^{M}
%     \coloneqq
%     \left\{\forall t>s\colon M_t=0\ \text{or}\ \err_t(\acal)=0\right\},
%   \end{equation*}
%   working conditionally on $\mathcal{I}_s$ and using the tail log-budget bound $\sum_{t=s+1}^{T} M_t s_t\le L$.
%   This gives
%   \begin{equation*}
%     \prob\Big(\exists t>s\colon M_t=1,\ \err_t(\acal)=1 \Big| \mathcal{I}_s\Big)\le \alpha.
%   \end{equation*}
%   Restart-validity follows by restricting this conditional bound to $\{\sigma_j^{M}=s\}$ and using $\{\exists t>s\colon M_t=1,\ \err_t(\acal)=1\}=\{\sigma_{j+1}^{M}\le T\}$ on that event, as in Definition~\ref{def:restart-valid}.
% \end{proof}\par\medskip

\subsection{Notations, Preliminaries, Assumptions, and Duality}
\label{app:replanning}

This appendix provides the formal definitions and complete proofs underlying the suboptimality bounds for the replanning policy.
We begin by establishing the formal notation, assumptions, and duality properties.
Then we state the general suboptimality theorem, and finally provide a specialized bound for the all-ones mask.

\subsubsection{Notation and Setup}
\label{app:replanning-prelim}

We work in the log-budget coordinate $s_t := -\log(1-\tau_t) \in [0, s_{\max}]$, where $s_{\max} := -\log(1-\tau_{\max})$.
Conceptually, our framework assigns an initial log-budget $L$ to every mask $M \in \mcal$. At each timestep $u$, a scheduling policy allocates a log-budget $s_u$, which is subtracted from the available budget of any mask $M$ that covers timestep $u$ (i.e., when $M_u = 1$).
Therefore, after allocating $s_1, \dots, s_{t-1}$, the remaining budget at time $t$ for a mask $M \in \mcal$ is $b_t(M) := L - \sum_{u < t} M_u s_u$.
Let $\mathbf{b}_t := (b_t(M))_{M \in \mcal} \in \R^{|\mcal|}$ denote the state vector tracking the remaining budgets for all masks, with initial state $\mathbf{b}_1 = L\1$, where $\1$ is the all-ones vector.
The mask-incidence direction at time $t$ is the vector $\mathbf{m}_t := (M_t)_{M \in \mcal} \in \{0,1\}^{|\mcal|}$, yielding the state transition $\mathbf{b}_{t+1} = \mathbf{b}_t - s_t \mathbf{m}_t$.
For $\mathbf{b} \in \R_{\ge 0}^{|\mcal|}$, the one-step feasible set is the interval $\scalt(\mathbf{b}) := [0, \min\{s_{\max}, \inf_{M\in\mcal: M_t = 1} b(M)\}]$, which prevents any active mask's budget from becoming negative.

We define the remaining offline value function as
\begin{equation}
  \label{eq:value-function}
  V_t(\mathbf{b})
  := \min_{\substack{0 \le s_{t:T} \le s_{\max} \\
  \sum_{u=t}^{T} M_u s_u \le b(M),\ \forall M \in \mcal}}
  \sum_{u=t}^{T} \tilde{c}_u(s_u),
\end{equation}
with terminal condition $V_{T+1}(\cdot) = 0$.
Similarly, for each planning time $t$, the forecasted value function $\hat{V}_{u \mid t}(\mathbf{b})$ is defined by replacing $\tilde{c}_v$ with the forecast $\hat{c}_{t,v}$ in \eqref{eq:value-function} for stages $v \ge u$, with $\hat{V}_{T+1 \mid t}(\cdot)=0$.
Under these definitions, the replanning action from Equation~\eqref{eq:replan-action} can be equivalently written as
\begin{equation}
  \label{eq:replan-action-value}
  s_t \in \argmin_{s \in \scalt(\mathbf{b}_t)} \left[ \hat{c}_{t,t}(s) + \hat{V}_{t+1 \mid t}(\mathbf{b}_t - s\mathbf{m}_t) \right].
\end{equation}

The realized costs $\tilde{c}_u$, the forecasts $\hat{c}_{t,u}$, and the replanning policy follow Section~\ref{sec:finite}.
We quantify forecast accuracy via the uniform derivative forecast error
\begin{equation*}
  e_{t,u} := \sup_{s\in[0,s_{\max}]} \bigl|\tilde{c}_u'(s) - \hat{c}_{t,u}'(s)\bigr|,
  \qquad
  e_t := e_{t,t}.
\end{equation*}

Define the tail-interior region
\begin{equation*}
  \bcal_{t}^{\circ} := \left\{\mathbf{b}\in[0,L]^{|\mcal|} : b(M)>0, \forall M\in\mcal \text{ such that } \sum_{u=t}^{T} M_u>0 \right\}.
\end{equation*}
This set represents the interior of the budget: it encompasses all states where no active mask has fully exhausted its budget allocation.

\subsubsection{Standing Assumptions}

We require formal regularity conditions on the costs and the realized replanning trajectory.

\begin{assumptions}
\item\label{ass:A1} (Strong convexity and smoothness.) There exist constants $0 < \mu \le \beta$ such that, for each $u$, the realized cost $\tilde{c}_u$ is continuously differentiable on $[0,s_{\max}]$, twice continuously differentiable on $(0,s_{\max})$, and satisfies $\mu \le \tilde{c}_u''(s) \le \beta$. The forecast curves $\hat{c}_{t,u}$ satisfy identical bounds. Write the condition number $\kappa := \beta/\mu$.
\item\label{ass:A2} (Trajectory interiority.)
  There exists a margin $\delta > 0$ such that, along the realized trajectory $\{(\mathbf{b}_t,s_t)\}_{t=1}^{T}$ generated by the replanning policy, for each $t \in \{1,\dots,T\}$, both the replanning action and the offline-optimal stage action lie in the interior of the one-step feasible interval:
  \begin{equation*}
    s_t,\ s_t^{\star}
    \in
    [\delta,\ \sup\scalt(\mathbf{b}_t)-\delta].
  \end{equation*}
  Furthermore, for each $t \in \{1,\dots,T-1\}$, every coordinate of the primal optima of both $V_{t+1}(\mathbf{b}_{t+1})$ and $\hat{V}_{t+1\mid t}(\mathbf{b}_{t+1})$ lies in $[\delta, s_{\max} - \delta]$.
\item\label{ass:A3} (Decreasing costs.)
  For each $t \in \{1,\dots,T\}$, $\tilde{c}_t'(s) < 0$ and $\hat{c}_{t,u}'(s) < 0$ for all $s \in [0,s_{\max}]$.
\end{assumptions}

The interiority condition (\ref{ass:A2}) ensures that the budget is neither entirely depleted nor perfectly abundant along the path actually taken by the algorithm, avoiding degenerate boundary constraints.
The decreasing costs condition (\ref{ass:A3}) reflects the practical reality that larger thresholds (which consume more log-budget) result in smaller prediction sets and lower costs.
Note that the strong convexity in \ref{ass:A1} inherently guarantees that the optimization problems are well-posed convex programs, ensures unique primal optima, are sufficient for strong duality.

\begin{remark}
  We provide a concrete example from batch-sequential conformal prediction satisfying these assumptions.
  Let the expected prediction set size be proportional to $(1-\tau)^{\gamma}$ for constants $C > 0$ and $\gamma > 0$.
  The log-transformed cost is $\tilde{c}_{t}(s) = c_t(1-e^{-s}) = C e^{-\gamma s}$.
  The derivatives are $\tilde{c}_{t}'(s) = -C\gamma e^{-\gamma s}$ and $\tilde{c}_{t}''(s) = C\gamma^2 e^{-\gamma s}$.
  The cost is strictly decreasing, satisfying \ref{ass:A3}.
  The domain $s \in [0, s_{\max}]$ is bounded, so the second derivative is uniformly bounded by $\mu = C\gamma^2 e^{-\gamma s_{\max}}$ and $\beta = C\gamma^2$.
  This establishes strong convexity and smoothness (\ref{ass:A1}).

  The strictly decreasing costs (\ref{ass:A3}) also relate to the interiority condition (\ref{ass:A2}).
  Because marginal costs $\tilde{c}_{t}'(s)$ strictly increase with $s$, an optimal policy avoids extreme single-stage allocations that would force other stages to $0$.
  For homogeneous cost curves, this forces the allocation into the interior of the feasible domain.
\end{remark}

\subsubsection{Lagrangian Duality}
\label{subsec:duality}

In our results, we exploit the duality structure of $V_t$ and its forecasted approximation. In particular, let $\Lambda_t$ be the cone of nonnegative multipliers $\boldsymbol{\lambda} \in \R_{\ge 0}^{|\mcal|}$ where $\lambda(M)=0$ if mask $M$ is inactive in the tail ($\sum_{u=t}^T M_u = 0$).
The tail Lagrangian is
\begin{equation*}
  \mathcal{L}_t(\mathbf{s},\boldsymbol{\lambda};\mathbf{b})
  :=
  \sum_{u=t}^{T} \tilde{c}_u(s_u)
  + \boldsymbol{\lambda}^\top (\mathbf{M}_{t:T} \mathbf{s} - \mathbf{b}),
  \qquad \mathbf{s} \in [0,s_{\max}]^{T-t+1},
\end{equation*}
where $\mathbf{M}_{t:T} := [\mathbf{m}_t, \dots, \mathbf{m}_T] \in \{0,1\}^{|\mcal| \times (T-t+1)}$ is the mask-incidence matrix whose columns are the incidence directions for the tail stages, such that its entry for mask $M$ and stage $u$ is $M_u$.
The dual function is $\mathcal{D}_t(\boldsymbol{\lambda};\mathbf{b}) := \min_{\mathbf{s}} \mathcal{L}_t(\mathbf{s},\boldsymbol{\lambda};\mathbf{b})$, and the dual-optimal set is $D_t(\mathbf{b}) := \argmax_{\boldsymbol{\lambda} \in \Lambda_t} \mathcal{D}_t(\boldsymbol{\lambda};\mathbf{b})$.

In the following sequence of results, we explicitly write the proofs for the exact value function $V_t$ and the realized costs $\tilde{c}$. Because the forecasted costs $\hat{c}$ satisfy identical structural conditions under Standing Assumptions \ref{ass:A1}--\ref{ass:A3}, every lemma holds symmetrically for the forecasted tail program $\hat{V}_{u \mid t}$. For clarity, we state both results in the lemma but omit the completely redundant proof for the forecast.
The auxiliary lemmas below (Lemmas~\ref{lem:budget-tight-app}, \ref{lem:interior-optimality}, and~\ref{lem:dual-unique-app}) record tightness, KKT, and uniqueness properties that are used implicitly throughout the subsequent sensitivity arguments.

\begin{lemma}[At least one budget constraint is tight]
  \label{lem:budget-tight-app}
  Under \ref{ass:A2} and \ref{ass:A3}, fix $t \in \{1,\dots,T-1\}$ and let $\mathbf{b}=\mathbf{b}_{t+1}\in\bcal_{t+1}^{\circ}$ be the realized remaining budget vector.
  Then at least one tail-active mask constraint is satisfied with strict equality at every primal optimum of $V_{t+1}(\mathbf{b})$ and $\hat{V}_{t+1\mid t}(\mathbf{b})$.
  That is, there exists a mask $M\in\mcal$ with $\sum_{u=t+1}^{T} M_u>0$ such that $\sum_{u=t+1}^{T} M_u s_u=b(M)$.
\end{lemma}
\begin{proof}
  By the trajectory interiority assumption (\ref{ass:A2}), the primal optima $\mathbf{s}^\star$ and $\hat{\mathbf{s}}$ lie strictly in $(0, s_{\max})^{T-t}$. If all active budget constraints were strictly loose (i.e., $\sum_{u=t+1}^{T} M_u s_u < b(M)$ for all active $M$), one could strictly increase every $s_u$ by some infinitesimal amount $\varepsilon > 0$ without violating any budget constraint or the upper bound $s_{\max}$. Since the cost functions are strictly decreasing (\ref{ass:A3}), this uniform perturbation would strictly decrease the objective value, which contradicts the optimality of $\mathbf{s}^\star$. Therefore, at least one active constraint must bind.
\end{proof}

\begin{lemma}[Existence and strong duality]
  \label{lem:slater-attainment}
  Under \ref{ass:A1}, for $\mathbf{b} \in \bcal_t^\circ$, Slater's condition holds for the offline tail program. Strong duality holds, and the dual optimum $D_t(\mathbf{b})$ is non-empty.
\end{lemma}
\begin{proof}
  Since $b(M)>0$ for all active masks, setting $s_u = \epsilon$ for a sufficiently small $\epsilon>0$ strictly satisfies all active linear constraints. The result follows by standard convex duality for compact feasible sets.
\end{proof}

\begin{lemma}[First-order optimality]
  \label{lem:interior-optimality}
  Under \ref{ass:A1} and \ref{ass:A2}, fix $t \in \{1,\dots,T\}$ and $\mathbf{b} \in \bcal_t^\circ$.
  Let $\mathbf{s}^\star = (s_v^\star)_{v=t}^{T}$ be a primal optimizer attaining $V_t(\mathbf{b})$, and fix any $u \in \{t,\dots,T\}$.
  If the coordinate $s_u^\star \in (0, s_{\max})$, then
  \begin{equation*}
    \tilde{c}_u'(s_u^\star) = - \sum_{M \in \mcal} M_{u} \lambda^\star(M)
  \end{equation*}
  for any $\boldsymbol{\lambda}^\star \in D_t(\mathbf{b})$.

  Similarly, fix a planning time $t \in \{1,\dots,T\}$, a tail-start index $u \in \{1,\dots,T\}$, and $\mathbf{b} \in \bcal_u^\circ$.
  Let $\hat{\mathbf{s}} = (\hat{s}_v)_{v=u}^{T}$ be a primal optimizer attaining $\hat{V}_{u \mid t}(\mathbf{b})$.
  If the coordinate $\hat{s}_u \in (0, s_{\max})$, then
  \begin{equation*}
    \hat{c}_{t,u}'(\hat{s}_u) = - \sum_{M \in \mcal} M_{u} \hat{\lambda}^\star(M)
  \end{equation*}
  for any $\hat{\boldsymbol{\lambda}}^\star \in \hat{D}_{u \mid t}(\mathbf{b})$.
\end{lemma}
\begin{proof}
  By restricted strong duality (Lemma~\ref{lem:slater-attainment}), the primal-dual pair forms a saddle point of the tail Lagrangian.
  Equivalently, the optimal sequence $\mathbf{s}^\star$ solves the decoupled minimization of the Lagrangian over the box constraints $0 \le s_u \le s_{\max}$.
  Introducing nonnegative dual multipliers $\nu_u^-$ for the lower bounds ($-s_u \le 0$) and $\nu_u^+$ for the upper bounds ($s_u - s_{\max} \le 0$), the Karush--Kuhn--Tucker (KKT) stationarity condition states:
  \begin{equation}
    \label{eq:kkt-stationarity}
    \tilde{c}_u'(s_u^\star) + \sum_{M \in \mcal} M_{u} \lambda^\star(M) - \nu_u^- + \nu_u^+ = 0.
  \end{equation}
  By hypothesis, the optimal coordinate $s_u^\star$ lies strictly in the interior of its domain. The complementary slackness conditions ($\nu_u^- s_u^\star = 0$ and $\nu_u^+(s_u^\star - s_{\max}) = 0$) therefore dictate that the bound multipliers must be exactly zero ($\nu_u^- = \nu_u^+ = 0$). Substituting these into \eqref{eq:kkt-stationarity} yields the exact relation. An identical KKT argument applies to the forecasted sequence.
\end{proof}

\begin{lemma}[Primal-dual uniqueness and strict complementarity]
  \label{lem:dual-unique-app}
  Assume \ref{ass:A1}--\ref{ass:A3}. Fix $t \in \{1,\dots,T-1\}$ and let $\mathbf{b} = \mathbf{b}_{t+1} \in \bcal_{t+1}^\circ$ be the realized remaining budget vector.
  Then the tail program $V_{t+1}(\mathbf{b})$ possesses a unique primal optimum $\mathbf{s}^{\star} \in (0,s_{\max})^{T-t}$.
  If the mask incidence vectors corresponding to the active budget constraints are linearly independent, the associated optimal dual multiplier vector $\boldsymbol{\lambda}(\mathbf{b}) \in \Lambda_{t+1}$ is unique.
  Furthermore, the inner product $\mathbf{M}_{t+1:T}^\top \boldsymbol{\lambda}(\mathbf{b})$ is strictly positive component-wise. If exactly one mask constraint is binding, its associated dual multiplier is strictly positive.
  Identical properties hold for the forecasted tail program $\hat{V}_{t+1\mid t}(\mathbf{b})$, yielding a unique primal restricted action $\hat{\mathbf{s}} \in (0,s_{\max})^{T-t}$ and unique dual vector $\hat{\boldsymbol{\lambda}}(\mathbf{b}) \in \Lambda_{t+1}$ satisfying $\mathbf{M}_{t+1:T}^\top \hat{\boldsymbol{\lambda}}(\mathbf{b}) > 0$.
\end{lemma}
\begin{proof}
  We study the tail optimization programs $V_{t+1}(\mathbf{b})$ and $\hat{V}_{t+1\mid t}(\mathbf{b})$. The objectives are strongly convex (\ref{ass:A1}), which guarantees unique primal optima.
  By the interiority assumption (\ref{ass:A2}), these optima lie strictly in the interior $(0, s_{\max})^{T-t}$. Therefore, tight constraints arise exclusively from the mask limits $b(M)$, meaning the box-constraint multipliers are zero.
  The KKT stationarity condition reduces to $\nabla \tilde{c}(\mathbf{s}^\star) + \mathbf{M}_{t+1:T}^\top \boldsymbol{\lambda} = 0$. By complementary slackness, $\lambda(M) = 0$ for any non-binding mask.
  The binding dual multipliers are determined by this linear subsystem. When the rows of $\mathbf{M}_{t+1:T}$ corresponding to binding masks are linearly independent, the linear subsystem has a unique solution, making the mapping from the primal $\mathbf{s}^\star$ to the dual $\boldsymbol{\lambda}(\mathbf{b})$ unique.
  Because the stage costs are strictly decreasing (\ref{ass:A3}), we have $\tilde{c}_u' < 0$. The stationarity condition then implies $\mathbf{M}_{t+1:T}^\top \boldsymbol{\lambda} = -\nabla \tilde{c}(\mathbf{s}^\star) > 0$ component-wise.
\end{proof}

\begin{lemma}[Value-Function Subgradients]
  \label{lem:dual-subgrad}
  For $\mathbf{b} \in \bcal_t^\circ$, the exact value function satisfies $\subdiff V_t(\mathbf{b}) = \{-\boldsymbol{\lambda} : \boldsymbol{\lambda} \in D_t(\mathbf{b})\}$.
  Similarly, for the forecasted value function, $\subdiff \hat{V}_{u\mid t}(\mathbf{b}) = \{-\hat{\boldsymbol{\lambda}} : \hat{\boldsymbol{\lambda}} \in \hat{D}_{u\mid t}(\mathbf{b})\}$.
\end{lemma}
\begin{proof}
  We state the proof for the exact value function; the argument for the forecast is identical. The tail program defines $V_t(\mathbf{b})$ as the optimal value of a convex minimization problem parameterized by right-hand side perturbations $\mathbf{b}$.
  Under Slater's condition, which holds for $\mathbf{b} \in \bcal_t^\circ$ by Lemma~\ref{lem:slater-attainment}, a standard result in convex programming duality states that the subdifferential of the optimal value function with respect to these right-hand side constraints is exactly the negative of the dual optimal set.
  Since $D_t(\mathbf{b})$ is convex, this set is equal to its convex hull.
\end{proof}

\subsection{The Replanning Suboptimality Theorem}
\label{sec:suboptimality-theorem}

We quantify the performance loss of replanning relative to the offline optimum for the realized costs.
The bound is controlled by a stage-cost forecasting error and a directional error in the forecasted tail value function.
Recall from Appendix~\ref{app:replanning-prelim} that $\mathbf{m}_t := (M_t)_{M \in \mcal} \in \{0,1\}^{|\mcal|}$ is the mask-incidence direction at time $t$.

Define the directional dual-multiplier error
\begin{equation}
  \label{eq:eps-def-app}
  \epsilon_t(\mathbf{b})
  :=
  \sup_{\substack{\hat{\boldsymbol{\lambda}}\in\hat{D}_{t+1\mid t}(\mathbf{b})\\ \boldsymbol{\lambda}\in D_{t+1}(\mathbf{b})}}
  \left| (\hat{\boldsymbol{\lambda}}-\boldsymbol{\lambda})^\top\mathbf{m}_t \right|,
  \qquad \mathbf{b}\in\bcal_{t+1}^{\circ},
\end{equation}
and let $\epsilon_t := \sup_{\mathbf{b}\in\bcal_{t+1}^{\circ}} \epsilon_t(\mathbf{b})$.

\begin{theorem}[Closed-loop suboptimality bound]
  \label{thm:replan-regret}
  Assume each realized cost curve $\tilde{c}_t$ and each forecast curve $\hat{c}_{t,u}$ is convex and continuously differentiable on $[0,s_{\max}]$, and finite on an open neighborhood of $[0,s_{\max}]$.
  Assume also that the replanning trajectory satisfies $\mathbf{b}_{t+1}\in\bcal_{t+1}^{\circ}$ for all $t\in\{1,\dots,T\}$.
  Let $(s_t)_{t=1}^{T}$ be generated by \eqref{eq:replan-action} and define the suboptimality gap
  \begin{equation*}
    \mathrm{Gap} := \sum_{t=1}^{T} \tilde{c}_t(s_t) - V_1(L\1).
  \end{equation*}
  Then $\mathrm{Gap}=\sum_{t=1}^{T} \delta_t$, where $\delta_t:=\tilde{c}_t(s_t)+V_{t+1}(\mathbf{b}_{t+1})-V_t(\mathbf{b}_t)$, and
  \begin{equation*}
    0 \le \delta_t \le \left(\epsilon_t(\mathbf{b}_{t+1})+e_t\right)|s_t-s_t^{\star}|
    \le (\epsilon_t+e_t)|s_t-s_t^{\star}|,
  \end{equation*}
  where $s_t^{\star}\in\argmin_{s\in\scalt(\mathbf{b}_t)}[\tilde{c}_t(s)+V_{t+1}(\mathbf{b}_t-s\mathbf{m}_t)]$ denotes the optimal action for the true costs.

  If, in addition, each $\tilde{c}_t$ is $\mu$-strongly convex on $[0,s_{\max}]$ for some $\mu>0$, then
  \begin{equation*}
    \mathrm{Gap} \le \frac{1}{2\mu}\sum_{t=1}^{T} (\epsilon_t+e_t)^2.
  \end{equation*}
\end{theorem}

\begin{proof}[Proof of Theorem~\ref{thm:replan-regret}]
  We first recall why $\scalt(\mathbf{b})$ is the correct one-step feasible interval.
  Fix $t$ and $\mathbf{b}\ge 0$ coordinatewise.
  If $0\le s\le \min\{s_{\max},\inf_{M\in\mcal:M_t=1} b(M)\}$, then for every mask $M$ with $M_t=1$ we have $s\le b(M)$, and for every mask with $M_t=0$ we have $M_t s = 0\le b(M)$.
  Taking the completion $s_{t+1}=\cdots=s_T=0$ shows that $s\in\scalt(\mathbf{b})$ indeed admits a feasible tail completion.
  Conversely, if $s>\min\{s_{\max},\inf_{M\in\mcal:M_t=1} b(M)\}$, then either $s>s_{\max}$ or there exists a mask $M$ with $M_t=1$ and $b(M)<s$.
  In either case, feasibility fails for any completion with $s_{t+1:T}\ge 0$ because $\sum_{u=t}^{T} M_u s_u \ge M_t s = s > b(M)$.

  \paragraph{Replaning recursion and telescoping.}
  Define the one-step cost-to-go
  \begin{equation*}
    h_t(s;\mathbf{b})
    \coloneqq
    \tilde{c}_t(s)+V_{t+1}(\mathbf{b}-s\mathbf{m}_t).
  \end{equation*}
  Splitting the offline tail program at time $t$ yields the dynamic programming recursion
  \begin{equation*}
    V_t(\mathbf{b})=\min_{s\in\scalt(\mathbf{b})} h_t(s;\mathbf{b}).
  \end{equation*}
  Therefore, writing $\mathbf{b}_{t+1}=\mathbf{b}_t-s_t\mathbf{m}_t$, the one-step suboptimality is
  \begin{equation*}
    \delta_t
    =
    h_t(s_t;\mathbf{b}_t)-\min_{s\in\scalt(\mathbf{b}_t)}h_t(s;\mathbf{b}_t)
    \ge 0.
  \end{equation*}
  Summing $\delta_t=\tilde{c}_t(s_t)+V_{t+1}(\mathbf{b}_{t+1})-V_t(\mathbf{b}_t)$ over $t=1,\dots,T$ telescopes:
  \begin{equation*}
    \sum_{t=1}^{T}\delta_t
    =
    \sum_{t=1}^{T}\tilde{c}_t(s_t)
    +
    \sum_{t=1}^{T}\left[V_{t+1}(\mathbf{b}_{t+1})-V_t(\mathbf{b}_t)\right]
    =
    \sum_{t=1}^{T}\tilde{c}_t(s_t)-V_1(L\1),
  \end{equation*}
  since $V_{T+1}(\cdot)=0$ and $\mathbf{b}_1=L\1$.

  \paragraph{One-step suboptimality bound.}
  Fix $t$.
  Let $s_t^{\star}\in\argmin_{s\in\scalt(\mathbf{b}_t)}h_t(s;\mathbf{b}_t)$ denote an optimal stage action for the true costs.
  By the subgradient inequality, for any $g\in\partial h_t(s_t;\mathbf{b}_t)$,
  \begin{equation*}
    h_t(s_t^{\star};\mathbf{b}_t)\ge h_t(s_t;\mathbf{b}_t)+g(s_t^{\star}-s_t),
  \end{equation*}
  hence $\delta_t=h_t(s_t;\mathbf{b}_t)-h_t(s_t^{\star};\mathbf{b}_t)\le g(s_t-s_t^{\star})$.

  Define the forecasted cost-to-go
  \begin{equation*}
    \hat{h}_t(s;\mathbf{b})
    \coloneqq
    \hat{c}_{t,t}(s)+\hat{V}_{t+1\mid t}(\mathbf{b}-s\mathbf{m}_t).
  \end{equation*}
  Since $s_t$ minimizes $\hat{h}_t(\cdot;\mathbf{b}_t)$ over the interval $\scalt(\mathbf{b}_t)$, the first-order optimality condition for convex minimization over a convex set implies that there exists $\hat{g}\in\partial \hat{h}_t(s_t;\mathbf{b}_t)$ such that
  \begin{equation*}
    \hat{g}(s-s_t)\ge 0,\qquad \forall s\in\scalt(\mathbf{b}_t).
  \end{equation*}
  In particular, taking $s=s_t^{\star}$ yields $\hat{g}(s_t^{\star}-s_t)\ge 0$, i.e., $\hat{g}(s_t-s_t^{\star})\le 0$.
  Therefore,
  \begin{equation*}
    \delta_t
    \le g(s_t-s_t^{\star})
    = (g-\hat{g})(s_t-s_t^{\star})+\hat{g}(s_t-s_t^{\star})
    \le |g-\hat{g}|\cdot |s_t-s_t^{\star}|.
  \end{equation*}

  It remains to bound $|g-\hat{g}|$.
  By the subdifferential sum rule and affine chain rule applied at $s_t$ (using that all functions are convex and finite on a neighborhood of $[0,s_{\max}]$ and $\mathbf{b}_{t+1}\in\bcal_{t+1}^{\circ}$), every $g\in\partial h_t(s_t;\mathbf{b}_t)$ admits a decomposition
  \begin{equation*}
    g = g_c - g_V^{\top}\mathbf{m}_t,
    \qquad
    g_c\in\partial \tilde{c}_t(s_t),
    \qquad
    g_V\in\partial V_{t+1}(\mathbf{b}_{t+1}),
  \end{equation*}
  and similarly every $\hat{g}\in\partial \hat{h}_t(s_t;\mathbf{b}_t)$ can be written as
  \begin{equation*}
    \hat{g} = \hat{g}_c - \hat{g}_V^{\top}\mathbf{m}_t,
    \qquad
    \hat{g}_c\in\partial \hat{c}_{t,t}(s_t),
    \qquad
    \hat{g}_V\in\partial \hat{V}_{t+1\mid t}(\mathbf{b}_{t+1}).
  \end{equation*}
  Therefore,
  \begin{equation*}
    |g-\hat{g}|
    \le |\hat{g}_c-g_c| + |(\hat{g}_V-g_V)^{\top}\mathbf{m}_t|
    \le \bigl|\hat{c}_{t,t}'(s_t)-\tilde{c}_t'(s_t)\bigr|+\epsilon_t(\mathbf{b}_{t+1})
    \le e_t+\epsilon_t(\mathbf{b}_{t+1}),
  \end{equation*}
  where the last inequality uses the definition of $e_t$ together with Lemma~\ref{lem:dual-subgrad}.
  Substituting back gives
  \begin{equation*}
    0\le \delta_t \le \left(e_t+\epsilon_t(\mathbf{b}_{t+1})\right)|s_t-s_t^{\star}|
    \le (e_t+\epsilon_t)|s_t-s_t^{\star}|.
  \end{equation*}

  \paragraph{Strong convexity refinement.}
  Now assume each $\tilde{c}_t$ is $\mu$-strongly convex on $[0,s_{\max}]$.
  Since $s\mapsto V_{t+1}(\mathbf{b}_t-s\mathbf{m}_t)$ is convex as a composition of a convex function with an affine map, $h_t(\cdot;\mathbf{b}_t)$ is also $\mu$-strongly convex on $\scalt(\mathbf{b}_t)$.
  Fix any $g\in\partial h_t(s_t;\mathbf{b}_t)$ and write $d\coloneqq |s_t-s_t^{\star}|$.
  Strong convexity implies
  \begin{equation*}
    h_t(s_t^{\star};\mathbf{b}_t)
    \ge h_t(s_t;\mathbf{b}_t)+g(s_t^{\star}-s_t)+\frac{\mu}{2}d^2,
  \end{equation*}
  hence
  \begin{equation*}
    \delta_t
    \le g(s_t-s_t^{\star})-\frac{\mu}{2}d^2.
  \end{equation*}
  Using $g=\hat{g}+(g-\hat{g})$ and $\hat{g}(s_t-s_t^{\star})\le 0$ as in Step~2 gives
  \begin{equation*}
    \delta_t
    \le |g-\hat{g}|\cdot d-\frac{\mu}{2}d^2.
  \end{equation*}
  Let $A_t\coloneqq |g-\hat{g}|\le e_t+\epsilon_t(\mathbf{b}_{t+1})$.
  Maximizing the quadratic upper bound $A_t d-(\mu/2)d^2$ over $d\ge 0$ yields
  \begin{equation*}
    \delta_t
    \le \frac{A_t^2}{2\mu}
    \le \frac{1}{2\mu}\left(e_t+\epsilon_t(\mathbf{b}_{t+1})\right)^2
    \le \frac{1}{2\mu}(\epsilon_t+e_t)^2.
  \end{equation*}
  Summing over $t$ gives the strong-convexity bound on $\mathrm{Gap}$.
\end{proof}

\subsection{Analysis of the All-Ones Mask}
\label{app:dual-bounds}

We now prove the multiplier error bound for the all-ones mask $\mcal=\{\1\}$.
The state is the scalar remaining budget $b_t := L - \sum_{u<t} s_u$,
and the feasible set is $\scalt(b_t) = [0,\min\{b_t,s_{\max}\}]$.
Write $N := T-t$.
Let $\lambda(b)$ and $\hat{\lambda}(b)$ denote the (scalar) Lagrange multipliers associated with the tail budget constraint at the optimizers of $V_{t+1}(b)$ and $\hat{V}_{t+1\mid t}(b)$, respectively.

\begin{theorem}[Multiplier error bound for the all-ones mask]
  \label{thm:epsilon-allones-app}%
  Assume Standing Assumptions~\ref{ass:A1}--\ref{ass:A3}.
  For $t \in \{1,\dots,T-1\}$ and $b = b_{t+1} > 0$,
  \begin{equation}
    \label{eq:epsilon-allones-app}
    \begin{aligned}
      \epsilon_t(b)
      &= |\lambda(b)-\hat{\lambda}(b)|\\
      &\le \kappa\bar{e}_{t,>t},\\
      \bar{e}_{t,>t}
      &:= \frac{1}{N}\sum_{u=t+1}^{T} e_{t,u}.
    \end{aligned}
  \end{equation}
\end{theorem}

\begin{proof}
  Fix $t \in \{1,\dots,T-1\}$ and $b = b_{t+1} > 0$.
  Let $\mathbf{s}^\star = (s_u^\star)_{u=t+1}^{T}$ and $\hat{\mathbf{s}} = (\hat{s}_u)_{u=t+1}^{T}$ denote the unique primal optimizers attaining $V_{t+1}(b)$ and $\hat{V}_{t+1\mid t}(b)$, respectively, which are unique by \ref{ass:A1}.
  By the interiority assumption \ref{ass:A2}, these optima satisfy $s_u^\star \in (0,s_{\max})$ and $\hat{s}_u \in (0,s_{\max})$ for all $u$, so the box constraints are inactive.

  We first note that the all-ones budget constraint is tight for both tail problems:
  \begin{equation*}
    \sum_{u=t+1}^{T} s_u^\star = b,
    \qquad
    \sum_{u=t+1}^{T} \hat{s}_u = b.
  \end{equation*}
  Indeed, if $\sum_{u=t+1}^{T} s_u^\star < b$, then for sufficiently small $\varepsilon>0$ the uniform perturbation $s_u^\star+\varepsilon$ remains feasible (since $s_u^\star<s_{\max}$ for all $u$) and strictly decreases the objective because $\tilde{c}_u'(s)<0$ for all $s \in [0,s_{\max}]$ by \ref{ass:A3}. This contradicts optimality. The forecasted problem is identical.

  By KKT stationarity for the single (tight) budget constraint, there exist multipliers $\lambda(b)$ and $\hat{\lambda}(b)$ such that
  \begin{equation*}
    \tilde{c}_u'(s_u^\star)+\lambda(b)=0,
    \qquad
    \hat{c}_{t,u}'(\hat{s}_u)+\hat{\lambda}(b)=0,
    \qquad
    \forall u \in \{t+1,\dots,T\}.
  \end{equation*}
  In particular, since $\tilde{c}_u'(s_u^\star)<0$ and $\hat{c}_{t,u}'(\hat{s}_u)<0$ by \ref{ass:A3}, we have $\lambda(b)>0$ and $\hat{\lambda}(b)>0$.
  Because $\mathbf{s}^\star$ and $\hat{\mathbf{s}}$ are unique and the stationarity equations identify the multipliers, the associated dual sets are singletons, hence $\epsilon_t(b)=|\lambda(b)-\hat{\lambda}(b)|$.

  For each $u \in \{t+1,\dots,T\}$, define
  \begin{equation*}
    \Delta_u \coloneqq \tilde{c}_u'(\hat{s}_u)-\hat{c}_{t,u}'(\hat{s}_u).
  \end{equation*}
  By definition of $e_{t,u}$, we have $|\Delta_u| \le e_{t,u}$.
  Subtracting the stationarity conditions and adding/subtracting $\tilde{c}_u'(\hat{s}_u)$ yields
  \begin{equation*}
    \tilde{c}_u'(s_u^\star)-\tilde{c}_u'(\hat{s}_u)
    =
    -\left(\lambda(b)-\hat{\lambda}(b)\right)-\Delta_u.
  \end{equation*}
  By the mean value theorem, there exists $\xi_u$ between $s_u^\star$ and $\hat{s}_u$ such that
  \begin{equation*}
    \tilde{c}_u''(\xi_u)(s_u^\star-\hat{s}_u)
    =
    -\left(\lambda(b)-\hat{\lambda}(b)\right)-\Delta_u.
  \end{equation*}
  Summing over $u=t+1,\dots,T$ and using $\sum_{u=t+1}^{T}(s_u^\star-\hat{s}_u)=0$ gives
  \begin{equation*}
    \left(\lambda(b)-\hat{\lambda}(b)\right)\sum_{u=t+1}^{T}\frac{1}{\tilde{c}_u''(\xi_u)}
    =
    -\sum_{u=t+1}^{T}\frac{\Delta_u}{\tilde{c}_u''(\xi_u)}.
  \end{equation*}
  Using \ref{ass:A1}, we have $\mu \le \tilde{c}_u''(\xi_u) \le \beta$, hence
  \begin{equation*}
    \left|\sum_{u=t+1}^{T}\frac{\Delta_u}{\tilde{c}_u''(\xi_u)}\right|
    \le
    \frac{1}{\mu}\sum_{u=t+1}^{T}|\Delta_u|
    \le
    \frac{1}{\mu}\sum_{u=t+1}^{T}e_{t,u},
    \qquad
    \sum_{u=t+1}^{T}\frac{1}{\tilde{c}_u''(\xi_u)}
    \ge
    \frac{N}{\beta}.
  \end{equation*}
  Therefore,
  \begin{equation*}
    |\lambda(b)-\hat{\lambda}(b)|
    \le
    \frac{(1/\mu)\sum_{u=t+1}^{T}e_{t,u}}{N/\beta}
    =
    \frac{\beta}{\mu}\bar{e}_{t,>t}
    =
    \kappa\bar{e}_{t,>t}.
    \qedhere
  \end{equation*}
\end{proof}

\begin{corollary}[Total suboptimality gap for the all-ones mask --- formal version]
  \label{cor:regret-allones-app}%
  Under the assumptions of Theorem~\ref{thm:epsilon-allones-app},
  the replanning suboptimality gap satisfies
  \begin{equation*}
    \mathrm{Gap}
    \le \frac{\kappa^{2}}{\mu}\sum_{t=1}^{T-1} \bar{e}_{t,>t}^{2}
    + \frac{1}{\mu}\sum_{t=1}^{T} e_{t,t}^{2}.
  \end{equation*}
\end{corollary}

\begin{proof}
  Theorem~\ref{thm:replan-regret} (strongly convex case) gives
  $\mathrm{Gap}\le (2\mu)^{-1}\sum_{t=1}^{T}(\epsilon_t+e_t)^2$.
  Theorem~\ref{thm:epsilon-allones-app} gives $\epsilon_t \le \kappa \bar{e}_{t,>t}$ for $t \le T-1$ and $\epsilon_T=0$.
  Using $(a+b)^2\le 2a^2+2b^2$ and dividing by $2\mu$ gives the claim.
\end{proof}

\section{Proofs of Section 5}
\label{app:proofs_5}

\begin{proof}[Proof of Proposition~\ref{prop:log2-spending}]
  The normalizing constant $Z$ is finite by an integral test.
  Anytime validity follows from the independence of conformal $p$-values:
  \begin{equation*}
    \prob(\exists t \ge 1 \colon p_t \le \tau_t)
    \le 1 - \prod_{t=1}^{\infty} (1 - \tau_t)
    = 1 - \exp\left(-\sum_{t=1}^{\infty} s_t\right)
    = 1 - e^{-L}
    = \alpha.
  \end{equation*}
  For the informativeness bound, recall $T_p := \inf\{t \ge 1 : \tau_t < 1/(n{+}1)\}$.
  Let
  \begin{equation*}
    t_0 \coloneqq \left\lfloor \frac{L}{8Z}\cdot \frac{n{+}1}{(\log(n{+}1))^2} \right\rfloor.
  \end{equation*}
  Assume $t_0 \ge 1$ (which holds for all $n$ large enough).
  Since $t\mapsto s_t$ is decreasing, for all $t \le t_0$ we have $s_t \ge s_{t_0}$.
  Let $A := \frac{L}{8Z}\cdot \frac{n{+}1}{(\log(n{+}1))^2} \ge 1$.
  Then $t_0+1 \le A+1 \le 2A = \frac{L}{4Z}\cdot \frac{n{+}1}{(\log(n{+}1))^2}$.
  For all $n$ large enough, $2A \le n{+}1$ (equivalently, $(\log(n{+}1))^2 \ge L/(4Z)$), hence $\log(t_0{+}1)\le \log(n{+}1)$.
  Therefore,
  \begin{equation*}
    s_{t_0}
    = \frac{L}{Z}\cdot \frac{1}{(t_0{+}1)(\log(t_0{+}1))^2}
    \ge \frac{4}{n{+}1}.
  \end{equation*}
  Therefore $s_t \ge 4/(n{+}1)$ for all $t \le t_0$, and
  \begin{equation*}
    \tau_t = 1-e^{-s_t} \ge 1-e^{-4/(n{+}1)}.
  \end{equation*}
  For $n \ge 3$ we have $4/(n{+}1)\le 1$, and $1-e^{-x}\ge x/2$ for $x\in[0,1]$ implies $\tau_t \ge 2/(n{+}1) > 1/(n{+}1)$.
  Thus $T_p \ge t_0 = \Omega(n/(\log n)^2)$.
\end{proof}\par\medskip

\begin{proof}[Proof of Lemma~\ref{lem:beta-law}]
  Let $S_{0},\dots,S_{n}\stackrel{\text{i.i.d.}}{\sim}\mathrm{Gamma}(a,\theta)$.
  The vector of normalized proportions
  \begin{equation*}
    \Big(\frac{S_{0}}{\sum_{i=0}^{n} S_{i}},\dots,\frac{S_{n}}{\sum_{i=0}^{n} S_{i}}\Big)
  \end{equation*}
  has a $\mathrm{Dirichlet}(a,\dots,a)$ distribution by the standard Gamma--Dirichlet construction (the common scale $\theta$ cancels).
  Hence the first coordinate $U \coloneqq S_{0}/\sum_{i=0}^{n} S_{i}$ is $\mathrm{Beta}(a,na)$.
  Since $E_{t}=(n+1)U$, we obtain $E_{t}/(n+1)\sim \mathrm{Beta}(a,na)$.

  For the moment formula, let $U\sim\mathrm{Beta}(a,na)$.
  For any $q>-a$,
  \begin{equation*}
    \expec[U^q]=\frac{B(a+q,na)}{B(a,na)}
    =\frac{\Gamma(a+q)\Gamma((n+1)a)}{\Gamma(a)\Gamma((n+1)a+q)}.
  \end{equation*}
  Multiplying by $(n+1)^q$ yields Equation~\eqref{eq:Eq-moment}.
\end{proof}\par\medskip

\begin{proof}[Proof of Theorem~\ref{thm:Tfull-bound}]
  If $T_e>t$, then in particular $W_{t-1}\ge (\alpha(n+1))^{-1}$.
  Therefore,
  \begin{equation*}
    \prob(T_e>t)
    \le \prob\Big(W_{t-1}\ge \frac{1}{\alpha(n+1)}\Big).
  \end{equation*}
  Fix any $q\in(0,1)$.
  By Markov's inequality,
  \begin{equation*}
    \prob\Big(W_{t-1}\ge \frac{1}{\alpha(n+1)}\Big)
    =\prob\Big(W_{t-1}^q\ge (\alpha(n+1))^{-q}\Big)
    \le (\alpha(n+1))^{q} \expec[W_{t-1}^q].
  \end{equation*}
  Since $W_{t-1}=\prod_{u=1}^{t-1} E_{u}$ and the $E_{u}$ are i.i.d., $\expec[W_{t-1}^q]=(\expec[E_{1}^q])^{t-1}$.
  By Lemma~\ref{lem:beta-law},
  \begin{equation*}
    c_{a,q,n}\coloneqq \expec[E_{1}^q]
    =(n+1)^q\frac{\Gamma(a+q)\Gamma((n+1)a)}{\Gamma(a)\Gamma((n+1)a+q)}.
  \end{equation*}
  Therefore,
  \begin{equation*}
    \prob(T_e>t)\le (\alpha(n+1))^{q} c_{a,q,n}^{t-1},
  \end{equation*}
  which is Equation~\eqref{eq:Tfull-tail}.

  It remains to justify that $c_{a,q,n}<1$ and to extract the $O(\log n)$ scaling.
  Let $U\coloneqq E_{1}/(n+1)\sim\mathrm{Beta}(a,na)$, so $\expec[U]=1/(n+1)$ and $U$ is non-degenerate.
  Since $q\in(0,1)$, strict concavity of $x\mapsto x^q$ implies $\expec[U^q]<(\expec[U])^q$, hence
  \begin{equation*}
    c_{a,q,n}
    =(n+1)^q\expec[U^q]
    < (n+1)^q\Big(\frac{1}{n+1}\Big)^q
    =1.
  \end{equation*}

  For the limit, let $x\coloneqq (n+1)a$.
  By Wendel's inequality, for all $q\in(0,1)$,
  \begin{equation*}
    1
    \le x^q\frac{\Gamma(x)}{\Gamma(x+q)}
    \le \Big(1+\frac{q}{x}\Big)^{1-q}.
  \end{equation*}
  Since $(1+q/x)^{1-q}\to 1$ as $x\to\infty$, we obtain
  \begin{equation*}
    c_{a,q,n}
    =\frac{\Gamma(a+q)}{\Gamma(a)a^q}\cdot \Big(x^q\frac{\Gamma(x)}{\Gamma(x+q)}\Big)
    \to c_{a,q}\coloneqq \frac{\Gamma(a+q)}{\Gamma(a)a^q}.
  \end{equation*}
  Since $X\sim\mathrm{Gamma}(a,1)$ is non-degenerate, strict concavity again gives $\expec[X^q]<(\expec[X])^q=a^q$, hence $c_{a,q}\in(0,1)$.

  Finally, for any $\delta\in(0,1)$, choosing $t$ such that $(\alpha(n+1))^{q} c_{a,q,n}^{t-1}\le \delta$ ensures $\prob(T_e>t)\le \delta$.
  Equivalently, it suffices that
  \begin{equation*}
    t
    \ge 1+\frac{\log((\alpha(n+1))^q/\delta)}{-\log(c_{a,q,n})},
  \end{equation*}
  which yields $T_e=O(\log n)$ with probability at least $1-\delta$ as $n\to\infty$ since $c_{a,q,n}\to c_{a,q}\in(0,1)$.
\end{proof}

% \section{Additional Results and Experimental Details} \label{app:extra_exeriments}
% \subsection{Dataset and Experimental Details}
% \label{app:datasets}

\section{Additional Appendix Material}
\label{app:additional_material}

\subsection{Experimental Details}
\label{app:extra_experiments}

This section provides a comprehensive description of the benchmarks and the sequential evaluation protocol used in our study. We evaluate our methods on two distinct tasks: a regression task and an image classification task, both of which are subjected to synthetic distribution drift to evaluate robustness under non-stationary conditions.

\paragraph{Regression: California Housing.}
We utilize the California Housing dataset~\citep{california_housing}, originally derived from the 1990 U.S. Census. The dataset consists of $20{,}640$ observations, where each entry represents a specific census tract. The model's objective is to predict the \emph{median house value} ($Y \in \mathbb{R}$, measured in units of \$100{,}000). For each tract, we consider eight predictive features:
\begin{itemize}
  \item \textbf{Demographics}: Median income (\texttt{MedInc}), house age (\texttt{HouseAge}), and average occupancy (\texttt{AveOccup}).
  \item \textbf{Housing Characteristics}: Average number of rooms (\texttt{AveRooms}) and bedrooms (\texttt{AveBedrms}) per household.
  \item \textbf{Geographic/Scale}: Total population (\texttt{Population}), as well as the \texttt{Latitude} and \texttt{Longitude} of the tract.
\end{itemize}

A Gradient Boosted Regressor (GBR) was trained on a random $70\%$ split ($14{,}448$ samples), leaving $30\%$ ($6{,}192$ samples) for sequential testing. We use the absolute residual as the nonconformity score: $S(x,y) = |y - \hat{f}(x)|$, where $\hat{f}(x)$ is the point estimate from the GBR. The spatial distribution and histograms of these features are visualized in Figure~\ref{fig:california-preview}.

\begin{figure}[ht]
  \centering
  \includegraphics[width=0.95\textwidth]{california_housing_preview_two-panel}
  \caption{California Housing (stationary). Left: geographic distribution (longitude vs latitude, colored by median house value). Right: histograms of the eight features and target.}
  \label{fig:california-preview}
\end{figure}

\paragraph{Classification: CIFAR-10.}
The CIFAR-10 benchmark~\citep{cifar10} comprises $60{,}000$ color images of size $32 \times 32$ distributed equally across 10 classes: \emph{airplane, automobile, bird, cat, deer, dog, frog, horse, ship,} and \emph{truck}. We split the data into the standard $50{,}000$ training and $10{,}000$ test images (see Figure~\ref{fig:cifar10-preview}).

As a base predictor, we employ a ResNet-18 architecture~\citep{resnet18} pre-trained on ImageNet~\citep{imagenet}. We fine-tuned the model on the CIFAR-10 training set for 10 epochs, replacing the final fully connected layer to output a 10-dimensional logit vector. To ensure numerical stability in sequential e-value martingales, we adopt the scoring construction from~\citet{gauthier2025values}:
\begin{equation}
  \label{eq:clf-score}
  S(x, k) = \frac{1}{\bigl(\log(1 + p_f(k \mid x))\bigr)^{1/4}},
\end{equation}
where $p_f(k \mid x)$ is the softmax probability assigned to class $k$. The $\log(1+p)$ transformation bounds the scores away from zero ($S > 1$ for all $p \in (0,1]$), preventing the product-martingale of the e-value baseline from collapsing prematurely across batches.

\begin{figure}[ht]
  \centering
  \includegraphics[width=0.87\textwidth]{cifar10_preview_multi-per-class}
  \caption{CIFAR-10 (stationary). Five sample images per class across the ten classes.}
  \label{fig:cifar10-preview}
\end{figure}

\paragraph{Batch Structure and Distribution Drift.}
Both tasks share a sequential protocol reflecting the framework of Section~\ref{sec:batch}. The $30\%$ held-out data (Regression) and the test set (Classification) are partitioned into $T = 20$ consecutive batches. At each timestep $t \in \{1, \dots, 20\}$, a batch arrives containing $m_t$ labeled calibration points and a single unlabeled test point $X_{t,n_t}$.

To evaluate robustness, we introduce a progressive drift magnitude $\rho_t = (t-1)/(T-1) \in [0,1]$.
\begin{itemize}
  \item \textbf{Regression Drift}: We simulate geographic/economic shifts by adding a constant offset of $0.5 \cdot \rho_t$ to all eight features. We also introduce \emph{target noise} by adding $\varepsilon \sim \mathcal{N}(0,\, (0.3 \cdot \rho_t)^2)$ to the house values, simulating increasing market volatility (see Figure~\ref{fig:california-drift}).
  \item \textbf{Classification Drift}: We simulate sensor degradation by increasing pixel brightness by $0.5 \cdot \rho_t$ and introduce \emph{label noise} by flipping a fraction $0.3 \cdot \rho_t$ of labels to a random different class (see Figure~\ref{fig:cifar10-drift}).
\end{itemize}

\begin{figure}[ht]
  \centering
  \includegraphics[width=0.78\textwidth]{california_housing_drift_preview_two-panel}
  \caption{California Housing under drift. Overlay of original (blue) and drifted (coral) distributions for each feature and the target at the final batch.}
  \label{fig:california-drift}
\end{figure}

\begin{figure}[ht]
  \centering
  \includegraphics[width=0.95\textwidth]{cifar10_drift_preview_multi-per-class}
  \caption{CIFAR-10 under drift. Original (left) vs drifted (right) image pairs per class; drift applies brightness increase, and random label flips.}
  \label{fig:cifar10-drift}
\end{figure}

% \subsection{Complete Results Tables}
% \label{app:complete_results}

\subsection{Per-Batch Visualisations}
\label{app:figures}

Figures~\ref{fig:reg-drift-traces} and~\ref{fig:clf-drift-traces} display per-run traces for a single representative run under drift on the regression and classification tasks, respectively.
Each figure is organised as a $6{\times}3$ grid: rows correspond to Mask Valid with $K\in\{1,3,5,7,9\}$ and Anytime CP; columns show (i) the prediction sets, (ii) the $p$-value or martingale trajectory, and (iii) the distribution of set sizes per batch.

\paragraph{Prediction sets (column 1).}
For regression (Figure~\ref{fig:reg-drift-traces}), each panel overlays the true label, model prediction, and conformal interval over 20 batches.
Green dots indicate covered test points; red crosses indicate miscoverages.
Intervals widen as drift accumulates, but the density of red crosses remains visually consistent across batches, reflecting that mask validity is a \emph{finite-sample} guarantee that does not require stationarity.
For classification (Figure~\ref{fig:clf-drift-traces}), the first column shows a heatmap of conformal set membership across labels and batches (black = label in set, white = excluded).
As $K$ increases, the prediction sets become larger to meet the stricter window guarantee.
In all rows, the ground-truth label (red cross) is included in the set in the large majority of batches, consistent with validity.

\paragraph{P-value or martingale trajectories (column 2).}
For Mask Valid methods, the second column shows the realised conformal $p$-value $p_t(X_t, Y_t)$ (blue dots) against the threshold $\tau_t$ (red dashed line).
A miscoverage occurs when a blue dot falls below the dashed line; the sparsity of such events confirms that W is close to but below $\alpha$.
The Anytime row replaces the $p$-value trace with the wealth process $W_t$.
In both figures, the martingale collapses from its initial value of approximately $1$ to near zero within the first few batches and remains effectively flat for the remaining steps, precisely as predicted by Theorem~\ref{thm:Tfull-bound}.
Once $W_t < 1/(\alpha(n{+}1))$, the prediction set must include the entire label space, and the e-value approach provides no further discrimination.

\paragraph{Set sizes per batch (column 3).}
Box plots summarise the distribution of prediction set sizes across calibration draws at each batch.
For windowed Mask Valid methods, set size grows smoothly with $K$ at every batch, forming a visible family of curves that interpolates between near point-prediction ($K{=}1$) and near-full-set ($K{=}9$ under drift).
For Anytime CP, the set sizes in column 3 tell the same collapse story as column 2: they start moderate in early batches but quickly rise to the maximum ($300$ for regression, $10$ for classification) and remain there, confirming that the wealth depletion shown in column 2 translates directly into uninformative prediction sets.

\begin{figure}[p]
  \centering
  \includegraphics[width=\textwidth]{gathered_reg_closed_drift}
  \caption{%
    \textbf{Per-batch traces, regression (California Housing, drift).}
    Rows (top to bottom): Mask Valid $K{=}1$, $K{=}3$, $K{=}5$, $K{=}7$, $K{=}9$, Anytime CP.
    \emph{Column 1 (Conformal Prediction Bands)}: conformal intervals over 20 batches for a single representative run; green dots = covered test labels, red crosses = miscoverages, dashed line = model prediction.
    Intervals widen progressively as drift accumulates.
    \emph{Column 2 (P-value / Martingale Trajectory)}: blue dots are $p_t(X_t, Y_t)$; red dashed line is the threshold $\tau_t$ (Mask Valid) or $1/\alpha$ (Anytime martingale $1/\alpha \approx 6.7$).
    For Mask Valid methods, errors occur when a $p$-value falls below the threshold.
    For Anytime, the martingale $W_t$ collapses to near zero within the first few batches and remains there, yielding no further discriminative power.
    \emph{Column 3 (Interval Sizes per Batch)}: box plots of prediction interval widths across calibration draws at each batch.
    Anytime interval widths jump to the dataset maximum shortly after the wealth collapses, while Mask Valid methods sustain informative set sizes throughout the horizon.
  }
  \label{fig:reg-drift-traces}
\end{figure}

\begin{figure}[p]
  \centering
  \includegraphics[width=\textwidth]{gathered_clf_closed_drift}
  \caption{%
    \textbf{Per-batch traces, classification (CIFAR-10, drift).}
    Same layout as Figure~\ref{fig:reg-drift-traces}.
    \emph{Column 1 (Conformal Set Sequence)}: heatmap of set membership across 10 labels (rows)
    and 20 batches (columns); white = label excluded, black = included, red cross = ground truth and green triangle = model prediction.
    Decreasing $K$ produces sparser heatmaps (fewer included labels per batch), reflecting the smaller, more informative sets.
    \emph{Column 2}: same as regression; for Anytime CP, $W_t$ again collapses within a few batches, after which the prediction set must include all 10 classes.
    \emph{Column 3}: set cardinality (number of classes included) per batch.
    For Anytime, cardinality jumps to $10$ (all classes) immediately after wealth collapse, while Mask Valid $K{=}1$ maintains a median set size of $3$--$4$ classes throughout the horizon.
  }
  \label{fig:clf-drift-traces}
\end{figure}

\subsection{Additional Experiments}
\label{app:additional_experiments}

Table~\ref{tab:results-full} extends the main Table~\ref{tab:results} by including the \emph{open-loop} (OL) variant of Mask Valid CP, in which the threshold schedule is computed once via the static log-budget program~\eqref{eq:design-problem-log} before any data are observed, using only the prior cost model.
Comparing open-loop (OL) and closed-loop (CL, replanning) variants isolates the benefit of adaptive budget allocation.
% leveraging the observed calibration quantiles at each step, recovers slightly more informative
% sets ($\mathrm{S}{=}287$ for regression ND), confirming that replanning provides a decisive
% efficiency advantage for the global mask at long horizons.

% At $T{=}15$, the shorter horizon gives the OL schedule a larger per-step budget allocation
% even without calibration data, so it already produces informative sets
% ($\mathrm{S}{=}175$ for regression ND, $\mathrm{S}{=}1.9$ for classification ND).
% Here CL all-ones does not improve over OL and can even yield larger sets
% (e.g., $\mathrm{S}{=}205$ vs $175$ for regression ND, $\mathrm{S}{=}9.1$ vs $1.9$ for classification ND),
% suggesting that with fewer steps the prior cost model is already sufficiently accurate and
% replanning introduces estimation variance rather than efficiency gains.

\paragraph{Windowed masks: OL and CL are nearly equivalent for regression; CL inflates classification sets.}
For $K \in \{1,3,5,7,9\}$ and both regression tasks, OL and CL produce essentially
identical set sizes (differences within one standard deviation in all cases).
The prior cost model is accurate enough for short windows that replanning offers no benefit.

For classification, a different pattern emerges at larger $K$: CL produces substantially larger sets with much higher variance than OL.
For instance at $K{=}3$ Clf.\ ND: OL gives $\mathrm{S}{=}1.0{\pm}0.2$ while CL gives $6.4{\pm}4.4$; at $K{=}5$ Clf.\ ND, OL gives $1.1{\pm}0.5$ vs CL $6.4{\pm}4.4$.
This indicates that the closed-loop cost estimator is sensitive to the highly variable softmax
confidence scores of the ResNet-18, leading to erratic budget allocations across runs.
Notably, these larger CL sets are paired with similar or lower W values, confirming that
CL is spending its budget more conservatively, not incorrectly---it is simply less efficient.
For windowed classification monitoring, the simpler OL schedule is therefore preferable.

\paragraph{Validity.}
Window failure rates are at or below $\alpha{=}0.15$ across all configurations, confirming empirical validity. Marginal exceedances (e.g., $W{=}0.17$ for CL $K{=}5$ Clf.\ ND) lie well within Monte Carlo noise for 100 runs.

\begin{table}[ht]
  \centering
  \caption{%
    Full results with $T = 20$ batches, 100 runs, $\alpha = 0.15$.
    \textbf{W}: window failure rate (target $\le 0.15$).
    \textbf{S}: mean prediction set size $\pm$ std.\ across runs.
    Open-loop (OL) methods fix thresholds before deployment; closed-loop (CL) methods replan using observed calibration data and an exponentially weighted cost forecaster.
  }
  \label{tab:results-full}
  \small
  \begin{tabular}{@{}l|cc|cc|cc|cc@{}}
    \toprule
    & \multicolumn{2}{c|}{Reg.\ ND}
    & \multicolumn{2}{c|}{Reg.\ D}
    & \multicolumn{2}{c|}{Clf.\ ND}
    & \multicolumn{2}{c}{Clf.\ D}                                                                                                                       \\
    \cmidrule(lr){2-3}\cmidrule(lr){4-5}\cmidrule(lr){6-7}\cmidrule(lr){8-9}
    Method                        & W                             & S            & W               & S            & W               & S             & W               & S             \\
    \midrule
    Anytime                       & --                            & $273{\pm}68$ & --              & $279{\pm}57$ & --              & $4.0{\pm}3.4$ & --              & $7.8{\pm}3.5$ \\
    \midrule
    Mask Valid OL ($K{=}1$)  & $0.15{\pm}0.08$               & $45{\pm}5$   & $0.13{\pm}0.08$ & $98{\pm}25$  & $0.15{\pm}0.08$ & $0.9{\pm}0.3$ & $0.15{\pm}0.08$ & $3.9{\pm}2.9$ \\
    Mask Valid OL ($K{=}3$)  & $0.16{\pm}0.14$               & $76{\pm}11$  & $0.14{\pm}0.13$ & $158{\pm}34$ & $0.18{\pm}0.14$ & $1.0{\pm}0.2$ & $0.14{\pm}0.12$ & $6.8{\pm}3.3$ \\
    Mask Valid OL ($K{=}5$)  & $0.13{\pm}0.18$               & $102{\pm}19$ & $0.13{\pm}0.17$ & $193{\pm}38$ & $0.16{\pm}0.18$ & $1.1{\pm}0.5$ & $0.17{\pm}0.20$ & $7.7{\pm}3.1$ \\
    Mask Valid OL ($K{=}7$)  & $0.12{\pm}0.21$               & $114{\pm}22$ & $0.13{\pm}0.21$ & $204{\pm}38$ & $0.15{\pm}0.22$ & $1.3{\pm}0.7$ & $0.17{\pm}0.25$ & $8.1{\pm}2.8$ \\
    Mask Valid OL ($K{=}9$)  & $0.10{\pm}0.21$               & $133{\pm}28$ & $0.13{\pm}0.25$ & $218{\pm}37$ & $0.15{\pm}0.26$ & $1.4{\pm}0.9$ & $0.20{\pm}0.30$ & $8.3{\pm}2.7$ \\
    \midrule
    Mask Valid CL ($K{=}1$)  & $0.14{\pm}0.08$               & $51{\pm}28$  & $0.13{\pm}0.08$ & $104{\pm}27$ & $0.15{\pm}0.08$ & $0.9{\pm}0.3$ & $0.15{\pm}0.08$ & $3.9{\pm}2.9$ \\
    Mask Valid CL ($K{=}3$)  & $0.14{\pm}0.14$               & $83{\pm}24$  & $0.13{\pm}0.13$ & $166{\pm}31$ & $0.18{\pm}0.14$ & $1.0{\pm}0.2$ & $0.14{\pm}0.13$ & $6.8{\pm}3.3$ \\
    Mask Valid CL ($K{=}5$)  & $0.15{\pm}0.18$               & $101{\pm}26$ & $0.13{\pm}0.17$ & $191{\pm}35$ & $0.16{\pm}0.18$ & $1.1{\pm}0.5$ & $0.17{\pm}0.21$ & $7.7{\pm}3.1$ \\
    Mask Valid CL ($K{=}7$)  & $0.12{\pm}0.20$               & $118{\pm}29$ & $0.14{\pm}0.21$ & $206{\pm}36$ & $0.15{\pm}0.22$ & $1.3{\pm}0.7$ & $0.17{\pm}0.25$ & $8.1{\pm}2.8$ \\
    Mask Valid CL ($K{=}9$)  & $0.10{\pm}0.21$               & $131{\pm}32$ & $0.14{\pm}0.26$ & $216{\pm}37$ & $0.15{\pm}0.26$ & $1.4{\pm}0.9$ & $0.20{\pm}0.30$ & $8.3{\pm}2.7$ \\
    \bottomrule
  \end{tabular}
\end{table}

% \subsection{\texorpdfstring{$T = 15$}{T=15} batches, 50 runs}
% \label{app:results-15-50}

% \begin{table}[ht]
%   \centering
%   \caption{Full results with $T = 15$ batches, 50 runs, $\alpha = 0.1$.
%     \textbf{W}: window failure rate (target $\le 0.1$).
%     \textbf{S}: mean prediction set size $\pm$ std.\ across runs.
%     Open-loop methods fix thresholds before deployment; closed-loop methods replan using observed
%     calibration data.}
%   \label{tab:results-15-50-full}
%   \small
%   \begin{tabular}{@{}l|cc|cc|cc|cc@{}}
%     \toprule
%                              & \multicolumn{2}{c|}{Reg.\ ND}
%                              & \multicolumn{2}{c|}{Reg.\ D}
%                              & \multicolumn{2}{c|}{Clf.\ ND}
%                              & \multicolumn{2}{c}{Clf.\ D}                                                                                                                         \\
%     \cmidrule(lr){2-3}\cmidrule(lr){4-5}\cmidrule(lr){6-7}\cmidrule(lr){8-9}
%     Method                   & W                             & S            & W               & S            & W               & S               & W               & S             \\
%     \midrule
%     Anytime                  & --                            & $273{\pm}62$ & --              & $287{\pm}41$ & --              & $4.3{\pm}3.5$   & --              & $9.4{\pm}1.9$ \\
%     Anyhow OL (all-ones)        & $0.10{\pm}0.30$               & $175{\pm}38$ & $0.14{\pm}0.35$ & $240{\pm}33$ & $0.04{\pm}0.20$ & $1.9{\pm}1.8$   & $0.00{\pm}0.00$ & $9.1{\pm}2.1$ \\
%     Anyhow OL  (win. K=1)        & $0.09{\pm}0.08$               & $56{\pm}6$   & $0.09{\pm}0.08$ & $122{\pm}29$ & $0.08{\pm}0.09$ & $0.95{\pm}0.22$ & $0.07{\pm}0.06$ & $4.0{\pm}2.8$ \\
%     Anyhow OL  (win. K=3)        & $0.11{\pm}0.14$               & $91{\pm}14$  & $0.12{\pm}0.17$ & $178{\pm}36$ & $0.12{\pm}0.13$ & $1.14{\pm}0.47$ & $0.06{\pm}0.10$ & $6.8{\pm}3.3$ \\
%     Anyhow OL  (win. K=5)        & $0.09{\pm}0.19$               & $113{\pm}19$ & $0.10{\pm}0.17$ & $201{\pm}37$ & $0.05{\pm}0.13$ & $1.30{\pm}0.73$ & $0.02{\pm}0.08$ & $7.8{\pm}3.0$ \\
%     Anyhow CL (all-ones) & $0.10{\pm}0.30$               & $205{\pm}64$ & $0.00{\pm}0.00$ & $248{\pm}38$ & $0.08{\pm}0.27$ & $9.1{\pm}2.6$   & $0.12{\pm}0.32$ & $9.4{\pm}1.9$ \\
%     Anyhow CL (win. K=1) & $0.10{\pm}0.09$               & $56{\pm}6$   & $0.09{\pm}0.09$ & $122{\pm}28$ & $0.08{\pm}0.07$ & $0.97{\pm}0.21$ & $0.04{\pm}0.04$ & $4.1{\pm}2.6$ \\
%     Anyhow CL (win. K=3) & $0.13{\pm}0.15$               & $91{\pm}15$  & $0.10{\pm}0.13$ & $180{\pm}37$ & $0.06{\pm}0.11$ & $4.3{\pm}4.2$   & $0.04{\pm}0.09$ & $7.5{\pm}3.2$ \\
%     Anyhow CL (win. K=5) & $0.12{\pm}0.19$               & $113{\pm}20$ & $0.13{\pm}0.19$ & $202{\pm}38$ & $0.10{\pm}0.17$ & $6.8{\pm}4.3$   & $0.02{\pm}0.07$ & $8.4{\pm}2.8$ \\
%     \bottomrule
%   \end{tabular}
% \end{table}

% \subsection{\texorpdfstring{$T = 10$}{T=10} batches, 20 runs}
% \label{app:results-10-20}

% \begin{table}[ht]
%   \centering
%   \caption{Full results with $T = 10$ batches, 10 runs, $\alpha = 0.1$.
%     \textbf{W}: window failure rate (target $\le 0.1$).
%     \textbf{S}: mean prediction set size $\pm$ std.\ across runs.
%     Open-loop methods fix thresholds before deployment; closed-loop methods replan using observed
%     calibration data.}
%   \label{tab:results-10-20-full}
%   \small
%   \begin{tabular}{@{}l|cc|cc|cc|cc@{}}
%     \toprule
%                              & \multicolumn{2}{c|}{Reg.\ ND}
%                              & \multicolumn{2}{c|}{Reg.\ D}
%                              & \multicolumn{2}{c|}{Clf.\ ND}
%                              & \multicolumn{2}{c}{Clf.\ D}                                                                                                                         \\
%     \cmidrule(lr){2-3}\cmidrule(lr){4-5}\cmidrule(lr){6-7}\cmidrule(lr){8-9}
%     Method                   & W                             & S            & W               & S            & W               & S               & W               & S             \\
%     \midrule
%     Anytime                   & -- & $274{\pm}61$ & -- & $283{\pm}41$ & -- & $3.8{\pm}3.3$ & -- & $9.3{\pm}2.0$ \\
%   Anyhow OL (all-ones)         & $0.00{\pm}0.00$ & $143{\pm}22$ & $0.1{\pm}0.30$ & $221.09{\pm}34.61$ & $0.20{\pm}0.40$ & $1.56{\pm}1.28$ & $0.00{\pm}0.00$ & $8.20{\pm}2.92$ \\
%   Anyhow OL (win. K=1)         & $0.10{\pm}0.10$ & $56.35{\pm}4.94$ & $0.10{\pm}0.09$ & $120.95{\pm}28.68$ & $0.09{\pm}0.11$ & $0.95{\pm}0.21$ & $0.05{\pm}0.06$ & $3.98{\pm}2.99$ \\
%   Anyhow OL (win. K=3)         & $0.05{\pm}0.12$ & $91.59{\pm}11.09$ & $0.07{\pm}0.17$ & $178.19{\pm}36.19$ & $0.08{\pm}0.14$ & $1.09{\pm}0.37$ & $0.03{\pm}0.08$ & $6.33{\pm}3.48$ \\
%   Anyhow OL (win. K=5)         & $0.06{\pm}0.19$ & $113.72{\pm}15.32$ & $0.11{\pm}0.24$ & $200.65{\pm}37.91$ & $0.09{\pm}0.19$ & $1.25{\pm}0.73$ & $0.00{\pm}0.00$ & $7.38{\pm}3.35$ \\
%   Anyhow CL (all-ones)  & $0.00{\pm}0.00$ & $152.86{\pm}46.32$ & $0.00{\pm}0.00$ & $222.66{\pm}35.06$ & $0.10{\pm}0.30$ & $8.52{\pm}3.33$ & $0.20{\pm}0.40$ & $9.01{\pm}2.51$ \\
%   Anyhow CL (win. K=1)  & $0.10{\pm}0.10$ & $55.75{\pm}4.89$ & $0.10{\pm}0.09$ & $119.47{\pm}28.48$ & $0.08{\pm}0.10$ & $0.97{\pm}0.19$ & $0.04{\pm}0.05$ & $4.04{\pm}2.94$ \\
%   Anyhow CL (win. K=3)  & $0.08{\pm}0.14$ & $91.27{\pm}13.88$ & $0.07{\pm}0.17$ & $178.25{\pm}40.11$ & $0.07{\pm}0.12$ & $5.07{\pm}4.43$ & $0.04{\pm}0.10$ & $7.11{\pm}3.46$ \\
%   Anyhow CL (win. K=5)  & $0.06{\pm}0.19$ & $112.69{\pm}15.58$ & $0.11{\pm}0.24$ & $199.83{\pm}38.37$ & $0.03{\pm}0.11$ & $6.53{\pm}4.35$ & $0.05{\pm}0.11$ & $8.27{\pm}3.01$ \\
%   \bottomrule
%   \end{tabular}
% \end{table}

% \subsection{$T = 35$ batches, 20 runs, old replanning}
% \label{app:results-35-20-old}

% \begin{table}[ht]
%   \centering
%   \caption{Results with $T = 35$ batches, 20 runs, old replanning implementation.}
%   \label{tab:results-35-20-old}
%   \small
%   \begin{tabular}{@{}l|cc|cc|cc|cc@{}}
%     \toprule
%                              & \multicolumn{2}{c|}{Reg.\ ND}
%                              & \multicolumn{2}{c|}{Reg.\ D}
%                              & \multicolumn{2}{c|}{Clf.\ ND}
%                              & \multicolumn{2}{c}{Clf.\ D}                                                                                                                          \\
%     \cmidrule(lr){2-3}\cmidrule(lr){4-5}\cmidrule(lr){6-7}\cmidrule(lr){8-9}
%     Method                   & W                             & S            & W               & S            & W               & S               & W               & S              \\
%     \midrule
%     Anytime                  & --                            & $291{\pm}37$ & --              & $292{\pm}35$ & --              & $6.2{\pm}4.0$   & --              & $9.8{\pm}1.3$  \\
%     Anyhow (all-ones)        & $0.00{\pm}0.00$               & $300{\pm}0$  & $0.00{\pm}0.00$ & $300{\pm}0$  & $0.00{\pm}0.00$ & $10.0{\pm}0.0$  & $0.00{\pm}0.00$ & $10.0{\pm}0.0$ \\
%     Anyhow (win. K=1)        & $0.10{\pm}0.04$               & $59{\pm}9$   & $0.09{\pm}0.04$ & $127{\pm}32$ & $0.09{\pm}0.05$ & $0.96{\pm}0.23$ & $0.07{\pm}0.04$ & $4.0{\pm}2.7$  \\
%     Anyhow (win. K=3)        & $0.14{\pm}0.10$               & $95{\pm}21$  & $0.10{\pm}0.08$ & $183{\pm}40$ & $0.10{\pm}0.09$ & $1.23{\pm}0.68$ & $0.03{\pm}0.06$ & $7.2{\pm}3.2$  \\
%     Anyhow (win. K=5)        & $0.11{\pm}0.12$               & $147{\pm}40$ & $0.08{\pm}0.08$ & $224{\pm}38$ & $0.10{\pm}0.10$ & $1.59{\pm}1.43$ & $0.00{\pm}0.02$ & $8.5{\pm}2.7$  \\
%     Anyhow replan (all-ones) & $0.00{\pm}0.00$               & $300{\pm}7$  & $0.00{\pm}0.00$ & $300{\pm}7$  & $0.00{\pm}0.00$ & $10.0{\pm}0.0$  & $0.00{\pm}0.00$ & $10.0{\pm}0.3$ \\
%     Anyhow replan (win. K=1) & $0.10{\pm}0.04$               & $59{\pm}9$   & $0.09{\pm}0.04$ & $127{\pm}32$ & $0.06{\pm}0.04$ & $1.02{\pm}0.24$ & $0.06{\pm}0.04$ & $4.3{\pm}2.7$  \\
%     Anyhow replan (win. K=3) & $0.12{\pm}0.10$               & $104{\pm}32$ & $0.10{\pm}0.09$ & $193{\pm}44$ & $0.10{\pm}0.08$ & $1.32{\pm}0.78$ & $0.04{\pm}0.06$ & $7.3{\pm}3.1$  \\
%     Anyhow replan (win. K=5) & $0.12{\pm}0.10$               & $160{\pm}65$ & $0.08{\pm}0.08$ & $224{\pm}43$ & $0.13{\pm}0.11$ & $1.46{\pm}1.08$ & $0.02{\pm}0.05$ & $8.4{\pm}2.7$  \\
%     \bottomrule
%   \end{tabular}
% \end{table}

% \subsection{$T = 15$ batches, 50 runs, old replanning}
% \label{app:results-15-50-old}

% \begin{table}[ht]
%   \centering
%   \caption{Results with $T = 15$ batches, 50 runs, old replanning implementation.}
%   \label{tab:results-15-50-old}
%   \small
%   \begin{tabular}{@{}l|cc|cc|cc|cc@{}}
%     \toprule
%                              & \multicolumn{2}{c|}{Reg.\ ND}
%                              & \multicolumn{2}{c|}{Reg.\ D}
%                              & \multicolumn{2}{c|}{Clf.\ ND}
%                              & \multicolumn{2}{c}{Clf.\ D}                                                                                                                         \\
%     \cmidrule(lr){2-3}\cmidrule(lr){4-5}\cmidrule(lr){6-7}\cmidrule(lr){8-9}
%     Method                   & W                             & S            & W               & S            & W               & S               & W               & S             \\
%     \midrule
%     Anytime                  & --                            & $276{\pm}59$ & --              & $286{\pm}44$ & --              & $4.4{\pm}3.6$   & --              & $9.6{\pm}1.6$ \\
%     Anyhow (all-ones)        & $0.04{\pm}0.20$               & $177{\pm}38$ & $0.10{\pm}0.30$ & $240{\pm}33$ & $0.10{\pm}0.30$ & $2.0{\pm}2.0$   & $0.00{\pm}0.00$ & $9.3{\pm}1.9$ \\
%     Anyhow (win. K=1)        & $0.10{\pm}0.06$               & $56{\pm}6$   & $0.10{\pm}0.09$ & $121{\pm}28$ & $0.12{\pm}0.09$ & $0.93{\pm}0.25$ & $0.05{\pm}0.06$ & $4.1{\pm}2.7$ \\
%     Anyhow (win. K=3)        & $0.12{\pm}0.14$               & $90{\pm}14$  & $0.10{\pm}0.14$ & $179{\pm}36$ & $0.14{\pm}0.18$ & $1.14{\pm}0.43$ & $0.03{\pm}0.09$ & $7.0{\pm}3.2$ \\
%     Anyhow (win. K=5)        & $0.11{\pm}0.20$               & $113{\pm}19$ & $0.12{\pm}0.22$ & $203{\pm}37$ & $0.12{\pm}0.19$ & $1.34{\pm}0.86$ & $0.01{\pm}0.07$ & $8.1{\pm}2.9$ \\
%     Anyhow replan (all-ones) & $0.02{\pm}0.14$               & $179{\pm}41$ & $0.10{\pm}0.30$ & $240{\pm}33$ & $0.10{\pm}0.30$ & $2.0{\pm}2.0$   & $0.00{\pm}0.00$ & $9.2{\pm}2.0$ \\
%     Anyhow replan (win. K=1) & $0.10{\pm}0.06$               & $56{\pm}6$   & $0.10{\pm}0.09$ & $121{\pm}28$ & $0.08{\pm}0.07$ & $1.02{\pm}0.23$ & $0.05{\pm}0.06$ & $4.5{\pm}2.6$ \\
%     Anyhow replan (win. K=3) & $0.11{\pm}0.14$               & $92{\pm}16$  & $0.10{\pm}0.14$ & $181{\pm}38$ & $0.12{\pm}0.15$ & $1.23{\pm}0.59$ & $0.02{\pm}0.07$ & $7.1{\pm}3.1$ \\
%     Anyhow replan (win. K=5) & $0.10{\pm}0.19$               & $114{\pm}22$ & $0.12{\pm}0.22$ & $204{\pm}39$ & $0.11{\pm}0.19$ & $1.37{\pm}0.84$ & $0.02{\pm}0.09$ & $8.0{\pm}2.8$ \\
%     \bottomrule
%   \end{tabular}
% \end{table}

% \subsection{$T = 15$ batches, 50 runs, new replanning}
% \label{app:results-15-50-new}

% \begin{table}[ht]
%   \centering
%   \caption{Results with $T = 15$ batches, 50 runs, new replanning implementation.}
%   \label{tab:results-15-50-new}
%   \small
%   \begin{tabular}{@{}l|cc|cc|cc|cc@{}}
%     \toprule
%                              & \multicolumn{2}{c|}{Reg.\ ND}
%                              & \multicolumn{2}{c|}{Reg.\ D}
%                              & \multicolumn{2}{c|}{Clf.\ ND}
%                              & \multicolumn{2}{c}{Clf.\ D}                                                                                                                         \\
%     \cmidrule(lr){2-3}\cmidrule(lr){4-5}\cmidrule(lr){6-7}\cmidrule(lr){8-9}
%     Method                   & W                             & S            & W               & S            & W               & S               & W               & S             \\
%     \midrule
%     Anytime                  & --                            & $273{\pm}62$ & --              & $287{\pm}41$ & --              & $4.3{\pm}3.5$   & --              & $9.4{\pm}1.9$ \\
%     Anyhow (all-ones)        & $0.10{\pm}0.30$               & $175{\pm}38$ & $0.14{\pm}0.35$ & $240{\pm}33$ & $0.04{\pm}0.20$ & $1.9{\pm}1.8$   & $0.00{\pm}0.00$ & $9.1{\pm}2.1$ \\
%     Anyhow (win. K=1)        & $0.09{\pm}0.08$               & $56{\pm}6$   & $0.09{\pm}0.08$ & $122{\pm}29$ & $0.08{\pm}0.09$ & $0.95{\pm}0.22$ & $0.07{\pm}0.06$ & $4.0{\pm}2.8$ \\
%     Anyhow (win. K=3)        & $0.11{\pm}0.14$               & $91{\pm}14$  & $0.12{\pm}0.17$ & $178{\pm}36$ & $0.12{\pm}0.13$ & $1.14{\pm}0.47$ & $0.06{\pm}0.10$ & $6.8{\pm}3.3$ \\
%     Anyhow (win. K=5)        & $0.09{\pm}0.19$               & $113{\pm}19$ & $0.10{\pm}0.17$ & $201{\pm}37$ & $0.05{\pm}0.13$ & $1.30{\pm}0.73$ & $0.02{\pm}0.08$ & $7.8{\pm}3.0$ \\
%     Anyhow replan (all-ones) & $0.10{\pm}0.30$               & $205{\pm}64$ & $0.00{\pm}0.00$ & $248{\pm}38$ & $0.08{\pm}0.27$ & $9.1{\pm}2.6$   & $0.12{\pm}0.32$ & $9.4{\pm}1.9$ \\
%     Anyhow replan (win. K=1) & $0.10{\pm}0.09$               & $56{\pm}6$   & $0.09{\pm}0.09$ & $122{\pm}28$ & $0.08{\pm}0.07$ & $0.97{\pm}0.21$ & $0.04{\pm}0.04$ & $4.1{\pm}2.6$ \\
%     Anyhow replan (win. K=3) & $0.13{\pm}0.15$               & $91{\pm}15$  & $0.10{\pm}0.13$ & $180{\pm}37$ & $0.06{\pm}0.11$ & $4.3{\pm}4.2$   & $0.04{\pm}0.09$ & $7.5{\pm}3.2$ \\
%     Anyhow replan (win. K=5) & $0.12{\pm}0.19$               & $113{\pm}20$ & $0.13{\pm}0.19$ & $202{\pm}38$ & $0.10{\pm}0.17$ & $6.8{\pm}4.3$   & $0.02{\pm}0.07$ & $8.4{\pm}2.8$ \\
%     \bottomrule
%   \end{tabular}
% \end{table}

% \section{Additional Experiments}
% \label{app:add_experiments}

% \subsection{Regression}
\subsubsection{Without drift --- All-ones.}
\begin{figure}[H]
  \centering
  \begin{subfigure}[b]{0.48\textwidth}
    \includegraphics[width=\textwidth]{figures/appendix/regression/nodrift/anyhow_all_ones/coverage_distribution}
    \caption{Coverage distribution}
  \end{subfigure}
  \hfill
  \begin{subfigure}[b]{0.48\textwidth}
    \includegraphics[width=\textwidth]{figures/appendix/regression/nodrift/anyhow_all_ones/pvalue_trajectory}
    \caption{P-value trajectory}
  \end{subfigure}
  \medskip
  \begin{subfigure}[b]{0.48\textwidth}
    \includegraphics[width=\textwidth]{figures/appendix/regression/nodrift/anyhow_all_ones/regression_bands}
    \caption{Regression bands}
  \end{subfigure}
  \hfill
  \begin{subfigure}[b]{0.48\textwidth}
    \includegraphics[width=\textwidth]{figures/appendix/regression/nodrift/anyhow_all_ones/set_size_histogram}
    \caption{Set size histogram}
  \end{subfigure}
  \medskip
  \centering
  \begin{subfigure}[b]{0.48\textwidth}
    \includegraphics[width=\textwidth]{figures/appendix/regression/nodrift/anyhow_all_ones/set_sizes_per_batch}
    \caption{Set sizes per batch}
  \end{subfigure}
  \caption{Regression, no drift, all-ones.}
  \label{fig:app-reg-nd-all-ones}
\end{figure}

\subsubsection{Without drift --- Window \texorpdfstring{$K{=}1$}{K=1}.}
\begin{figure}[H]
  \centering
  \begin{subfigure}[b]{0.48\textwidth}
    \includegraphics[width=\textwidth]{figures/appendix/regression/nodrift/anyhow_window_1/coverage_distribution}
    \caption{Coverage distribution}
  \end{subfigure}
  \hfill
  \begin{subfigure}[b]{0.48\textwidth}
    \includegraphics[width=\textwidth]{figures/appendix/regression/nodrift/anyhow_window_1/pvalue_trajectory}
    \caption{P-value trajectory}
  \end{subfigure}
  \medskip
  \begin{subfigure}[b]{0.48\textwidth}
    \includegraphics[width=\textwidth]{figures/appendix/regression/nodrift/anyhow_window_1/regression_bands}
    \caption{Regression bands}
  \end{subfigure}
  \hfill
  \begin{subfigure}[b]{0.48\textwidth}
    \includegraphics[width=\textwidth]{figures/appendix/regression/nodrift/anyhow_window_1/set_size_histogram}
    \caption{Set size histogram}
  \end{subfigure}
  \medskip
  \centering
  \begin{subfigure}[b]{0.48\textwidth}
    \includegraphics[width=\textwidth]{figures/appendix/regression/nodrift/anyhow_window_1/set_sizes_per_batch}
    \caption{Set sizes per batch}
  \end{subfigure}
  \caption{Regression, no drift, window $K{=}1$.}
  \label{fig:app-reg-nd-window-1}
\end{figure}

\subsubsection{Without drift --- Window \texorpdfstring{$K{=}3$}{K=3}.}
\begin{figure}[H]
  \centering
  \begin{subfigure}[b]{0.48\textwidth}
    \includegraphics[width=\textwidth]{figures/appendix/regression/nodrift/anyhow_window_3/coverage_distribution}
    \caption{Coverage distribution}
  \end{subfigure}
  \hfill
  \begin{subfigure}[b]{0.48\textwidth}
    \includegraphics[width=\textwidth]{figures/appendix/regression/nodrift/anyhow_window_3/pvalue_trajectory}
    \caption{P-value trajectory}
  \end{subfigure}
  \medskip
  \begin{subfigure}[b]{0.48\textwidth}
    \includegraphics[width=\textwidth]{figures/appendix/regression/nodrift/anyhow_window_3/regression_bands}
    \caption{Regression bands}
  \end{subfigure}
  \hfill
  \begin{subfigure}[b]{0.48\textwidth}
    \includegraphics[width=\textwidth]{figures/appendix/regression/nodrift/anyhow_window_3/set_size_histogram}
    \caption{Set size histogram}
  \end{subfigure}
  \medskip
  \centering
  \begin{subfigure}[b]{0.48\textwidth}
    \includegraphics[width=\textwidth]{figures/appendix/regression/nodrift/anyhow_window_3/set_sizes_per_batch}
    \caption{Set sizes per batch}
  \end{subfigure}
  \caption{Regression, no drift, window $K{=}3$.}
  \label{fig:app-reg-nd-window-3}
\end{figure}

\subsubsection{Without drift --- Window \texorpdfstring{$K{=}5$}{K=5}.}
\begin{figure}[H]
  \centering
  \begin{subfigure}[b]{0.48\textwidth}
    \includegraphics[width=\textwidth]{figures/appendix/regression/nodrift/anyhow_window_5/coverage_distribution}
    \caption{Coverage distribution}
  \end{subfigure}
  \hfill
  \begin{subfigure}[b]{0.48\textwidth}
    \includegraphics[width=\textwidth]{figures/appendix/regression/nodrift/anyhow_window_5/pvalue_trajectory}
    \caption{P-value trajectory}
  \end{subfigure}
  \medskip
  \begin{subfigure}[b]{0.48\textwidth}
    \includegraphics[width=\textwidth]{figures/appendix/regression/nodrift/anyhow_window_5/regression_bands}
    \caption{Regression bands}
  \end{subfigure}
  \hfill
  \begin{subfigure}[b]{0.48\textwidth}
    \includegraphics[width=\textwidth]{figures/appendix/regression/nodrift/anyhow_window_5/set_size_histogram}
    \caption{Set size histogram}
  \end{subfigure}
  \medskip
  \centering
  \begin{subfigure}[b]{0.48\textwidth}
    \includegraphics[width=\textwidth]{figures/appendix/regression/nodrift/anyhow_window_5/set_sizes_per_batch}
    \caption{Set sizes per batch}
  \end{subfigure}
  \caption{Regression, no drift, window $K{=}5$.}
  \label{fig:app-reg-nd-window-5}
\end{figure}

\subsubsection{Without drift --- Window \texorpdfstring{$K{=}10$}{K=10}.}
\begin{figure}[H]
  \centering
  \begin{subfigure}[b]{0.48\textwidth}
    \includegraphics[width=\textwidth]{figures/appendix/regression/nodrift/anyhow_window_10/coverage_distribution}
    \caption{Coverage distribution}
  \end{subfigure}
  \hfill
  \begin{subfigure}[b]{0.48\textwidth}
    \includegraphics[width=\textwidth]{figures/appendix/regression/nodrift/anyhow_window_10/pvalue_trajectory}
    \caption{P-value trajectory}
  \end{subfigure}
  \medskip
  \begin{subfigure}[b]{0.48\textwidth}
    \includegraphics[width=\textwidth]{figures/appendix/regression/nodrift/anyhow_window_10/regression_bands}
    \caption{Regression bands}
  \end{subfigure}
  \hfill
  \begin{subfigure}[b]{0.48\textwidth}
    \includegraphics[width=\textwidth]{figures/appendix/regression/nodrift/anyhow_window_10/set_size_histogram}
    \caption{Set size histogram}
  \end{subfigure}
  \medskip
  \centering
  \begin{subfigure}[b]{0.48\textwidth}
    \includegraphics[width=\textwidth]{figures/appendix/regression/nodrift/anyhow_window_10/set_sizes_per_batch}
    \caption{Set sizes per batch}
  \end{subfigure}
  \caption{Regression, no drift, window $K{=}10$.}
  \label{fig:app-reg-nd-window-10}
\end{figure}

\subsubsection{Without drift --- Window \texorpdfstring{$K{=}15$}{K=15}.}
\begin{figure}[H]
  \centering
  \begin{subfigure}[b]{0.48\textwidth}
    \includegraphics[width=\textwidth]{figures/appendix/regression/nodrift/anyhow_window_15/coverage_distribution}
    \caption{Coverage distribution}
  \end{subfigure}
  \hfill
  \begin{subfigure}[b]{0.48\textwidth}
    \includegraphics[width=\textwidth]{figures/appendix/regression/nodrift/anyhow_window_15/pvalue_trajectory}
    \caption{P-value trajectory}
  \end{subfigure}
  \medskip
  \begin{subfigure}[b]{0.48\textwidth}
    \includegraphics[width=\textwidth]{figures/appendix/regression/nodrift/anyhow_window_15/regression_bands}
    \caption{Regression bands}
  \end{subfigure}
  \hfill
  \begin{subfigure}[b]{0.48\textwidth}
    \includegraphics[width=\textwidth]{figures/appendix/regression/nodrift/anyhow_window_15/set_size_histogram}
    \caption{Set size histogram}
  \end{subfigure}
  \medskip
  \centering
  \begin{subfigure}[b]{0.48\textwidth}
    \includegraphics[width=\textwidth]{figures/appendix/regression/nodrift/anyhow_window_15/set_sizes_per_batch}
    \caption{Set sizes per batch}
  \end{subfigure}
  \caption{Regression, no drift, window $K{=}15$.}
  \label{fig:app-reg-nd-window-15}
\end{figure}

\subsubsection{Without drift --- Window \texorpdfstring{$K{=}25$}{K=25}.}
\begin{figure}[H]
  \centering
  \begin{subfigure}[b]{0.48\textwidth}
    \includegraphics[width=\textwidth]{figures/appendix/regression/nodrift/anyhow_window_25/coverage_distribution}
    \caption{Coverage distribution}
  \end{subfigure}
  \hfill
  \begin{subfigure}[b]{0.48\textwidth}
    \includegraphics[width=\textwidth]{figures/appendix/regression/nodrift/anyhow_window_25/pvalue_trajectory}
    \caption{P-value trajectory}
  \end{subfigure}
  \medskip
  \begin{subfigure}[b]{0.48\textwidth}
    \includegraphics[width=\textwidth]{figures/appendix/regression/nodrift/anyhow_window_25/regression_bands}
    \caption{Regression bands}
  \end{subfigure}
  \hfill
  \begin{subfigure}[b]{0.48\textwidth}
    \includegraphics[width=\textwidth]{figures/appendix/regression/nodrift/anyhow_window_25/set_size_histogram}
    \caption{Set size histogram}
  \end{subfigure}
  \medskip
  \centering
  \begin{subfigure}[b]{0.48\textwidth}
    \includegraphics[width=\textwidth]{figures/appendix/regression/nodrift/anyhow_window_25/set_sizes_per_batch}
    \caption{Set sizes per batch}
  \end{subfigure}
  \caption{Regression, no drift, window $K{=}25$.}
  \label{fig:app-reg-nd-window-25}
\end{figure}

\subsubsection{With drift --- All-ones.}
\begin{figure}[H]
  \centering
  \begin{subfigure}[b]{0.48\textwidth}
    \includegraphics[width=\textwidth]{figures/appendix/regression/drift/anyhow_all_ones/coverage_distribution}
    \caption{Coverage distribution}
  \end{subfigure}
  \hfill
  \begin{subfigure}[b]{0.48\textwidth}
    \includegraphics[width=\textwidth]{figures/appendix/regression/drift/anyhow_all_ones/pvalue_trajectory}
    \caption{P-value trajectory}
  \end{subfigure}
  \medskip
  \begin{subfigure}[b]{0.48\textwidth}
    \includegraphics[width=\textwidth]{figures/appendix/regression/drift/anyhow_all_ones/regression_bands}
    \caption{Regression bands}
  \end{subfigure}
  \hfill
  \begin{subfigure}[b]{0.48\textwidth}
    \includegraphics[width=\textwidth]{figures/appendix/regression/drift/anyhow_all_ones/set_size_histogram}
    \caption{Set size histogram}
  \end{subfigure}
  \medskip
  \centering
  \begin{subfigure}[b]{0.48\textwidth}
    \includegraphics[width=\textwidth]{figures/appendix/regression/drift/anyhow_all_ones/set_sizes_per_batch}
    \caption{Set sizes per batch}
  \end{subfigure}
  \caption{Regression, drift, all-ones.}
  \label{fig:app-reg-d-all-ones}
\end{figure}

\subsubsection{With drift --- Window \texorpdfstring{$K{=}1$}{K=1}.}
\begin{figure}[H]
  \centering
  \begin{subfigure}[b]{0.48\textwidth}
    \includegraphics[width=\textwidth]{figures/appendix/regression/drift/anyhow_window_1/coverage_distribution}
    \caption{Coverage distribution}
  \end{subfigure}
  \hfill
  \begin{subfigure}[b]{0.48\textwidth}
    \includegraphics[width=\textwidth]{figures/appendix/regression/drift/anyhow_window_1/pvalue_trajectory}
    \caption{P-value trajectory}
  \end{subfigure}
  \medskip
  \begin{subfigure}[b]{0.48\textwidth}
    \includegraphics[width=\textwidth]{figures/appendix/regression/drift/anyhow_window_1/regression_bands}
    \caption{Regression bands}
  \end{subfigure}
  \hfill
  \begin{subfigure}[b]{0.48\textwidth}
    \includegraphics[width=\textwidth]{figures/appendix/regression/drift/anyhow_window_1/set_size_histogram}
    \caption{Set size histogram}
  \end{subfigure}
  \medskip
  \centering
  \begin{subfigure}[b]{0.48\textwidth}
    \includegraphics[width=\textwidth]{figures/appendix/regression/drift/anyhow_window_1/set_sizes_per_batch}
    \caption{Set sizes per batch}
  \end{subfigure}
  \caption{Regression, drift, window $K{=}1$.}
  \label{fig:app-reg-d-window-1}
\end{figure}

\subsubsection{With drift --- Window \texorpdfstring{$K{=}3$}{K=3}.}
\begin{figure}[H]
  \centering
  \begin{subfigure}[b]{0.48\textwidth}
    \includegraphics[width=\textwidth]{figures/appendix/regression/drift/anyhow_window_3/coverage_distribution}
    \caption{Coverage distribution}
  \end{subfigure}
  \hfill
  \begin{subfigure}[b]{0.48\textwidth}
    \includegraphics[width=\textwidth]{figures/appendix/regression/drift/anyhow_window_3/pvalue_trajectory}
    \caption{P-value trajectory}
  \end{subfigure}
  \medskip
  \begin{subfigure}[b]{0.48\textwidth}
    \includegraphics[width=\textwidth]{figures/appendix/regression/drift/anyhow_window_3/regression_bands}
    \caption{Regression bands}
  \end{subfigure}
  \hfill
  \begin{subfigure}[b]{0.48\textwidth}
    \includegraphics[width=\textwidth]{figures/appendix/regression/drift/anyhow_window_3/set_size_histogram}
    \caption{Set size histogram}
  \end{subfigure}
  \medskip
  \centering
  \begin{subfigure}[b]{0.48\textwidth}
    \includegraphics[width=\textwidth]{figures/appendix/regression/drift/anyhow_window_3/set_sizes_per_batch}
    \caption{Set sizes per batch}
  \end{subfigure}
  \caption{Regression, drift, window $K{=}3$.}
  \label{fig:app-reg-d-window-3}
\end{figure}

\subsubsection{With drift --- Window \texorpdfstring{$K{=}5$}{K=5}.}
\begin{figure}[H]
  \centering
  \begin{subfigure}[b]{0.48\textwidth}
    \includegraphics[width=\textwidth]{figures/appendix/regression/drift/anyhow_window_5/coverage_distribution}
    \caption{Coverage distribution}
  \end{subfigure}
  \hfill
  \begin{subfigure}[b]{0.48\textwidth}
    \includegraphics[width=\textwidth]{figures/appendix/regression/drift/anyhow_window_5/pvalue_trajectory}
    \caption{P-value trajectory}
  \end{subfigure}
  \medskip
  \begin{subfigure}[b]{0.48\textwidth}
    \includegraphics[width=\textwidth]{figures/appendix/regression/drift/anyhow_window_5/regression_bands}
    \caption{Regression bands}
  \end{subfigure}
  \hfill
  \begin{subfigure}[b]{0.48\textwidth}
    \includegraphics[width=\textwidth]{figures/appendix/regression/drift/anyhow_window_5/set_size_histogram}
    \caption{Set size histogram}
  \end{subfigure}
  \medskip
  \centering
  \begin{subfigure}[b]{0.48\textwidth}
    \includegraphics[width=\textwidth]{figures/appendix/regression/drift/anyhow_window_5/set_sizes_per_batch}
    \caption{Set sizes per batch}
  \end{subfigure}
  \caption{Regression, drift, window $K{=}5$.}
  \label{fig:app-reg-d-window-5}
\end{figure}

\subsubsection{With drift --- Window \texorpdfstring{$K{=}10$}{K=10}.}
\begin{figure}[H]
  \centering
  \begin{subfigure}[b]{0.48\textwidth}
    \includegraphics[width=\textwidth]{figures/appendix/regression/drift/anyhow_window_10/coverage_distribution}
    \caption{Coverage distribution}
  \end{subfigure}
  \hfill
  \begin{subfigure}[b]{0.48\textwidth}
    \includegraphics[width=\textwidth]{figures/appendix/regression/drift/anyhow_window_10/pvalue_trajectory}
    \caption{P-value trajectory}
  \end{subfigure}
  \medskip
  \begin{subfigure}[b]{0.48\textwidth}
    \includegraphics[width=\textwidth]{figures/appendix/regression/drift/anyhow_window_10/regression_bands}
    \caption{Regression bands}
  \end{subfigure}
  \hfill
  \begin{subfigure}[b]{0.48\textwidth}
    \includegraphics[width=\textwidth]{figures/appendix/regression/drift/anyhow_window_10/set_size_histogram}
    \caption{Set size histogram}
  \end{subfigure}
  \medskip
  \centering
  \begin{subfigure}[b]{0.48\textwidth}
    \includegraphics[width=\textwidth]{figures/appendix/regression/drift/anyhow_window_10/set_sizes_per_batch}
    \caption{Set sizes per batch}
  \end{subfigure}
  \caption{Regression, drift, window $K{=}10$.}
  \label{fig:app-reg-d-window-10}
\end{figure}

\subsubsection{With drift --- Window \texorpdfstring{$K{=}15$}{K=15}.}
\begin{figure}[H]
  \centering
  \begin{subfigure}[b]{0.48\textwidth}
    \includegraphics[width=\textwidth]{figures/appendix/regression/drift/anyhow_window_15/coverage_distribution}
    \caption{Coverage distribution}
  \end{subfigure}
  \hfill
  \begin{subfigure}[b]{0.48\textwidth}
    \includegraphics[width=\textwidth]{figures/appendix/regression/drift/anyhow_window_15/pvalue_trajectory}
    \caption{P-value trajectory}
  \end{subfigure}
  \medskip
  \begin{subfigure}[b]{0.48\textwidth}
    \includegraphics[width=\textwidth]{figures/appendix/regression/drift/anyhow_window_15/regression_bands}
    \caption{Regression bands}
  \end{subfigure}
  \hfill
  \begin{subfigure}[b]{0.48\textwidth}
    \includegraphics[width=\textwidth]{figures/appendix/regression/drift/anyhow_window_15/set_size_histogram}
    \caption{Set size histogram}
  \end{subfigure}
  \medskip
  \centering
  \begin{subfigure}[b]{0.48\textwidth}
    \includegraphics[width=\textwidth]{figures/appendix/regression/drift/anyhow_window_15/set_sizes_per_batch}
    \caption{Set sizes per batch}
  \end{subfigure}
  \caption{Regression, drift, window $K{=}15$.}
  \label{fig:app-reg-d-window-15}
\end{figure}

\subsubsection{With drift --- Window \texorpdfstring{$K{=}25$}{K=25}.}
\begin{figure}[H]
  \centering
  \begin{subfigure}[b]{0.48\textwidth}
    \includegraphics[width=\textwidth]{figures/appendix/regression/drift/anyhow_window_25/coverage_distribution}
    \caption{Coverage distribution}
  \end{subfigure}
  \hfill
  \begin{subfigure}[b]{0.48\textwidth}
    \includegraphics[width=\textwidth]{figures/appendix/regression/drift/anyhow_window_25/pvalue_trajectory}
    \caption{P-value trajectory}
  \end{subfigure}
  \medskip
  \begin{subfigure}[b]{0.48\textwidth}
    \includegraphics[width=\textwidth]{figures/appendix/regression/drift/anyhow_window_25/regression_bands}
    \caption{Regression bands}
  \end{subfigure}
  \hfill
  \begin{subfigure}[b]{0.48\textwidth}
    \includegraphics[width=\textwidth]{figures/appendix/regression/drift/anyhow_window_25/set_size_histogram}
    \caption{Set size histogram}
  \end{subfigure}
  \medskip
  \centering
  \begin{subfigure}[b]{0.48\textwidth}
    \includegraphics[width=\textwidth]{figures/appendix/regression/drift/anyhow_window_25/set_sizes_per_batch}
    \caption{Set sizes per batch}
  \end{subfigure}
  \caption{Regression, drift, window $K{=}25$.}
  \label{fig:app-reg-d-window-25}
\end{figure}

\subsection{Classification}
\subsubsection{Without drift --- All-ones.}
\begin{figure}[H]
  \centering
  \begin{subfigure}[b]{0.48\textwidth}
    \includegraphics[width=\textwidth]{figures/appendix/classification/nodrift/anyhow_all_ones/coverage_distribution}
    \caption{Coverage distribution}
  \end{subfigure}
  \hfill
  \begin{subfigure}[b]{0.48\textwidth}
    \includegraphics[width=\textwidth]{figures/appendix/classification/nodrift/anyhow_all_ones/pvalue_trajectory}
    \caption{P-value trajectory}
  \end{subfigure}
  \medskip
  \begin{subfigure}[b]{0.48\textwidth}
    \includegraphics[width=\textwidth]{figures/appendix/classification/nodrift/anyhow_all_ones/conformal_sequence_matrix}
    \caption{Conformal sequence matrix}
  \end{subfigure}
  \hfill
  \begin{subfigure}[b]{0.48\textwidth}
    \includegraphics[width=\textwidth]{figures/appendix/classification/nodrift/anyhow_all_ones/set_size_histogram}
    \caption{Set size histogram}
  \end{subfigure}
  \medskip
  \centering
  \begin{subfigure}[b]{0.48\textwidth}
    \includegraphics[width=\textwidth]{figures/appendix/classification/nodrift/anyhow_all_ones/set_sizes_per_batch}
    \caption{Set sizes per batch}
  \end{subfigure}
  \caption{Classification, no drift, all-ones.}
  \label{fig:app-cla-nd-all-ones}
\end{figure}

\subsubsection{Without drift --- Window \texorpdfstring{$K{=}1$}{K=1}.}
\begin{figure}[H]
  \centering
  \begin{subfigure}[b]{0.48\textwidth}
    \includegraphics[width=\textwidth]{figures/appendix/classification/nodrift/anyhow_window_1/coverage_distribution}
    \caption{Coverage distribution}
  \end{subfigure}
  \hfill
  \begin{subfigure}[b]{0.48\textwidth}
    \includegraphics[width=\textwidth]{figures/appendix/classification/nodrift/anyhow_window_1/pvalue_trajectory}
    \caption{P-value trajectory}
  \end{subfigure}
  \medskip
  \begin{subfigure}[b]{0.48\textwidth}
    \includegraphics[width=\textwidth]{figures/appendix/classification/nodrift/anyhow_window_1/conformal_sequence_matrix}
    \caption{Conformal sequence matrix}
  \end{subfigure}
  \hfill
  \begin{subfigure}[b]{0.48\textwidth}
    \includegraphics[width=\textwidth]{figures/appendix/classification/nodrift/anyhow_window_1/set_size_histogram}
    \caption{Set size histogram}
  \end{subfigure}
  \medskip
  \centering
  \begin{subfigure}[b]{0.48\textwidth}
    \includegraphics[width=\textwidth]{figures/appendix/classification/nodrift/anyhow_window_1/set_sizes_per_batch}
    \caption{Set sizes per batch}
  \end{subfigure}
  \caption{Classification, no drift, window $K{=}1$.}
  \label{fig:app-cla-nd-window-1}
\end{figure}

\subsubsection{Without drift --- Window \texorpdfstring{$K{=}3$}{K=3}.}
\begin{figure}[H]
  \centering
  \begin{subfigure}[b]{0.48\textwidth}
    \includegraphics[width=\textwidth]{figures/appendix/classification/nodrift/anyhow_window_3/coverage_distribution}
    \caption{Coverage distribution}
  \end{subfigure}
  \hfill
  \begin{subfigure}[b]{0.48\textwidth}
    \includegraphics[width=\textwidth]{figures/appendix/classification/nodrift/anyhow_window_3/pvalue_trajectory}
    \caption{P-value trajectory}
  \end{subfigure}
  \medskip
  \begin{subfigure}[b]{0.48\textwidth}
    \includegraphics[width=\textwidth]{figures/appendix/classification/nodrift/anyhow_window_3/conformal_sequence_matrix}
    \caption{Conformal sequence matrix}
  \end{subfigure}
  \hfill
  \begin{subfigure}[b]{0.48\textwidth}
    \includegraphics[width=\textwidth]{figures/appendix/classification/nodrift/anyhow_window_3/set_size_histogram}
    \caption{Set size histogram}
  \end{subfigure}
  \medskip
  \centering
  \begin{subfigure}[b]{0.48\textwidth}
    \includegraphics[width=\textwidth]{figures/appendix/classification/nodrift/anyhow_window_3/set_sizes_per_batch}
    \caption{Set sizes per batch}
  \end{subfigure}
  \caption{Classification, no drift, window $K{=}3$.}
  \label{fig:app-cla-nd-window-3}
\end{figure}

\subsubsection{Without drift --- Window \texorpdfstring{$K{=}5$}{K=5}.}
\begin{figure}[H]
  \centering
  \begin{subfigure}[b]{0.48\textwidth}
    \includegraphics[width=\textwidth]{figures/appendix/classification/nodrift/anyhow_window_5/coverage_distribution}
    \caption{Coverage distribution}
  \end{subfigure}
  \hfill
  \begin{subfigure}[b]{0.48\textwidth}
    \includegraphics[width=\textwidth]{figures/appendix/classification/nodrift/anyhow_window_5/pvalue_trajectory}
    \caption{P-value trajectory}
  \end{subfigure}
  \medskip
  \begin{subfigure}[b]{0.48\textwidth}
    \includegraphics[width=\textwidth]{figures/appendix/classification/nodrift/anyhow_window_5/conformal_sequence_matrix}
    \caption{Conformal sequence matrix}
  \end{subfigure}
  \hfill
  \begin{subfigure}[b]{0.48\textwidth}
    \includegraphics[width=\textwidth]{figures/appendix/classification/nodrift/anyhow_window_5/set_size_histogram}
    \caption{Set size histogram}
  \end{subfigure}
  \medskip
  \centering
  \begin{subfigure}[b]{0.48\textwidth}
    \includegraphics[width=\textwidth]{figures/appendix/classification/nodrift/anyhow_window_5/set_sizes_per_batch}
    \caption{Set sizes per batch}
  \end{subfigure}
  \caption{Classification, no drift, window $K{=}5$.}
  \label{fig:app-cla-nd-window-5}
\end{figure}

\subsubsection{Without drift --- Window \texorpdfstring{$K{=}10$}{K=10}.}
\begin{figure}[H]
  \centering
  \begin{subfigure}[b]{0.48\textwidth}
    \includegraphics[width=\textwidth]{figures/appendix/classification/nodrift/anyhow_window_10/coverage_distribution}
    \caption{Coverage distribution}
  \end{subfigure}
  \hfill
  \begin{subfigure}[b]{0.48\textwidth}
    \includegraphics[width=\textwidth]{figures/appendix/classification/nodrift/anyhow_window_10/pvalue_trajectory}
    \caption{P-value trajectory}
  \end{subfigure}
  \medskip
  \begin{subfigure}[b]{0.48\textwidth}
    \includegraphics[width=\textwidth]{figures/appendix/classification/nodrift/anyhow_window_10/conformal_sequence_matrix}
    \caption{Conformal sequence matrix}
  \end{subfigure}
  \hfill
  \begin{subfigure}[b]{0.48\textwidth}
    \includegraphics[width=\textwidth]{figures/appendix/classification/nodrift/anyhow_window_10/set_size_histogram}
    \caption{Set size histogram}
  \end{subfigure}
  \medskip
  \centering
  \begin{subfigure}[b]{0.48\textwidth}
    \includegraphics[width=\textwidth]{figures/appendix/classification/nodrift/anyhow_window_10/set_sizes_per_batch}
    \caption{Set sizes per batch}
  \end{subfigure}
  \caption{Classification, no drift, window $K{=}10$.}
  \label{fig:app-cla-nd-window-10}
\end{figure}

\subsubsection{Without drift --- Window \texorpdfstring{$K{=}15$}{K=15}.}
\begin{figure}[H]
  \centering
  \begin{subfigure}[b]{0.48\textwidth}
    \includegraphics[width=\textwidth]{figures/appendix/classification/nodrift/anyhow_window_15/coverage_distribution}
    \caption{Coverage distribution}
  \end{subfigure}
  \hfill
  \begin{subfigure}[b]{0.48\textwidth}
    \includegraphics[width=\textwidth]{figures/appendix/classification/nodrift/anyhow_window_15/pvalue_trajectory}
    \caption{P-value trajectory}
  \end{subfigure}
  \medskip
  \begin{subfigure}[b]{0.48\textwidth}
    \includegraphics[width=\textwidth]{figures/appendix/classification/nodrift/anyhow_window_15/conformal_sequence_matrix}
    \caption{Conformal sequence matrix}
  \end{subfigure}
  \hfill
  \begin{subfigure}[b]{0.48\textwidth}
    \includegraphics[width=\textwidth]{figures/appendix/classification/nodrift/anyhow_window_15/set_size_histogram}
    \caption{Set size histogram}
  \end{subfigure}
  \medskip
  \centering
  \begin{subfigure}[b]{0.48\textwidth}
    \includegraphics[width=\textwidth]{figures/appendix/classification/nodrift/anyhow_window_15/set_sizes_per_batch}
    \caption{Set sizes per batch}
  \end{subfigure}
  \caption{Classification, no drift, window $K{=}15$.}
  \label{fig:app-cla-nd-window-15}
\end{figure}

\subsubsection{Without drift --- Window \texorpdfstring{$K{=}25$}{K=25}.}
\begin{figure}[H]
  \centering
  \begin{subfigure}[b]{0.48\textwidth}
    \includegraphics[width=\textwidth]{figures/appendix/classification/nodrift/anyhow_window_25/coverage_distribution}
    \caption{Coverage distribution}
  \end{subfigure}
  \hfill
  \begin{subfigure}[b]{0.48\textwidth}
    \includegraphics[width=\textwidth]{figures/appendix/classification/nodrift/anyhow_window_25/pvalue_trajectory}
    \caption{P-value trajectory}
  \end{subfigure}
  \medskip
  \begin{subfigure}[b]{0.48\textwidth}
    \includegraphics[width=\textwidth]{figures/appendix/classification/nodrift/anyhow_window_25/conformal_sequence_matrix}
    \caption{Conformal sequence matrix}
  \end{subfigure}
  \hfill
  \begin{subfigure}[b]{0.48\textwidth}
    \includegraphics[width=\textwidth]{figures/appendix/classification/nodrift/anyhow_window_25/set_size_histogram}
    \caption{Set size histogram}
  \end{subfigure}
  \medskip
  \centering
  \begin{subfigure}[b]{0.48\textwidth}
    \includegraphics[width=\textwidth]{figures/appendix/classification/nodrift/anyhow_window_25/set_sizes_per_batch}
    \caption{Set sizes per batch}
  \end{subfigure}
  \caption{Classification, no drift, window $K{=}25$.}
  \label{fig:app-cla-nd-window-25}
\end{figure}

\subsubsection{With drift --- All-ones.}
\begin{figure}[H]
  \centering
  \begin{subfigure}[b]{0.48\textwidth}
    \includegraphics[width=\textwidth]{figures/appendix/classification/drift/anyhow_all_ones/coverage_distribution}
    \caption{Coverage distribution}
  \end{subfigure}
  \hfill
  \begin{subfigure}[b]{0.48\textwidth}
    \includegraphics[width=\textwidth]{figures/appendix/classification/drift/anyhow_all_ones/pvalue_trajectory}
    \caption{P-value trajectory}
  \end{subfigure}
  \medskip
  \begin{subfigure}[b]{0.48\textwidth}
    \includegraphics[width=\textwidth]{figures/appendix/classification/drift/anyhow_all_ones/conformal_sequence_matrix}
    \caption{Conformal sequence matrix}
  \end{subfigure}
  \hfill
  \begin{subfigure}[b]{0.48\textwidth}
    \includegraphics[width=\textwidth]{figures/appendix/classification/drift/anyhow_all_ones/set_size_histogram}
    \caption{Set size histogram}
  \end{subfigure}
  \medskip
  \centering
  \begin{subfigure}[b]{0.48\textwidth}
    \includegraphics[width=\textwidth]{figures/appendix/classification/drift/anyhow_all_ones/set_sizes_per_batch}
    \caption{Set sizes per batch}
  \end{subfigure}
  \caption{Classification, drift, all-ones.}
  \label{fig:app-cla-d-all-ones}
\end{figure}

\subsubsection{With drift --- Window \texorpdfstring{$K{=}1$}{K=1}.}
\begin{figure}[H]
  \centering
  \begin{subfigure}[b]{0.48\textwidth}
    \includegraphics[width=\textwidth]{figures/appendix/classification/drift/anyhow_window_1/coverage_distribution}
    \caption{Coverage distribution}
  \end{subfigure}
  \hfill
  \begin{subfigure}[b]{0.48\textwidth}
    \includegraphics[width=\textwidth]{figures/appendix/classification/drift/anyhow_window_1/pvalue_trajectory}
    \caption{P-value trajectory}
  \end{subfigure}
  \medskip
  \begin{subfigure}[b]{0.48\textwidth}
    \includegraphics[width=\textwidth]{figures/appendix/classification/drift/anyhow_window_1/conformal_sequence_matrix}
    \caption{Conformal sequence matrix}
  \end{subfigure}
  \hfill
  \begin{subfigure}[b]{0.48\textwidth}
    \includegraphics[width=\textwidth]{figures/appendix/classification/drift/anyhow_window_1/set_size_histogram}
    \caption{Set size histogram}
  \end{subfigure}
  \medskip
  \centering
  \begin{subfigure}[b]{0.48\textwidth}
    \includegraphics[width=\textwidth]{figures/appendix/classification/drift/anyhow_window_1/set_sizes_per_batch}
    \caption{Set sizes per batch}
  \end{subfigure}
  \caption{Classification, drift, window $K{=}1$.}
  \label{fig:app-cla-d-window-1}
\end{figure}

\subsubsection{With drift --- Window \texorpdfstring{$K{=}3$}{K=3}.}
\begin{figure}[H]
  \centering
  \begin{subfigure}[b]{0.48\textwidth}
    \includegraphics[width=\textwidth]{figures/appendix/classification/drift/anyhow_window_3/coverage_distribution}
    \caption{Coverage distribution}
  \end{subfigure}
  \hfill
  \begin{subfigure}[b]{0.48\textwidth}
    \includegraphics[width=\textwidth]{figures/appendix/classification/drift/anyhow_window_3/pvalue_trajectory}
    \caption{P-value trajectory}
  \end{subfigure}
  \medskip
  \begin{subfigure}[b]{0.48\textwidth}
    \includegraphics[width=\textwidth]{figures/appendix/classification/drift/anyhow_window_3/conformal_sequence_matrix}
    \caption{Conformal sequence matrix}
  \end{subfigure}
  \hfill
  \begin{subfigure}[b]{0.48\textwidth}
    \includegraphics[width=\textwidth]{figures/appendix/classification/drift/anyhow_window_3/set_size_histogram}
    \caption{Set size histogram}
  \end{subfigure}
  \medskip
  \centering
  \begin{subfigure}[b]{0.48\textwidth}
    \includegraphics[width=\textwidth]{figures/appendix/classification/drift/anyhow_window_3/set_sizes_per_batch}
    \caption{Set sizes per batch}
  \end{subfigure}
  \caption{Classification, drift, window $K{=}3$.}
  \label{fig:app-cla-d-window-3}
\end{figure}

\subsubsection{With drift --- Window \texorpdfstring{$K{=}5$}{K=5}.}
\begin{figure}[H]
  \centering
  \begin{subfigure}[b]{0.48\textwidth}
    \includegraphics[width=\textwidth]{figures/appendix/classification/drift/anyhow_window_5/coverage_distribution}
    \caption{Coverage distribution}
  \end{subfigure}
  \hfill
  \begin{subfigure}[b]{0.48\textwidth}
    \includegraphics[width=\textwidth]{figures/appendix/classification/drift/anyhow_window_5/pvalue_trajectory}
    \caption{P-value trajectory}
  \end{subfigure}
  \medskip
  \begin{subfigure}[b]{0.48\textwidth}
    \includegraphics[width=\textwidth]{figures/appendix/classification/drift/anyhow_window_5/conformal_sequence_matrix}
    \caption{Conformal sequence matrix}
  \end{subfigure}
  \hfill
  \begin{subfigure}[b]{0.48\textwidth}
    \includegraphics[width=\textwidth]{figures/appendix/classification/drift/anyhow_window_5/set_size_histogram}
    \caption{Set size histogram}
  \end{subfigure}
  \medskip
  \centering
  \begin{subfigure}[b]{0.48\textwidth}
    \includegraphics[width=\textwidth]{figures/appendix/classification/drift/anyhow_window_5/set_sizes_per_batch}
    \caption{Set sizes per batch}
  \end{subfigure}
  \caption{Classification, drift, window $K{=}5$.}
  \label{fig:app-cla-d-window-5}
\end{figure}

\subsubsection{With drift --- Window \texorpdfstring{$K{=}10$}{K=10}.}
\begin{figure}[H]
  \centering
  \begin{subfigure}[b]{0.48\textwidth}
    \includegraphics[width=\textwidth]{figures/appendix/classification/drift/anyhow_window_10/coverage_distribution}
    \caption{Coverage distribution}
  \end{subfigure}
  \hfill
  \begin{subfigure}[b]{0.48\textwidth}
    \includegraphics[width=\textwidth]{figures/appendix/classification/drift/anyhow_window_10/pvalue_trajectory}
    \caption{P-value trajectory}
  \end{subfigure}
  \medskip
  \begin{subfigure}[b]{0.48\textwidth}
    \includegraphics[width=\textwidth]{figures/appendix/classification/drift/anyhow_window_10/conformal_sequence_matrix}
    \caption{Conformal sequence matrix}
  \end{subfigure}
  \hfill
  \begin{subfigure}[b]{0.48\textwidth}
    \includegraphics[width=\textwidth]{figures/appendix/classification/drift/anyhow_window_10/set_size_histogram}
    \caption{Set size histogram}
  \end{subfigure}
  \medskip
  \centering
  \begin{subfigure}[b]{0.48\textwidth}
    \includegraphics[width=\textwidth]{figures/appendix/classification/drift/anyhow_window_10/set_sizes_per_batch}
    \caption{Set sizes per batch}
  \end{subfigure}
  \caption{Classification, drift, window $K{=}10$.}
  \label{fig:app-cla-d-window-10}
\end{figure}

\subsubsection{With drift --- Window \texorpdfstring{$K{=}15$}{K=15}.}
\begin{figure}[H]
  \centering
  \begin{subfigure}[b]{0.48\textwidth}
    \includegraphics[width=\textwidth]{figures/appendix/classification/drift/anyhow_window_15/coverage_distribution}
    \caption{Coverage distribution}
  \end{subfigure}
  \hfill
  \begin{subfigure}[b]{0.48\textwidth}
    \includegraphics[width=\textwidth]{figures/appendix/classification/drift/anyhow_window_15/pvalue_trajectory}
    \caption{P-value trajectory}
  \end{subfigure}
  \medskip
  \begin{subfigure}[b]{0.48\textwidth}
    \includegraphics[width=\textwidth]{figures/appendix/classification/drift/anyhow_window_15/conformal_sequence_matrix}
    \caption{Conformal sequence matrix}
  \end{subfigure}
  \hfill
  \begin{subfigure}[b]{0.48\textwidth}
    \includegraphics[width=\textwidth]{figures/appendix/classification/drift/anyhow_window_15/set_size_histogram}
    \caption{Set size histogram}
  \end{subfigure}
  \medskip
  \centering
  \begin{subfigure}[b]{0.48\textwidth}
    \includegraphics[width=\textwidth]{figures/appendix/classification/drift/anyhow_window_15/set_sizes_per_batch}
    \caption{Set sizes per batch}
  \end{subfigure}
  \caption{Classification, drift, window $K{=}15$.}
  \label{fig:app-cla-d-window-15}
\end{figure}

\subsubsection{With drift --- Window \texorpdfstring{$K{=}25$}{K=25}.}
\begin{figure}[H]
  \centering
  \begin{subfigure}[b]{0.48\textwidth}
    \includegraphics[width=\textwidth]{figures/appendix/classification/drift/anyhow_window_25/coverage_distribution}
    \caption{Coverage distribution}
  \end{subfigure}
  \hfill
  \begin{subfigure}[b]{0.48\textwidth}
    \includegraphics[width=\textwidth]{figures/appendix/classification/drift/anyhow_window_25/pvalue_trajectory}
    \caption{P-value trajectory}
  \end{subfigure}
  \medskip
  \begin{subfigure}[b]{0.48\textwidth}
    \includegraphics[width=\textwidth]{figures/appendix/classification/drift/anyhow_window_25/conformal_sequence_matrix}
    \caption{Conformal sequence matrix}
  \end{subfigure}
  \hfill
  \begin{subfigure}[b]{0.48\textwidth}
    \includegraphics[width=\textwidth]{figures/appendix/classification/drift/anyhow_window_25/set_size_histogram}
    \caption{Set size histogram}
  \end{subfigure}
  \medskip
  \centering
  \begin{subfigure}[b]{0.48\textwidth}
    \includegraphics[width=\textwidth]{figures/appendix/classification/drift/anyhow_window_25/set_sizes_per_batch}
    \caption{Set sizes per batch}
  \end{subfigure}
  \caption{Classification, drift, window $K{=}25$.}
  \label{fig:app-cla-d-window-25}
\end{figure}


\end{document}